\documentclass[11pt]{article}

\usepackage[a4paper,margin=1in]{geometry}
\usepackage{amsmath,amssymb,amsthm,mathtools}
\usepackage{booktabs,tabularx,longtable,array}
\usepackage{graphicx}
\usepackage{xcolor}
\usepackage{hyperref}
\hypersetup{hidelinks,hypertexnames=false}
\usepackage{microtype}
\usepackage{enumitem}

\newtheorem{theorem}{Theorem}[section]
\newtheorem{proposition}[theorem]{Proposition}
\newtheorem{lemma}[theorem]{Lemma}
\newtheorem{corollary}[theorem]{Corollary}
\newtheorem{definition}[theorem]{Definition}
\newtheorem{remark}[theorem]{Remark}
\newtheorem{example}[theorem]{Example}

\newcommand{\ord}{\operatorname{ord}}
\newcommand{\rad}{\operatorname{rad}}
\newcommand{\Spec}{\operatorname{Spec}}
\newcommand{\Eset}{\mathcal E}
\newcommand{\Pset}{\mathcal P}

\title{\textbf{Reciprocal Periods and Primality}\\
\large A period--primality toolbox and the rarity of small-index pseudoprimes}
\author{Extended arXiv version}
\date{}

\begin{document}
\maketitle

\begin{abstract}
Suppose that the exact reciprocal period
\[
L=L_b(n)=\ord_n(b)
\]
is known for an odd integer $n$ and a base $b$ coprime to $n$.  How much
information about the primality of $n$ is already contained in $L$?  This
expository question is the starting point of the present extended version.  We
assemble a period--primality toolbox connecting multiplicative order with
Lucas' criterion, Fermat and strong pseudoprimality, Midy's property,
overpseudoprimes, prime-only tests for total Midy, $q$-Midy congruence
forcing, cyclotomic factors, and $N\pm1$ certificates of Pocklington and
Brillhart--Lehmer--Selfridge type.  The toolbox assumes that the exact order is
already available; it is not proposed as a general-purpose primality algorithm.

The particularly important case $L\mid n-1$ leads to the reciprocal-period
index
\[
d=\frac{n-1}{L}.
\]
For composite $n$ this is Somer's Fermat $d$-pseudoprime condition; for
prime moduli the same $d$ is the standard residual index.  The special
$\kappa=1$ composite stratum is the prescribed Carmichael-index problem.
Somer proved fixed-$d$ finiteness and classified the nontrivial-base cases for
$1\le d\le8$.  Our spectral investigation gives two direct primality
consequences.  First, for $d\le64$ an odd integer with
$L=(n-1)/d$ is prime except for the explicitly certified composite table,
which contains $234$ distinct integers.  Second, if
\[
S(d)=1+\lfloor\log_2(d-1)\rfloor,\qquad
M(d)=2^{S(d)-1},
\]
then every composite spectral point satisfies
\[
n<(dS(d))^{2M(d)}.
\]
Consequently, for every fixed $\varepsilon>0$, all sufficiently large $n$
with
\[
d\le\left(\frac12-\varepsilon\right)
       \frac{\log n}{\log\log n}
\]
are prime.

The base-independent genuine spectrum is also shown to be a partition of all
odd composite integers: each odd composite appears exactly once at its genuine
index
\[
d_\lambda(n)=\frac{n-1}{\gcd(n-1,\lambda(n))},
\]
while all later appearances are inherited.  We certify the complete inclusive
spectrum through $d=64$, study the fixed-base prime spectrum via the
Wagstaff--Moree near-primitive-root densities, and prove under GRH that the
same range contains at least $98.2709160681865\ldots\%$ of the prime mass for
every positive non-perfect-power integer base.  We also introduce the counting functions $e(d)$ and $g(d)$.  The
$\kappa$--$H$ search decomposition has
$\frac12d(\log d)^2+O(d\log d)$ overlap-resolved branches, while collapsing to the Diophantine master equation gives a relaxed parameter count of only
$d\log d+O(d)$ candidate pairs.  A divisor-pair identity for the
last two prime factors then yields the uniform counting estimate
\[
\log e(d)\le
\frac d2\bigl(\log d+\log\log d+O(1)\bigr).
\]
In particular the one-, two-, and three-distinct-prime strata satisfy
$e_1(d)\le d^{o(1)}$ and $e_2(d),e_3(d)\le d^{1+o(1)}$.
Under a natural uniform sampling model on a spectral line, for fixed $d$ we
prove
\[
\Pr_X(n\text{ prime}\mid d)
=1-e(d)\phi(d)\frac{\log X}{X}
+o\!\left(\frac{\log X}{X}\right).
\]
More strongly, for every fixed $0<\varepsilon<2$ this probability tends to
$1$ uniformly throughout
\[
1\le d\le(2-\varepsilon)\frac{\log X}{\log\log X}.
\]
Thus the deterministic certainty range with leading constant $1/2$ is
complemented by a substantially wider high-probability range, in the stated
sampling model, with leading constant $2$.  Four proposed OEIS sequences are
printed in full through $d=64$.
\end{abstract}

\noindent\textbf{Keywords.} Fermat $d$-pseudoprimes; residual index;
near-primitive roots; Carmichael index; reciprocal periods; Carmichael function;
Miller--Rabin; computer-assisted proof.

\medskip
\tableofcontents
\bigskip


\section{Introduction}
\label{sec:introduction}

The elementary decimal identity
\[
\frac17=0.\overline{142857}
\]
contains a surprisingly strong arithmetic statement.  The six-digit period
means
\[
\ord_7(10)=6=7-1,
\]
and the classical Lucas criterion immediately proves that $7$ is prime.  More
generally, if the exact period
\[
L=L_b(n)=\ord_n(b)
\]
is known, then the natural question is:

\begin{center}
\fbox{\parbox{0.88\textwidth}{
\centering
\textbf{How much information about primality is encoded in the exact period
$L$?}
}}
\end{center}

This question, rather than a pre-existing pseudoprime classification, was the
motivation that led to the present project.  The first goal was to collect in
one place the most useful consequences of knowing $L$: direct Lucas-type
criteria, Fermat and strong-pseudoprime implications, Midy decompositions,
overpseudoprimes, cyclotomic information, and positive $N-1$ certificates of
Pocklington and Brillhart--Lehmer--Selfridge type.  The resulting toolbox is
given in Sections~\ref{sec:toolbox-prelim} and~\ref{sec:toolbox} with explicit
examples.

A second question emerged naturally.  When
\[
L\mid n-1,
\]
write
\begin{equation}
 n-1=dL,
 \qquad d=\frac{n-1}{L}.
 \label{eq:indexdef-intro}
\end{equation}
For a prime $p$, $d$ is the index of the cyclic subgroup generated by $b$ in
$(\mathbb Z/p\mathbb Z)^\times$.  A small value of $d$ therefore says that the
known period is close to the maximal prime value $p-1$.  It is natural to ask
whether sufficiently small $d$ forces primality.

For example,
\[
91=7\cdot13,\qquad \ord_{91}(10)=6,
\]
so $91-1=15\cdot6$.  Thus a composite can have $L\mid n-1$, but here the
index is already $d=15$.  The present paper studies systematically how small
such an index can be for a composite.

\subsection*{The rediscovery of Somer's $d$-pseudoprimes}

During the development of this project, after the reciprocal index
\eqref{eq:indexdef-intro} and the first small-index arguments had already been
obtained independently, an AI-assisted literature search located Lawrence
Somer's 1985 dissertation and his 1989 paper \emph{On Fermat
$d$-pseudoprimes} \cite{Somer1985,Somer1989}.  Somer had introduced precisely
the composite condition
\[
\ord_N(a)=\frac{N-1}{d},
\]
proved that for each fixed $d$ only finitely many such composites exist, and
classified the nontrivial-base cases for $1\le d\le8$.  Our parameter $d$
is therefore Somer's parameter.

Somer's work was not wholly forgotten, but the independent uptake that we
have been able to locate is surprisingly concentrated.  Two examples are
particularly clear.  Ribenboim did more than cite the paper: he devoted
Section~VIII.E, entitled ``Somer Pseudoprimes,'' of \emph{The New Book of Prime
Number Records} to the subject \cite[p.~117]{Ribenboim1996}.  A decade later,
Pinch's extended abstract \emph{Carmichael numbers with small index} opens by
stating that Somer had proved a result implying that the Carmichael index
$(N-1)/\lambda(N)$ tends to infinity, then gives a new proof and an enumeration
algorithm for bounded Carmichael index \cite{PinchSmallIndex}.  Pinch's 2006
tabulation through $10^{18}$ also cites Somer's Fermat $d$-pseudoprime paper
\cite{Pinch2006e18}.

A separate line of continuation came from Somer and his collaborators.  Somer
developed the analogous Lucas $d$-pseudoprime theory \cite{SomerLucas1998};
Carlip, Jacobson and Somer connected the same methods with Lehmer's problem and
multiperfect numbers \cite{CarlipJacobsonSomer1998}; and later work developed
Lucas $d$-pseudoprimes and Carmichael--Lucas numbers further
\cite{CarlipSomer2007}.

The Carlip--Jacobson--Somer paper is particularly relevant to the Fermat side
of the present work.  It develops Somer's methods in a unified setting and,
for a Fermat $d$-pseudoprime with $t$ distinct prime factors, proves the
factor-count restriction
\[
 t<\log_2 d+1,
\]
together with an explicit bound on repeated prime powers.  These are direct
antecedents of part of the structural theory below.  Our refinement retains the
additional overlap parameter $\kappa=\lambda(n)/L$: the exact divisibilities
$K\mid\kappa$ and $2^{s-1}\kappa\mid j$ sharpen the factor-count and
repeated-power restrictions, while the uniform closed size bound, the
last-two-prime divisor-pair identity, and the counting estimate for $e(d)$ are
additional quantitative consequences developed here.

Our targeted literature search has not revealed a comparably sustained
independent literature using the \emph{variable-base Fermat} $d$-spectrum as a
primary organizing framework.  We therefore describe Somer's Fermat theory as
under-recognized rather than forgotten.  This is not a bibliometric census or
a numerical citation claim: database coverage of an older conference chapter
is imperfect, and absence from a search is not evidence of absence.  What is
relevant for the present paper is that the same reciprocal index arose
independently here from a different motivation, namely the information content
of a known period.

This convergence of two routes is mathematically useful.  Somer's viewpoint
starts from generalized Fermat pseudoprimes.  Our route starts from reciprocal
periods and primality.  Both single out
\[
d=\frac{n-1}{L}
\]
as the natural coordinate measuring the deficiency of the period from the
Lucas extremal value $L=n-1$.

\subsection*{Two primality theorems produced by the spectral investigation}

The spectral part of the paper yields two answers to the motivating primality
question.

\begin{enumerate}[label=\textbf{P\arabic*.},leftmargin=2.4em]
\item \textbf{Exact low-index criterion.}
If $n$ is odd, $\gcd(b,n)=1$, $L=\ord_n(b)$, and
\[
L=\frac{n-1}{d},\qquad d\le64,
\]
then $n$ is prime unless it occurs in the certified exceptional table of
Section~\ref{sec:certified-search}.  Across all $64$ rows there are $234$ distinct composite exceptions in $351$ spectral incidences.

\item \textbf{Uniform size criterion.}
For $d\ge2$ put
\[
S(d)=1+\lfloor\log_2(d-1)\rfloor,
\qquad
M(d)=2^{S(d)-1}.
\]
Theorem~\ref{thm:uniform-U} proves that every composite spectral point obeys
\[
n<(dS(d))^{2M(d)}.
\]
Hence
\[
n\ge(dS(d))^{2M(d)},\qquad L=(n-1)/d
\quad\Longrightarrow\quad n\text{ is prime}.
\]
Asymptotically, for each fixed $\varepsilon>0$,
\[
d\le\left(\frac12-\varepsilon\right)
\frac{\log n}{\log\log n}
\]
forces primality for all sufficiently large $n$.
\end{enumerate}

The first statement is finite and exact; the second is uniform in $d$.  The
coefficient $1/2$ in the asymptotic statement is the worst leading coefficient
of our uniform bound, not a claim about the true optimal boundary.

They are complemented by a statistical consequence once a sampling model is
specified.  If one samples uniformly from all integers up to $X$ lying on one
fixed existential line $d$, Theorem~\ref{thm:conditional-prime-prob} gives
\[
\Pr_X(n\text{ prime}\mid d)
=1-e(d)\phi(d)\frac{\log X}{X}
+o\!\left(\frac{\log X}{X}\right).
\]
Thus fixed small index is not merely compatible with primality: in this precise
line-sampling sense it becomes asymptotically overwhelming.  This probability
statement is not one of the two deterministic primality theorems above and is
not an intrinsic Bayesian probability attached to a single integer.

\subsection*{Prime--pseudoprime spectral separation}

When the base is allowed to vary, the prime line at a fixed $d$ is simply
\[
\mathcal P(d)=\{p\text{ prime}:p\equiv1\pmod d\},
\]
and is therefore infinite.  Somer's theorem says that the corresponding
composite line is finite.  The present work strengthens the comparison in two
directions.  On the composite side we derive explicit structural bounds and
certify the complete spectrum through $d=64$.  On the prime side we compare
the same index range with the fixed-base near-primitive-root densities of
Wagstaff and Moree.  Under GRH, at least
$98.2709160681865\ldots\%$ of the prime mass lies in $d\le64$ for every
positive integer base that is not a perfect power, while the inclusive
composite union through $64$ contains $234$ integers.

A second structural point is equally important.  Allowing the trivial orders
$1$ and $2$ gives an inclusive group-theoretic spectrum in which \emph{every
odd composite integer appears at a unique genuine index}.  Thus the genuine
spectrum is not merely a collection of exceptional examples: it is a
partition of the odd composites.  Later spectral occurrences are inherited
from the genuine one.

\subsection*{Guide to the extended version}

This arXiv version is deliberately more expository than the compact working
paper.  Section~\ref{sec:period-groups} introduces multiplicative order,
group exponent, the Carmichael function, Chinese-remainder decomposition and
Sylow components.  Sections~\ref{sec:toolbox-prelim} and~\ref{sec:toolbox}
then develop the period--primality toolbox with examples.  The spectral
structure, genuine spectrum and uniform bounds follow next.  Somer's
small-index classification is reproved uniformly, and the exact computation through $d=64$ is tied to a frozen supplement combining the original Julia certification through $33$ with the exact extension through $64$.
The final sections treat counting functions, conditional primality
probabilities, Artin's constant and fixed-base prime spectra, and four natural
OEIS sequences.

\subsection*{Notation at a glance}

The following symbols recur throughout the paper.
\begin{center}
\small
\begin{tabularx}{0.94\textwidth}{@{}lX@{}}
\toprule
symbol & meaning\\
\midrule
$n$ & an odd integer under primality investigation\\
$b$ & a base with $\gcd(b,n)=1$\\
$L=L_b(n)=\ord_n(b)$ & the exact multiplicative order of $b$ modulo $n$, equivalently the reciprocal period\\
$\phi(n)$ & Euler's totient, the order of $(\mathbb Z/n\mathbb Z)^\times$\\
$\lambda(n)$ & Carmichael's function, the exponent of $(\mathbb Z/n\mathbb Z)^\times$\\
$d=(n-1)/L$ & the spectral or reciprocal-period index, when $L\mid n-1$\\
$L_{\max}(n)$ & $\gcd(n-1,\lambda(n))$, the largest element order compatible with $n-1$\\
$d_\lambda(n)$ & $(n-1)/L_{\max}(n)$, the unique genuine spectral index of a composite\\
$\mathcal E(d)$ & odd composites occurring on the inclusive existential line $d$\\
$\mathcal G(d)$ & odd composites whose genuine index equals $d$\\
$e(d),g(d)$ & $|\mathcal E(d)|$ and $|\mathcal G(d)|$\\
\bottomrule
\end{tabularx}
\end{center}

\section{Reciprocal periods and finite unit groups}
\label{sec:period-groups}

\subsection{Period length as multiplicative order}

Let $n>1$, $b\ge2$ and $\gcd(b,n)=1$.  The residue class of $b$ is then a
unit modulo $n$.  Its \emph{multiplicative order}
\[
\ord_n(b)
\]
is the least positive integer $L$ for which $b^L\equiv1\pmod n$.
Equivalently, $L$ is the period length of the purely recurring base-$b$
expansion of $1/n$.  We write
\[
L_b(n)=\ord_n(b).
\]

The units modulo $n$ form the finite abelian group
\[
U(n)=(\mathbb Z/n\mathbb Z)^\times.
\]
Its order is Euler's totient
\[
|U(n)|=\phi(n).
\]
Thus Lagrange's theorem already gives
\[
L_b(n)\mid\phi(n).
\]
The sharper universal divisor is the Carmichael function $\lambda(n)$.

\subsection{Exponent of a finite group and the Carmichael function}

\begin{definition}[Exponent]
If $G$ is a finite group, its \emph{exponent} is
\[
\exp(G)=\operatorname{lcm}\{\ord(g):g\in G\}.
\]
Equivalently, it is the least positive integer $e$ such that $g^e=1$ for every
$g\in G$.
\end{definition}

The Carmichael function is precisely the exponent of the unit group:
\[
\boxed{\lambda(n)=\exp U(n).}
\]
Consequently
\begin{equation}
L_b(n)\mid\lambda(n)\mid\phi(n).
\label{eq:ord-chain}
\end{equation}
The first divisibility says that every individual reciprocal period divides a
single base-independent invariant of $n$.

For later calculations it is useful to recall the explicit prime-power
formula.  If $p$ is odd, then
\[
\lambda(p^a)=\phi(p^a)=p^{a-1}(p-1).
\]
For powers of $2$,
\[
\lambda(2)=1,\qquad \lambda(4)=2,\qquad
\lambda(2^a)=2^{a-2}\quad(a\ge3).
\]
If $n=\prod_i p_i^{a_i}$, the Chinese remainder theorem gives
\[
U(n)\cong\prod_i U(p_i^{a_i}),
\]
and therefore
\[
\lambda(n)=\operatorname{lcm}_i\lambda(p_i^{a_i}).
\]

\subsection{Sylow components and attainable element orders}

A finite abelian group separates prime by prime.  Its \emph{Sylow
$q$-component} $G_q$ consists of the elements whose orders are powers of the
prime $q$, and
\[
G\cong\prod_{q\mid |G|}G_q.
\]
This simple decomposition is the reason that, in a finite abelian group,
there are no gaps among the divisors of the exponent.

\begin{lemma}[Orders dividing the exponent]
\label{lem:divisor-exponent}
Let $G$ be a finite abelian group of exponent $e$.  Every positive divisor of
$e$ occurs as the order of an element of $G$.
Consequently, for every $n>1$ and every $L\ge1$,
\[
\exists b:\ord_n(b)=L
\quad\Longleftrightarrow\quad
L\mid\lambda(n).
\]
\end{lemma}

\begin{proof}
Write
\[
e=\prod_q q^{a_q}.
\]
Because the exponent of the $q$-Sylow component is $q^{a_q}$, that component
contains an element of order $q^{a_q}$; taking powers produces elements of
every order $q^j$ with $0\le j\le a_q$.  If
$m=\prod_q q^{j_q}\mid e$, choose in each $G_q$ an element of order
$q^{j_q}$.  Their product, under the direct-product decomposition, has order
$m$.  Applying this to the abelian group $U(n)$ proves the second statement.
\end{proof}

\begin{remark}
This lemma is stronger than Lagrange's theorem.  In an arbitrary finite group,
a divisor of $|G|$ need not occur as an element order.  What matters here is
that the unit group is abelian and that we use divisors of its exponent, not
arbitrary divisors of its cardinality.
\end{remark}

\subsection{Spectral lines}

Whenever $L_b(n)\mid n-1$ we define
\[
d=\frac{n-1}{L_b(n)}.
\]
For primes, $U(p)$ is cyclic of order $p-1$, and $d$ is literally the index
of the subgroup generated by $b$:
\[
d=[U(p):\langle b\rangle].
\]
For this reason we call $d$ the \emph{reciprocal-period index} or
\emph{spectral index}.

\begin{definition}[Existential prime and composite spectra]
For $d\ge1$, define
\[
\Pset(d)
=
\left\{
p\text{ prime}:
\exists b\ge2,\ \gcd(b,p)=1,\
\ord_p(b)=\frac{p-1}{d}
\right\},
\]
and
\[
\Eset(d)
=
\left\{
n\text{ odd composite}:
\exists b\ge2,\ \gcd(b,n)=1,\
\ord_n(b)=\frac{n-1}{d}
\right\}.
\]
\end{definition}

For comparison with Somer's convention we also write
\[
\Eset^\ast(d)
=
\left\{
n\in\Eset(d):
\exists b\not\equiv\pm1\pmod n,\
\ord_n(b)=\frac{n-1}{d}
\right\}.
\]
Thus $\Eset^\ast(d)$ is the nontrivial-base Fermat $d$-pseudoprime line,
whereas $\Eset(d)$ is the inclusive group-theoretic spectral line.  The base
is allowed to depend on $n$.

\begin{theorem}[Prime spectral line]
\label{thm:prime-spectrum}
For every $d\ge1$,
\[
\Pset(d)=\{p\text{ prime}:p\equiv1\pmod d\}.
\]
Hence $\Pset(d)$ is infinite.  Its relative density among primes is
$1/\phi(d)$.
\end{theorem}

\begin{proof}
The cyclic group $U(p)$ contains an element of order $(p-1)/d$ if and only if
$d\mid p-1$.  Dirichlet's theorem gives infinitude and the prime number theorem
for arithmetic progressions gives the relative density.
\end{proof}

\begin{proposition}[Composite spectral line]
\label{prop:composite-spectrum}
For an odd composite $n$,
\[
n\in\Eset(d)
\quad\Longleftrightarrow\quad
d\mid n-1
\quad\text{and}\quad
\frac{n-1}{d}\mid\lambda(n).
\]
\end{proposition}

\begin{proof}
Apply Lemma~\ref{lem:divisor-exponent} with $L=(n-1)/d$.
\end{proof}

\section{Genuine and inherited spectral indices}

For a fixed composite $n$, one integer may occur on more than one
spectral line.  The natural distinguished line is determined by the
largest order compatible with $n-1$.

\begin{definition}
Define
\[
L_{\max}(n)=\gcd(n-1,\lambda(n)),
\qquad
d_\lambda(n)=\frac{n-1}{L_{\max}(n)}.
\]
We call $d_\lambda(n)$ the \emph{genuine spectral index} of $n$.
\end{definition}

\begin{theorem}[Spectral inheritance]
\label{thm:inheritance}
For an odd composite $n$,
\[
\Spec_L(n)
:=
\{D:n\in\Eset(D)\}
=
\{d_\lambda(n)r:r\mid L_{\max}(n)\}.
\]
In particular, every composite has one unique genuine spectral index,
namely the least index at which it occurs.
\end{theorem}

\begin{proof}
Every admissible order is a positive divisor of
$L_{\max}(n)$.  If
\[
L=\frac{L_{\max}(n)}r,
\qquad r\mid L_{\max}(n),
\]
then
\[
\frac{n-1}{L}
=
d_\lambda(n)r.
\]
Conversely, Lemma~\ref{lem:divisor-exponent} shows that every divisor
$L$ of $L_{\max}(n)$ occurs as the order of some unit modulo $n$.
\end{proof}

This theorem explains why the same pseudoprime reappears at higher
spectral indices.  The later occurrences are inherited from its genuine
line.

\begin{corollary}[The genuine spectrum partitions the odd composites]
\label{cor:genuine-partition}
Let
\[
\mathcal G(d)=\{n\text{ odd composite}:d_\lambda(n)=d\}.
\]
Then
\[
\boxed{
\{n\ge3:n\text{ odd composite}\}
=\bigsqcup_{d\ge1}\mathcal G(d).
}
\]
Thus every odd composite occurs somewhere in the genuine spectrum and occurs
there exactly once.
\end{corollary}

\begin{proof}
For every odd composite $n$, the integer $L_{\max}(n)=\gcd(n-1,\lambda(n))$
is positive and divides $n-1$, so $d_\lambda(n)$ is a positive integer.
Theorem~\ref{thm:inheritance} shows that it is the unique least spectral index
of $n$.  Hence $n$ belongs to exactly one set $\mathcal G(d)$.
\end{proof}

This partition property is one reason the genuine spectrum is more than a
list of exceptional pseudoprimes.  It supplies a canonical base-independent
coordinate to every odd composite integer.  The inclusive convention is
essential: the genuine witnessing order can be $1$ or $2$.

\section{A hierarchy of $L$-, $\lambda$-, and $\phi$-spectra}

The same denominator viewpoint can be applied directly to arithmetic
functions.

\begin{definition}
For $f\in\{\lambda,\phi\}$ define
\[
d_f(n)=\frac{n-1}{\gcd(n-1,f(n))}.
\]
Equivalently, $d_f(n)$ is the denominator of
$f(n)/(n-1)$ when written in lowest terms.
\end{definition}

Then
\[
d_\lambda(n)=\frac{n-1}{L_{\max}(n)}.
\]
Thus the genuine $L$-spectrum is the $\lambda$-spectrum.

Since $\lambda(n)\mid\phi(n)$,
\[
d_\phi(n)\mid d_\lambda(n).
\]
Writing
\[
H(n)=\frac{\phi(n)}{\lambda(n)},
\]
one also has
\[
\frac{d_\lambda(n)}{d_\phi(n)}\mid H(n).
\]

The classical Lehmer problem is the extreme $\phi$-spectral question
\[
\phi(n)\mid n-1.
\]
The original problem and classical restrictions on a hypothetical composite
solution are discussed in \cite{Lehmer1932,Pomerance1977}.  The more general
$\phi$-spectrum therefore contains Lehmer-type
questions as a special case.  In contrast, the $L$- and
$\lambda$-spectra considered here admit finite spectral lines by the
bounds proved below.

\section{Preliminaries for the period--primality toolbox}
\label{sec:toolbox-prelim}

This section records several elementary reductions that will be used both in
the proofs and in the extended toolbox.  They are stated in a form adapted to an
exact period $L=\ord_n(b)$.

\subsection{The repetend as an integer}

Suppose $\gcd(b,n)=1$ and $L=\ord_n(b)$.  If the repetend of $1/n$ in
base $b$ is regarded as an $L$-digit integer $R$ (allowing leading zeroes),
then
\begin{equation}
R=\frac{b^L-1}{n}.
\label{eq:repetend-integer}
\end{equation}
Indeed,
\[
\frac1n=0.\overline{R}_b=\frac{R}{b^L-1}.
\]
This simple identity is the bridge between multiplicative order and Midy's
block-sum phenomena.

\subsection{A reduction of Somer}

For
\[
n=\prod_{i=1}^s p_i^{a_i},
\]
define
\begin{equation}
\lambda_0(n):=\operatorname{lcm}_{p\mid n}(p-1)
 =\operatorname{lcm}(p_1-1,\ldots,p_s-1).
\label{eq:lambda0}
\end{equation}
This is Somer's $\lambda'(n)$.

\begin{lemma}[Somer's order reduction]
\label{lem:somer-lambda0}
If
\[
L=\ord_n(b)\mid n-1,
\]
then
\[
\gcd(n,L)=1,
\qquad
L\mid\lambda_0(n).
\]
In particular, if $n-1=dL$, then also $\gcd(n,d)=1$.
\end{lemma}

\begin{proof}
The coprimality assertions follow from $L\mid n-1$ and $d\mid n-1$.
For each prime power $p_i^{a_i}\Vert n$, the local order
$\ord_{p_i^{a_i}}(b)$ divides $p_i^{a_i-1}(p_i-1)$ and also divides $L$.
Since $L$ is coprime to $p_i$, the local order actually divides $p_i-1$.
Taking the least common multiple over $i$ gives
$L\mid\lambda_0(n)$.  This is essentially Lemma~1 of Somer
\cite{Somer1989}.
\end{proof}

\begin{remark}
The standard relation $L\mid\lambda(n)$ is valid without the assumption
$L\mid n-1$.  Lemma~\ref{lem:somer-lambda0} is stronger precisely in the
spectral situation: prime-power factors coming from $p_i^{a_i-1}$ cannot
contribute to $L$ because $L$ is coprime to $n$.
\end{remark}

\subsection{Block splitting and a cyclotomic criterion}

Let $r>1$ divide $L$ and put $h=L/r$.  Split the $L$ digits of the
repetend into $r$ consecutive blocks of $h$ digits, represented by integers
$A_0,\ldots,A_{r-1}$ with $0\le A_i<b^h$.  Since $b^h\equiv1\pmod{b^h-1}$,
\[
R\equiv A_0+\cdots+A_{r-1}\pmod{b^h-1}.
\]
Combining this with \eqref{eq:repetend-integer} gives the following useful
special form of Midy's property.  For broader treatments of Midy-type block
criteria and their order-theoretic formulations, see
\cite{GuptaSury2005,Lewittes2007,Martin2007,GarciaGiraldo2009}.

\begin{proposition}[Midy block criterion for $1/n$]
\label{prop:midy-block}
Let $r\mid L$, $r>1$, and $h=L/r$.  The sum of the $r$ equal-length
blocks is divisible by $b^h-1$ if and only if
\begin{equation}
n\mid 1+b^h+b^{2h}+\cdots+b^{(r-1)h}.
\label{eq:midy-geometric}
\end{equation}
If $r=q$ is prime, this is
\[
n\mid\Phi_q(b^h).
\]
For prime $n=p$, condition \eqref{eq:midy-geometric} holds for every
$r>1$ dividing $L$.
\end{proposition}

\begin{proof}
By \eqref{eq:repetend-integer},
\[
R=\frac{(b^h-1)(1+b^h+\cdots+b^{(r-1)h})}{n}.
\]
Thus $b^h-1\mid R$ exactly when $n$ divides the geometric sum.  If $n=p$
is prime, then $b^h$ has exact order $r$ modulo $p$; hence the geometric
sum vanishes modulo $p$.
\end{proof}

For $r=2$, the two blocks have sum exactly $b^h-1$ whenever the Midy
condition holds, because their sum is a positive multiple of $b^h-1$ but
is strictly less than $2(b^h-1)$.  Moreover,
\begin{equation}
\text{two-block Midy property}
\quad\Longleftrightarrow\quad
b^{L/2}\equiv-1\pmod n.
\label{eq:midy-minus1}
\end{equation}

\subsection{Two prime-only tests for the total Midy property}

The phrase \emph{total Midy property} means that the block condition of
Proposition~\ref{prop:midy-block} holds for every divisor $r>1$ of the exact
period $L$.  At first sight this asks for many block divisors and, in the
order-theoretic formulation, for every nontrivial divisor of the modulus.
Both requirements admit a prime-only reduction.

\begin{proposition}[Prime-only characterizations of total Midy]
\label{prop:total-midy-prime-only}
Let $N>1$ be odd, $\gcd(N,b)=1$, and $L=\ord_N(b)$.  The following are
equivalent.
\begin{enumerate}[label=(\roman*)]
\item the Midy block condition holds for every $r>1$ dividing $L$;
\item it is enough to check the block condition only for every \emph{prime}
      $q\mid L$;
\item
\[
\ord_p(b)=L\qquad\text{for every prime }p\mid N;
\]
\item
\[
\ord_D(b)=L\qquad\text{for every divisor }D>1\text{ of }N.
\]
\end{enumerate}
Thus total Midy can be checked on the prime divisors of the period, or,
equivalently, on the prime divisors of the modulus; there is no need to check
all divisors in either family.
\end{proposition}

\begin{proof}
The upward-closure property of the Midy set says that if $r_1\mid r_2\mid L$
and the property holds for $r_1$, then it holds for $r_2$.  Hence checking
all prime $q\mid L$ already checks every divisor $r>1$ of $L$.  The
equivalence of (i) and (iii) is Theorem~2.10 of
\cite{Castillo2015} (with the same result developed from the Midy-set
characterizations in the preceding results of that paper).

It remains only to observe that (iii) and (iv) are elementarily equivalent.
If $p^a\Vert N$, reduction modulo $p$ gives
\[
\ord_p(b)\mid\ord_{p^a}(b),
\]
whereas $p^a\mid N$ gives
\[
\ord_{p^a}(b)\mid\ord_N(b)=L.
\]
Thus $\ord_p(b)=L$ forces $\ord_{p^a}(b)=L$.  The order modulo an arbitrary
divisor $D>1$ is the least common multiple of its prime-power local orders,
so it is again $L$.  The converse follows by taking $D=p$.
\end{proof}

A caution is useful.  Merely checking that the orders modulo the prime
divisors are equal \emph{to one another} is not enough; they must equal the
global order $L=\ord_N(b)$.  For example,
\[
N=9,\qquad b=2,
\]
has the single prime divisor $3$, with $\ord_3(2)=2$, but
$\ord_9(2)=6$, so $9$ is not total Midy to base $2$.

\subsection{A $q$-Midy refinement of the Pocklington mechanism}

Midy's property can also force only a \emph{part} of a prime power in
$N-1$ into every prime divisor of $N$.  This is valuable precisely when the
usual Pocklington gcd for the full prime power fails.

\begin{proposition}[$q$-Midy congruence forcing; Castillo--Garc\'ia-Pulgar\'in--Vel\'asquez-Soto]
\label{prop:q-midy-forcing}
Let $N>1$ be odd, $\gcd(N,b)=1$, and write
\[
N-1=q^s t,
\qquad q\text{ prime},\qquad q\nmid t.
\]
If, for some $0\le i<s$,
\begin{equation}
N\mid\Phi_q\!\left(b^{q^i t}\right),
\label{eq:q-midy-force}
\end{equation}
then every prime divisor $p\mid N$ satisfies
\[
p\equiv1\pmod{q^{i+1}}.
\]
In particular, if
\begin{equation}
2(i+1)>s+\log_q t,
\label{eq:q-midy-prime-ineq}
\end{equation}
then $N$ is prime.
\end{proposition}

\begin{proof}
This is Theorem~3.10 of \cite{Castillo2015}; the short order argument is worth
recording.  Put $x=b^{q^i t}$.  If $p\mid N$, then $p\ne q$ because
$q\mid N-1$.  From $p\mid\Phi_q(x)$ one has $x\not\equiv1\pmod p$ (otherwise
$p\mid q$), while $x^q\equiv1\pmod p$.  Hence $x$ has exact order $q$
modulo $p$.  If $m=\ord_p(b)$, then $m\mid q^{i+1}t$ because $x^q=1$, but
$m\nmid q^it$ because $x\ne1$.  Since $q\nmid t$, this says precisely that
$v_q(m)=i+1$.  Therefore $q^{i+1}\mid m\mid p-1$.

If $N$ were composite, it would have a prime divisor $p\le\sqrt N$.  The
preceding congruence gives $q^{i+1}\le p-1<\sqrt N$, hence
$q^{2(i+1)}<N$.  On the other hand,
condition~\eqref{eq:q-midy-prime-ineq} gives
$q^{2(i+1)}>q^s t=N-1$; since both sides are integers this implies
$q^{2(i+1)}\ge N$, a contradiction.
\end{proof}

\begin{corollary}[Accumulating forced congruence factors]
\label{cor:forced-modulus}
Suppose independent order or Midy arguments prove, for a collection of prime
powers $q^{e_q}\mid N-1$, that
\[
q^{e_q}\mid p-1\qquad\text{for every prime }p\mid N.
\]
If a baseline order $L$ is already known to divide every $p-1$, define
\[
F_{\rm force}:=\operatorname{lcm}\!\left(L,\prod_q q^{e_q}\right).
\]
Then every prime divisor of $N$ is $1$ modulo $F_{\rm force}$.  Consequently
\[
F_{\rm force}\ge\sqrt N
\quad\Longrightarrow\quad N\text{ is prime}.
\]
\end{corollary}

The square-root threshold is not the only useful one.  The familiar BLS
cube-root discriminant argument depends only on the conclusion that every
prime factor is $1$ modulo a common forced modulus.

\begin{proposition}[Cube-root test for an accumulated forced modulus]
\label{prop:forced-modulus-cuberoot}
Assume that every prime divisor $p\mid N$ satisfies $p\equiv1\pmod F$, where
$F\mid N-1$, and suppose
\[
N^{1/3}<F<\sqrt N.
\]
Write the base-$F$ expansion
\[
N=c_2F^2+c_1F+1,
\qquad 0\le c_1,c_2<F.
\]
If
\[
c_1^2-4c_2
\]
is not a square, then $N$ is prime.
\end{proposition}

\begin{proof}
Suppose $N$ were composite and let $p\le\sqrt N$ be a prime divisor.  Write
\[
p=aF+1,
\qquad
N/p=bF+1,
\qquad a,b\ge1,
\]
using the forced congruence for $p$ and for the complementary divisor.  Then
\[
N=abF^2+(a+b)F+1.
\]
Since $N<F^3$, one has $ab<F$.  Also $a+b<F$: otherwise
$ab\ge a+b-1\ge F-1$, and hence
\[
N\ge(F-1)F^2+F^2+1=F^3+1,
\]
a contradiction.  Therefore the displayed expression is already the base-$F$
expansion, so $c_2=ab$ and $c_1=a+b$.  Consequently
\[
c_1^2-4c_2=(a-b)^2
\]
is a square.  The contrapositive proves the assertion.
\end{proof}

This gives a useful way to combine several auxiliary bases: the original
period supplies one forced modulus, and $q$-Midy tests can add prime-power
information that the original base does not expose.  Crossing $\sqrt N$
gives immediate primality; crossing $N^{1/3}$ may already finish the proof by
Proposition~\ref{prop:forced-modulus-cuberoot}.

\subsection{The part of $n-1$ immediately visible to $N-1$ tests}

Assume now that
\[
n-1=dL,
\qquad L=\ord_n(b).
\]
For each prime power $q^e\Vert n-1$, one has
\[
v_q(L)=e-v_q(d).
\]
Define the \emph{spectrally exposed factor}
\begin{equation}
F_L:=\prod_{\substack{q^e\Vert n-1\\q\nmid d}}q^e.
\label{eq:FL}
\end{equation}
Thus $F_L$ consists exactly of the full prime-power factors of $n-1$ whose
$q$-adic exponent is already present in $L$; equivalently, $F_L$ is the
largest divisor of $n-1$ supported on primes not dividing $d$.

For each prime $q\mid F_L$ put
\begin{equation}
g_q=\gcd\!\left(b^{(n-1)/q}-1,n\right).
\label{eq:gq}
\end{equation}
Since $v_q(L)=v_q(n-1)$, the order $L$ does not divide $(n-1)/q$, so
$g_q<n$.  Therefore either $g_q=1$, or $g_q$ is immediately a proper factor
of $n$.  If all $g_q=1$, the hypotheses of the $N-1$ Pocklington theorem
hold for the completely factored part $F_L$.  This observation is the basis
of the Pocklington and BLS entries in the final toolbox.


\subsection{Why the exact period is substantial input}
\label{subsec:order-complexity}

The implications collected in the final toolbox begin with the exact value
\[
L=\ord_n(b).
\]
This is important logically and computationally: the toolbox is not a
``formula'' which turns an arbitrary integer $n$ into a practical primality
test.  For a general pair $(b,n)$, finding $L$ may be much harder than merely
testing $n$ for primality.

In a generic group, the classical baby-step--giant-step and Pollard-rho order
algorithms use on the order of $L^{1/2}$ group operations.  Sutherland's
generic-order algorithms improve this for many orders, sometimes to roughly
$L^{1/3}$, but no classical polynomial-time algorithm in $\log n$ is known for
the unrestricted problem \cite{Sutherland2007}.  If a factored multiple of
$L$ is already known, order computation becomes much faster: one removes prime
factors successively by modular exponentiation.  For the multiplicative group
modulo an unfactored composite $n$, however, obtaining such an exponent
normally brings one back to integer factorization.

The standard general-purpose classical route is therefore factorization
assisted.  The heuristic complexity of the general number field sieve is
\[
\exp\!\left(
 \left((64/9)^{1/3}+o(1)\right)
 (\log n)^{1/3}(\log\log n)^{2/3}
\right),
\]
with faster constants for suitable special forms \cite{Pomerance1996}.
Conversely, knowledge of multiplicative orders of random units can be used to
factor composites by the classical post-processing at the heart of Shor's
reduction.  Quantum order finding, and hence factoring, is polynomial-time in
the input length \cite{Shor1997}; this does not presently make exact order a
routine classical input.

There is a useful conceptual asymmetry here.  Primality itself is known to be
in deterministic polynomial time \cite{AKS2004}, whereas exact order finding for an arbitrary
unknown composite modulus has no comparable classical polynomial-time method.
Thus one should not first compute $L$ merely in order to feed it into this
toolbox.  The toolbox is most useful when $L$ arises naturally from the
problem.

\subsection{Cyclotomic values: a natural source of known periods}
\label{subsec:cyclotomic-known-L}

Cyclotomic factor tables provide an important example in which the period is
known before primality is decided.

\begin{proposition}[Primitive cyclotomic factors are prime or overpseudoprime]
\label{prop:cyclotomic-over}
Let $a\ge2$, $L\ge1$, and let $N>1$ divide $\Phi_L(a)$.  Suppose that
\[
\gcd(N,aL)=1.
\]
Then every divisor $D>1$ of $N$ satisfies
\[
\ord_D(a)=L.
\]
Consequently $N$ is either prime or an overpseudoprime to base $a$.
Moreover $L\mid N-1$.
\end{proposition}

\begin{proof}
Let $p^e\Vert D$.  Since $p\nmid aL$ and $p\mid\Phi_L(a)$, the order of
$a$ modulo $p$ is exactly $L$.  Because $p^e\mid\Phi_L(a)$, one also has
$a^L\equiv1\pmod{p^e}$.  The order modulo $p^e$ is therefore both a
multiple of the order modulo $p$ and a divisor of $L$, hence equals $L$.
Taking least common multiples over the prime-power factors of $D$ gives
$\ord_D(a)=L$.  In particular every prime divisor $p$ of $N$ satisfies
$p\equiv1\pmod L$, and hence $N\equiv1\pmod L$.
\end{proof}

This proposition explains both the usefulness and the limitation of exact
period information in tables such as Kamada's factorizations of
$\Phi_L(10)$ \cite{KamadaPhi10}.  Once exceptional factors dividing $10L$
are removed, an unresolved cofactor is already known to have exactly the
same order $L$ on every nontrivial divisor.  Fermat, Midy, strong
pseudoprimality and overpseudoprimality therefore cannot distinguish a prime
cofactor from a composite one.

Within the present toolbox, the positive primality mechanism that can still
use $L$ is the $N-1$ machinery.  If $N-1=dL$ and $F_L$ is the spectrally
exposed factor of \eqref{eq:FL}, then for primitive cyclotomic factors the
Pocklington gcd conditions are automatic: for $q\mid F_L$ the exponent
$(N-1)/q$ is not divisible by $L$, so it is nontrivial modulo every prime
divisor of $N$.  Thus
\[
F_L\ge\sqrt N
\]
gives a Pocklington certificate, while the BLS cube-root regime may apply
when $F_L^3>N$ together with the usual nonsquare-discriminant condition.
Since $F_L\le L$, these conditions already show why only relatively small
factors of a large cyclotomic value can usually be certified from the period
alone.

\begin{example}[One line of Kamada's base-$10$ table]
Kamada's entry for $L=71401$ begins
\[
\Phi_{71401}(10)
=
856813\cdot195152311403370773\cdot C,
\]
where $C$ is an unresolved cofactor with $64877$ decimal digits.
Here
\[
71401=11\cdot6491.
\]
For the factor $N=856813$,
\[
N-1=12\cdot71401,
\qquad
F_L=71401,
\]
and
\[
71401^2=5\,098\,102\,801>N.
\]
Thus the known order $71401$ supplies a Pocklington proof of primality
using base $10$.  By contrast,
\[
195152311403370773>71401^3
 =364\,009\,638\,094\,201,
\]
so even the clean BLS cube-root threshold cannot be reached using only the
known period.  The enormous cofactor is farther outside the range still.
This is typical: the order information is perfect for recognizing the
cyclotomic structure but, except for small factors, does not by itself settle
primality.
\end{example}

\subsection{Beyond the period-only threshold: a qualification and proof ladder}
\label{subsec:cyclotomic-beyond-threshold}

Suppose now that $N$ is a probable-prime factor of a primitive cyclotomic
value $\Phi_L(a)$, with $\gcd(N,aL)=1$, but that $N$ is too large for the
factor $L$ (or $F_L$) by itself to reach the Pocklington or BLS thresholds.
The first conclusion should be negative but precise: \emph{no further test
that only rechecks the order of the same base $a$ can distinguish prime from
composite}.  If $N$ is composite, Proposition~\ref{prop:cyclotomic-over}
already forces every nontrivial divisor to have that same order.  Thus a
Fermat test, a strong test, or total Midy to the cyclotomic base is already
built into the hypothesis.

The period-only threshold is nevertheless not a ceiling for the rest of the
toolbox.  There are four natural next steps.

\medskip
\noindent\textbf{1. Enlarge the $N-1$ certificate with auxiliary bases.}
Write $N-1=dL$.  The factor $F_L$ is only the portion of $N-1$ automatically
certified by the original base.  It is not the largest factor that
Pocklington can use.  If $q^e\Vert N-1$ is found by factoring an additional
part of $d$, one may search for an auxiliary base $u$ satisfying
\[
u^{N-1}\equiv1\pmod N,
\qquad
\gcd\!\left(u^{(N-1)/q}-1,N\right)=1.
\]
Then every prime $p\mid N$ has $q^e\mid p-1$.  Different primes $q$ may use
different witnesses $u$.  Accumulating these prime powers enlarges the
certified factor of $N-1$; if it reaches $\sqrt N$, Pocklington finishes the
proof, and in the cube-root range the BLS criteria may finish it.  If one of
the gcds lies strictly between $1$ and $N$, the same computation has instead
found a proper factor.

\medskip
\noindent\textbf{2. Use partial $q$-adic forcing when full Pocklington fails.}
Proposition~\ref{prop:q-midy-forcing} is useful when an auxiliary base does
not expose the full $q^e\Vert N-1$.  A congruence
\[
N\mid\Phi_q\!\left(u^{q^i t}\right),
\qquad N-1=q^s t,
\]
still forces $q^{i+1}\mid p-1$ for every $p\mid N$.  Such partial powers can
be combined through Corollary~\ref{cor:forced-modulus}; once the accumulated
modulus exceeds $N^{1/3}$, Proposition~\ref{prop:forced-modulus-cuberoot}
may already finish the proof.  For the original
cyclotomic base $a$ this generally recovers only the $q$-part already present
in $L$; genuine amplification therefore requires auxiliary bases or
additional information.

\medskip
\noindent\textbf{3. Do not restrict BLS to $N-1$.}
The original Brillhart--Lehmer--Selfridge paper develops three families of
criteria: tests using factors of $N-1$, Lucas-sequence tests using factors of
$N+1$, and mixed tests using factors of both $N-1$ and $N+1$
\cite{BLS1975}.  Its stronger $N-1$ results also use lower bounds for the
prime factors of the still-unfactored remainder, so trial division or ECM
progress can itself strengthen a certificate even before another factor is
found.  Thus a stubborn factorization of $N-1$ is a reason both to exploit
such remainder bounds and to inspect $N+1$; the simplified cube-root entry in
our toolbox is only one member of the full BLS family.

\medskip
\noindent\textbf{4. Separate qualification from proof.}
If the aim is first to increase confidence that a large unresolved factor is
prime, use strong probable-prime tests to bases \emph{independent} of the
cyclotomic base.  For an odd composite, at most one quarter of the admissible
bases are strong Miller--Rabin liars, so $k$ independent uniformly random
strong tests have false-pass probability at most $4^{-k}$ with respect to the
random choice of bases; this is a statement about the randomized test, not a
Bayesian prior on $N$ \cite{CrandallPomerance}.  The original base $a$ must
not be counted among these independent tests: an overpseudoprime is forced to
pass it.  Strong Lucas/BPSW-type screening is another useful independent
compositeness filter, but passing such filters is still not a proof.

When a proof rather than a qualification is required and the $N\pm1$
machinery remains below threshold, one should switch to a general primality
prover.  APR--CL/Jacobi-sum methods descend from the algorithm of
Adleman--Pomerance--Rumely \cite{APR1983,CohenLenstra1984}, while ECPP supplies
compact elliptic-curve certificates \cite{AtkinMorain1993}.  These methods do
not exploit the known reciprocal period directly; their role is precisely to
finish cases in which the special structure has delivered all the information
it can.

Kamada maintains a separate list of probable-prime factors whose primality
has not yet been certified \cite{KamadaPRP}.  Such a ``p'' entry is different
from an unresolved cofactor already known to be composite: the former is a
candidate for the qualification/proof ladder above, whereas the latter needs
factorization rather than a primality proof.  As of August 2026 the current
list begins with a $3747$-digit probable-prime factor of
$\Phi_{12198}(10)$.  In such an example the five-digit period is enormously
below the cube-root scale of the candidate, so the period-only certificate is
not the realistic route; auxiliary $N\pm1$ information or a general
certificate method is needed.


\section{Extended period--primality toolbox}
\label{sec:toolbox}

This section is intentionally synthetic.  It is meant to be read as a
reference sheet rather than as a chain of proofs.  Every small tool is
accompanied by an elementary example.  Most entries are classical; the point
is to place them next to one another in the special notation
\[
L=L_b(n)=\ord_n(b).
\]

\begin{center}
\fbox{\parbox{0.91\textwidth}{
\textbf{Important warning.}
The toolbox assumes that the exact order $L$ is already known.  It is not a
general-purpose primality test from $n$ alone.  Generic classical order
computation is expensive, and for an unfactored composite modulus the
factorization-assisted route can be harder than primality testing itself;
see Subsection~\ref{subsec:order-complexity}.  The toolbox is most useful
when $L$ is supplied by the structure of the problem.
}}
\end{center}

\subsection{Fundamental order facts}

Let $n>1$, $b\ge2$, and $\gcd(b,n)=1$.

\begin{enumerate}[label=\textbf{T\arabic*.},leftmargin=2.2em]
\item \textbf{Period equals order.}
The period of $1/n$ in base $b$ is
\[
L=\ord_n(b),
\qquad
L\mid\lambda(n)\mid\phi(n).
\]

\smallskip
\noindent\emph{Example.}
\[
\frac17=0.\overline{142857},
\qquad
\ord_7(10)=6=\lambda(7)=\phi(7).
\]

\item \textbf{Every divisor of $\lambda(n)$ occurs.}
Because $(\mathbb Z/n\mathbb Z)^\times$ is finite abelian, every divisor of
its exponent $\lambda(n)$ is the order of some unit.  In particular
\[
\max_{\gcd(b,n)=1}L_b(n)=\lambda(n).
\]

\smallskip
\noindent\emph{Example.}
For $n=15$, $\lambda(15)=4$.  The divisors $1,2,4$ occur as orders:
\[
\ord_{15}(16)=1,\qquad
\ord_{15}(14)=2,\qquad
\ord_{15}(2)=4.
\]

\item \textbf{Full totient order.}
For odd $n$, an equality $L=\phi(n)$ is possible exactly when
$(\mathbb Z/n\mathbb Z)^\times$ is cyclic; hence $n=p^k$ for an odd prime
$p$.  Thus $L=\phi(n)$ forces an odd prime power, but not necessarily a
prime.

\smallskip
\noindent\emph{Example.}
\[
\phi(9)=6,\qquad \ord_9(2)=6,
\]
yet $9$ is composite.

\item \textbf{Spectral condition equals Fermat's congruence.}
\[
L\mid n-1
\quad\Longleftrightarrow\quad
b^{n-1}\equiv1\pmod n.
\]
For composite $n$ this is exactly Fermat pseudoprimality to base $b$.

\smallskip
\noindent\emph{Example.}
For $n=341=11\cdot31$ and $b=2$,
\[
L=\ord_{341}(2)=10\mid340,
\]
so $2^{340}\equiv1\pmod{341}$.
Carmichael numbers are the extreme all-base version of Fermat pseudoprimality;
their infinitude was proved by Alford, Granville and Pomerance
\cite{AGP1994}.

\item \textbf{Somer reduction under the spectral condition.}
If $L\mid n-1$, then
\[
L\mid\lambda_0(n)=\operatorname{lcm}_{p\mid n}(p-1).
\]
Prime-power factors of $n$ therefore do not supply extra factors to a
spectral period.

\smallskip
\noindent\emph{Example.}
For $n=25$, $b=7$,
\[
L=\ord_{25}(7)=4,
\qquad
\lambda_0(25)=5-1=4,
\]
even though $\lambda(25)=20$.
\end{enumerate}

\subsection{Direct primality and pseudoprimality consequences}

\begin{enumerate}[label=\textbf{T\arabic*.},leftmargin=2.2em,resume]
\item \textbf{Lucas criterion.}
For odd $n>1$,
\[
\exists b:\ \ord_n(b)=n-1
\quad\Longleftrightarrow\quad
n\text{ is prime}.
\]

\smallskip
\noindent\emph{Example.}
\[
\ord_7(3)=6=7-1,
\]
so Lucas' criterion certifies $7$.  Recursive primality certificates based
on a primitive-root witness and the prime factors of $p-1$ lead to the
classical Pratt certificate framework \cite{Pratt1975}.

\item \textbf{Small spectral index.}
Somer proved fixed-index finiteness.  The present paper adds explicit size
obstructions and exact low-frequency lines.

\smallskip
\noindent\emph{Example.}
At $d=4$ the complete spectral line is
\[
\mathcal E(4)=\{9\}.
\]
Thus any odd $n>9$ satisfying
\[
\ord_n(b)=\frac{n-1}{4}
\]
is prime.

\item \textbf{Repeated prime powers.}
In the notation
\[
\kappa=\lambda(n)/L,
\qquad
K=n/\rad(n),
\]
one has
\[
K\mid\kappa.
\]
Thus a prime divisor $p>\kappa$ occurs only to the first power.

\smallskip
\noindent\emph{Example.}
For the exceptional point $n=25$, $d=6$, $b=7$,
\[
L=4,\qquad \kappa=\frac{20}{4}=5,\qquad
K=\frac{25}{5}=5,
\]
and indeed $K\mid\kappa$.

\item \textbf{Number of distinct prime factors.}
If $s=\omega(n)$ and $n$ is an odd composite spectral point, then
\[
2^{s-1}\kappa\le d-1,
\qquad
s\le1+\left\lfloor\log_2\frac{d-1}{\kappa}\right\rfloor.
\]

\smallskip
\noindent\emph{Example.}
For $n=561$, one may take $b=5$ and $d=7$:
\[
L=80,\qquad \lambda(561)=80,\qquad \kappa=1,\qquad s=3.
\]
The inequality reads
\[
2^{3-1}=4\le6=d-1.
\]
\end{enumerate}

\subsection{Midy's property from the exact period}

Let $r>1$ divide $L$ and $h=L/r$.

\begin{enumerate}[label=\textbf{T\arabic*.},leftmargin=2.2em,resume]
\item \textbf{Exact block criterion.}
The sum of the $r$ equal-length repetend blocks is divisible by $b^h-1$
exactly when
\[
n\mid1+b^h+\cdots+b^{(r-1)h}.
\]
For prime $r=q$ this is
\[
n\mid\Phi_q(b^h).
\]

\smallskip
\noindent\emph{Example.}
For $1/7$ in base $10$, take $r=2$, $h=3$:
\[
1+10^3=1001=7\cdot143.
\]
Correspondingly,
\[
142+857=999=10^3-1.
\]

\item \textbf{Prime Midy property.}
If $n=p$ is prime, the Midy condition holds for every divisor $r>1$ of
$L$.

\smallskip
\noindent\emph{Example.}
Again $1/7=0.\overline{142857}$.  Taking $r=3$, $h=2$ gives
\[
14+28+57=99=10^2-1.
\]

\item \textbf{The nines property.}
If $2\mid L$ and $h=L/2$, the two halves add to $b^h-1$ precisely when
\[
b^h\equiv-1\pmod n.
\]

\smallskip
\noindent\emph{Example.}
For $n=7$, $b=10$,
\[
10^3\equiv-1\pmod7
\]
and the two halves $142$ and $857$ add to $999$.

\item \textbf{Two-block Midy implies strong pseudoprimality in the spectral
case.}
If $n$ is composite, $L\mid n-1$, and
\[
b^{L/2}\equiv-1\pmod n,
\]
then $n$ passes the Miller--Rabin strong pseudoprime test to base $b$.

\smallskip
\noindent\emph{Example.}
For
\[
n=3277,\qquad b=2,
\]
one has
\[
L=\ord_{3277}(2)=28,\qquad 28\mid3276,
\qquad
2^{14}\equiv-1\pmod{3277}.
\]
Thus $3277$ is a strong pseudoprime to base $2$.

\item \textbf{Total Midy needs only prime checks.}
The following conditions are equivalent:
\[
\begin{array}{c}
\text{Midy holds for every }r>1\text{ dividing }L,\\[1mm]
\text{Midy holds for every prime }q\mid L,\\[1mm]
\ord_p(b)=L\text{ for every prime }p\mid n,\\[1mm]
\ord_D(b)=L\text{ for every divisor }D>1\text{ of }n.
\end{array}
\]
Thus neither all block divisors of $L$ nor all divisors of $n$ have to be
checked; see Proposition~\ref{prop:total-midy-prime-only} and
\cite{Castillo2015}.  A composite satisfying these conditions is an
overpseudoprime to base $b$ and is automatically a strong pseudoprime to the
same base \cite{Shevelev2012,Castillo2012}.

\smallskip
\noindent\emph{Example.}
\[
2047=23\cdot89,\qquad
\ord_{23}(2)=\ord_{89}(2)=\ord_{2047}(2)=11.
\]
The two prime checks already prove that $2047$ is total Midy, hence an
overpseudoprime, to base $2$.

\item \textbf{$q$-Midy forcing can refine Pocklington.}
If
\[
n-1=q^s t,\qquad q\nmid t,
\]
and for some $0\le i<s$
\[
n\mid\Phi_q\!\left(b^{q^i t}\right),
\]
then every prime divisor $p\mid n$ satisfies
\[
p\equiv1\pmod{q^{i+1}}.
\]
If $2(i+1)>s+\log_q t$, this alone proves that $n$ is prime; see
Proposition~\ref{prop:q-midy-forcing} and \cite{Castillo2015}.

\smallskip
\noindent\emph{Example.}
For $n=1459$, $n-1=3^6\cdot2$.  Castillo, Garc\'ia-Pulgar\'in and
Vel\'asquez-Soto verify the criterion with $b=2$, $q=3$, and $i=4$; hence
every prime factor would be $1\pmod{3^5}$, which is already too large for a
proper factor not exceeding $\sqrt{1459}$.  Thus $1459$ is prime.
\end{enumerate}

\subsection{$N-1$ primality certificates extracted from $L$}

Suppose $L\mid n-1$, $n-1=dL$, and let $F_L$ be the spectrally exposed
factor of \eqref{eq:FL}.

\begin{enumerate}[label=\textbf{T\arabic*.},leftmargin=2.2em,resume]
\item \textbf{Automatic compositeness witness.}
For each prime $q\mid F_L$ put
\[
g_q=\gcd\!\left(b^{(n-1)/q}-1,n\right).
\]
One always has $g_q<n$; if $g_q>1$, it is a proper factor of $n$.

\smallskip
\noindent\emph{Example.}
For
\[
n=15,\quad b=4,\quad L=2,\quad d=7,
\]
the prime $q=2$ is spectrally exposed and
\[
\gcd(4^7-1,15)=3.
\]
The order information itself produces the factor $3$.

\item \textbf{Pocklington regime.}
If $g_q=1$ for every $q\mid F_L$, every prime divisor $p$ of $n$ satisfies
$p\equiv1\pmod{F_L}$.  Hence
\[
F_L\ge\sqrt n
\quad\Longrightarrow\quad n\text{ is prime}.
\]

\smallskip
\noindent\emph{Example.}
For $n=13$, $b=2$,
\[
L=12,\qquad F_L=12>\sqrt{13}.
\]
The two required checks are
\[
\gcd(2^6-1,13)=1,\qquad
\gcd(2^4-1,13)=1,
\]
so Pocklington certifies $13$.

\item \textbf{Certificate amplification beyond the original period.}
The exposed factor $F_L$ is only what the known base certifies automatically.
If $q^e\Vert n-1$ lies in the remaining cofactor $d=(n-1)/L$, an auxiliary
base $u$ satisfying
\[
u^{n-1}\equiv1\pmod n,
\qquad
\gcd\!\left(u^{(n-1)/q}-1,n\right)=1
\]
forces $q^e\mid p-1$ for every prime $p\mid n$.  Different $q$ may use
different bases.  Such prime powers may therefore be adjoined to the
certified $N-1$ factor until Pocklington or BLS reaches its threshold.  A
nontrivial gcd instead proves compositeness immediately.

\smallskip
\noindent\emph{Example.}
If the original order exposes $F_L=2^a$ but factoring $(n-1)/L$ also reveals
$q^e$, a successful auxiliary witness for $q$ upgrades the usable certified
factor from $2^a$ to $2^a q^e$; the auxiliary base need not be the base whose
reciprocal period supplied $L$.

\item \textbf{BLS cube-root regime.}
If
\[
n^{1/3}<F_L<\sqrt n
\]
and the Pocklington congruences hold, write
\[
n=c_2F_L^2+c_1F_L+1,\qquad0\le c_1,c_2<F_L.
\]
If
\[
c_1^2-4c_2
\]
is not a square, the Brillhart--Lehmer--Selfridge cube-root criterion
proves primality \cite{BLS1975}.  We use only this precise version here;
stronger BLS variants have additional hypotheses and are not folded into the
toolbox statement.  More generally, $F_L$ may be replaced by any accumulated
forced modulus $F$ for which every prime divisor of $n$ is known to be
$1\pmod F$; see Proposition~\ref{prop:forced-modulus-cuberoot}.

\smallskip
\noindent\emph{Example.}
Take $n=71$.  The element $b=20$ has
\[
\ord_{71}(20)=7,\qquad 71-1=10\cdot7,
\]
so $F_L=7$.  Since
\[
71=1\cdot7^2+3\cdot7+1
\]
and
\[
3^2-4\cdot1=5
\]
is not a square, the BLS cube-root criterion certifies $71$.

\item \textbf{The full BLS toolbox also uses $N+1$.}
The preceding entry is only the convenient $N-1$ cube-root form.  The
original BLS theory also contains Lucas-sequence criteria based on factors of
$N+1$ and stronger mixed criteria using factors of both $N-1$ and $N+1$
\cite{BLS1975}.  When the known period leaves the $N-1$ certificate too
small, factoring $N+1$ is therefore a natural second route.

\smallskip
\noindent\emph{Example.}
For a large cyclotomic probable prime with $L^3<N$, failure of the
period-only cube-root test says nothing about the factorization of $N+1$; a
large certified factor there can activate a Lucas or mixed BLS criterion.

\item \textbf{Primitive cyclotomic factors.}
If
\[
N\mid\Phi_L(b),\qquad \gcd(N,bL)=1,
\]
then every nontrivial divisor of $N$ has order $L$.  Thus $N$ is either
prime or overpseudoprime.  All order-only tests in this toolbox therefore
see the two cases in exactly the same way; the $N-1$ certificate is the
tool that can still turn the known $L$ into a positive primality proof.

\smallskip
\noindent\emph{Example.}
Kamada lists
\[
856813\mid\Phi_{71401}(10).
\]
Here
\[
856813-1=12\cdot71401,
\qquad
71401^2>856813.
\]
The primitive cyclotomic structure supplies the Pocklington gcd conditions,
so the known period proves that $856813$ is prime.  On the same Kamada line,
the $18$-digit factor $195152311403370773$ is already larger than
$71401^3$, so the period alone does not reach even the clean BLS cube-root
threshold.

\item \textbf{Independent bases qualify a cyclotomic PRP.}
An overpseudoprime is forced to pass the strong test to its defining
cyclotomic base, so that base supplies no statistical evidence.  Strong
Miller--Rabin tests to independently chosen bases do: an odd composite has at
most one quarter strong-liar bases, giving a false-pass probability at most
$4^{-k}$ for $k$ independent uniformly random choices.  This qualifies a
probable prime; it does not prove one \cite{CrandallPomerance}.

\smallskip
\noindent\emph{Example.}
For a Kamada factor of $\Phi_L(10)$, base $10$ should not be counted among
the independent Miller--Rabin bases, because the cyclotomic structure already
forces the required order behavior.

\item \textbf{General proof fallback: APR--CL or ECPP.}
If certificate amplification, $q$-Midy forcing, and the $N\pm1$ BLS routes
remain below threshold, a general primality prover is the correct final tool.
APR--CL uses Jacobi sums, while ECPP recursively constructs elliptic-curve
primality certificates \cite{APR1983,CohenLenstra1984,AtkinMorain1993}.
These methods do not need the reciprocal period; they finish precisely the
cases in which the special period structure no longer suffices.

\smallskip
\noindent\emph{Example.}
Kamada maintains a separate list of large probable-prime factors whose
primality has not yet been certified \cite{KamadaPRP}; those entries are
natural inputs for a general certificate generator after the inexpensive
structural and probable-prime filters have been exhausted.
\end{enumerate}

\subsection{Spectral distribution and related spectra}

\begin{enumerate}[label=\textbf{T\arabic*.},leftmargin=2.2em,resume]
\item \textbf{Existential prime spectrum.}
Allowing the base to vary,
\[
\mathcal P(d)=\{p\text{ prime}:p\equiv1\pmod d\}.
\]
Each line is infinite and has relative density $1/\phi(d)$ among primes.

\smallskip
\noindent\emph{Example.}
For $d=5$, the prime $31$ belongs to the line because
\[
31\equiv1\pmod5.
\]
Indeed,
\[
\ord_{31}(6)=6=\frac{31-1}{5}.
\]

\item \textbf{Fixed-base prime spectrum.}
For one fixed base $b$, the set
\[
\{p:\ord_p(b)=(p-1)/d\}
\]
is the near-primitive-root problem.  Under GRH it has the
Wagstaff--Moree density described in Section~\ref{sec:fixed-base}.

\smallskip
\noindent\emph{Example.}
For the decimal base,
\[
\ord_{13}(10)=6=\frac{13-1}{2},
\]
so $13$ lies on the fixed-base-$10$ spectral line $d=2$.

\item \textbf{$\lambda$- and $\phi$-spectral hierarchy.}
Define
\[
d_f(n)=\frac{n-1}{\gcd(n-1,f(n))},
\qquad f\in\{\lambda,\phi\}.
\]
Then
\[
d_\phi(n)\mid d_\lambda(n),
\]
and $d_\lambda(n)$ is the unique genuine $L$-spectral index of $n$.

\smallskip
\noindent\emph{Example.}
For $n=21$,
\[
\lambda(21)=6,\qquad \phi(21)=12,\qquad n-1=20.
\]
Hence
\[
d_\lambda(21)=\frac{20}{2}=10,
\qquad
d_\phi(21)=\frac{20}{4}=5,
\]
so indeed $d_\phi(21)\mid d_\lambda(21)$.
\end{enumerate}

\subsection{One-page decision map}

For quick use:
\[
\begin{array}{ccl}
L=n-1
&\Longrightarrow& n\text{ prime},\\[2mm]
L\mid n-1
&\Longleftrightarrow& b^{n-1}\equiv1\pmod n,\\[2mm]
L\mid n-1,\ n\text{ composite}
&\Longrightarrow& \text{Fermat pseudoprime to }b,\\[2mm]
2\mid L,\ \text{two-block Midy}
&\Longrightarrow& \text{strong pseudoprime to }b,\\[2mm]
\ord_p(b)=L\ \forall p\mid n
&\Longleftrightarrow& \text{total Midy / overpseudoprime},\\[2mm]
N\mid\Phi_q(b^{q^it})
&\Longrightarrow& q^{i+1}\mid(p-1)\ \forall p\mid N,\\[2mm]
F_L\ge\sqrt n,\ g_q=1\ \forall q\mid F_L
&\Longrightarrow& n\text{ prime},\\[2mm]
F_{\rm force}\ge\sqrt N
&\Longrightarrow& N\text{ prime},\\[2mm]
F_{\rm force}>N^{1/3},\ \text{nonsquare digit discriminant}
&\Longrightarrow& N\text{ prime},\\[2mm]
n\ge U_0(d),\ L=(n-1)/d
&\Longrightarrow& n\text{ prime}.
\end{array}
\]

The conceptual message is twofold.  First, many classical statements become
especially transparent once the exact reciprocal period $L$ is placed at the
center.  Second, none of this removes the computational cost of obtaining
$L$.  The toolbox is therefore a structural reference and a conditional
certificate generator, not a universal replacement for modern primality
testing.



\section{Structural bounds for composite spectral lines}

Assume throughout this section that $n$ is odd and composite and
\[
L=L_b(n)=\frac{n-1}{d}.
\]
Define
\begin{equation}
\kappa=\frac{\lambda(n)}L,
\qquad
H=\frac{\phi(n)}{\lambda(n)},
\qquad
j=\frac{\phi(n)}L=\kappa H.
\label{eq:dHj}
\end{equation}
Then
\[
d\lambda(n)=\kappa(n-1),
\qquad
d\phi(n)=j(n-1).
\]
Since $\phi(n)<n-1$,
\begin{equation}
1\le \kappa\le j\le d-1.
\label{eq:djrange}
\end{equation}

Write
\[
n=\prod_{i=1}^s p_i^{a_i},
\qquad
3\le p_1<\cdots<p_s,
\qquad
s=\omega(n),
\]
and put
\[
K=\prod_{i=1}^s p_i^{a_i-1}.
\]

\begin{lemma}[Repeated-prime restriction]
\label{lem:Kdivd}
One has
\[
K\mid \kappa.
\]
Consequently, if $a_i>1$, then
\[
p_i^{a_i-1}\mid \kappa.
\]
In particular every prime $p_i>\kappa$ occurs to the first power.
\end{lemma}

\begin{proof}
For odd prime powers,
\[
\lambda(p_i^{a_i})=p_i^{a_i-1}(p_i-1).
\]
Hence $K\mid\lambda(n)$.  From
$d\lambda(n)=\kappa(n-1)$ we get $K\mid \kappa(n-1)$.
But every prime divisor of $K$ divides $n$, and hence
$\gcd(K,n-1)=1$.  Therefore $K\mid \kappa$.
\end{proof}

\begin{lemma}[Two-adic overlap]
\label{lem:H2adic}
For odd $n$ with $s$ distinct prime divisors,
\[
2^{s-1}\mid H=\frac{\phi(n)}{\lambda(n)}.
\]
Consequently
\begin{equation}
2^{s-1}\kappa\mid j,
\qquad
2^{s-1}\kappa\le d-1.
\label{eq:sd-bound}
\end{equation}
\end{lemma}

\begin{proof}
If $e_i=v_2(p_i-1)$, then
\[
v_2(\phi(n))=\sum_i e_i,
\qquad
v_2(\lambda(n))=\max_i e_i.
\]
Since all $e_i\ge1$,
\[
v_2(H)
=
\sum_i e_i-\max_i e_i
\ge s-1.
\]
The second assertion follows from $j=\kappa H\le d-1$.
\end{proof}

\begin{corollary}[Number of distinct prime factors]
\label{cor:sbound}
Every composite on spectral line $d$ satisfies
\[
s
\le
1+\left\lfloor
\log_2\!\left(\frac{d-1}{\kappa}\right)
\right\rfloor
\le
1+\lfloor\log_2(d-1)\rfloor.
\]
\end{corollary}

\begin{proposition}[Smallest-prime bound]
\label{prop:p1}
One has
\[
p_1<
\frac{d+j(s-1)}{d-j},
\]
and therefore
\[
p_1\le s(d-1).
\]
\end{proposition}

\begin{proof}
From $d\phi(n)=j(n-1)$,
\[
\frac dj
=
\frac{n-1}{\phi(n)}
<
\frac n{\phi(n)}
=
\prod_{i=1}^s\frac{p_i}{p_i-1}.
\]
Since the primes are distinct and odd, $p_i\ge p_1+2(i-1)$, and hence a fortiori $p_i\ge p_1+i-1$,
\[
\prod_{i=1}^s\frac{p_i}{p_i-1}
\le
\frac{p_1+s-1}{p_1-1}.
\]
Solving for $p_1$ gives the result.
\end{proof}

\begin{lemma}[A uniform Euler-factor ratio bound]
\label{lem:euler-ratio}
If $3\le p_1<\cdots<p_r$ are distinct odd primes, then
\[
\prod_{i=1}^r\frac{p_i}{p_i-1}
\le
\prod_{i=1}^r\frac{i+2}{i+1}
=
\frac{r+2}{2}
<r+1.
\]
\end{lemma}

\begin{proof}
Since $p_i\ge i+2$ and $x/(x-1)$ is decreasing for $x>1$, each factor is at
most $(i+2)/(i+1)$; the latter product telescopes.
\end{proof}

For
\[
P_k=\prod_{i=1}^k p_i^{a_i},
\]
the following recursive restriction is useful.

\begin{proposition}[Squarefree-suffix bound]
\label{prop:suffix}
If $2\le k\le s$ and $p_k>\kappa$, then
$p_k,\ldots,p_s$ occur only to the first power and
\[
p_k
<
d\lambda(P_{k-1})(s-k+1).
\]
\end{proposition}

\begin{proof}
By Lemma~\ref{lem:Kdivd}, the suffix
\[
q:=\prod_{i=k}^{s}p_i
\]
is squarefree.  Put $P=P_{k-1}$ and
\[
J_k:=\frac{\lambda(P)\phi(q)}{\lambda(n)}.
\]
Since $\lambda(n)=\operatorname{lcm}(\lambda(P),\lambda(q))$ and
$\lambda(q)\mid\phi(q)$, the number $J_k$ is an integer.  Multiplying
\[
d\lambda(n)=\kappa(n-1)
\]
by $J_k$ gives
\[
d\lambda(P)\phi(q)=\kappa J_k(Pq-1).
\]
Write
\[
\phi(q)=q-\Delta_q.
\]
Then
\[
q\bigl(d\lambda(P)-\kappa J_kP\bigr)
=
d\lambda(P)\Delta_q-\kappa J_k.
\]
The coefficient on the left is a positive integer, because
\[
d\lambda(P)-\kappa J_kP
=
d\lambda(P)
\frac{P(q-\phi(q))-1}{Pq-1}>0.
\]
Hence it is at least $1$, and therefore
\[
q<d\lambda(P)\Delta_q.
\]
Finally,
\[
\frac{\Delta_q}{q}
=
1-\prod_{i=k}^{s}\left(1-\frac1{p_i}\right)
<
\sum_{i=k}^{s}\frac1{p_i}
\le\frac{s-k+1}{p_k}.
\]
Combining the last two inequalities proves
\[
p_k<d\lambda(P_{k-1})(s-k+1).
\]
\end{proof}

\begin{proposition}[Largest-prime bound]
\label{prop:largest}
For $s\ge2$,
\[
p_s\le \kappa(P_{s-1}-1)+1.
\]
More precisely, if $p_s>\kappa$, then
\[
p_s-1\mid \kappa(P_{s-1}-1).
\]
\end{proposition}

\begin{proof}
If $p_s\le\kappa$, then the displayed upper bound is immediate because
$P_{s-1}\ge3$.  Suppose therefore that $p_s>\kappa$.  Lemma~\ref{lem:Kdivd}
gives $a_s=1$.  Since $p_s-1\mid\lambda(n)$ and
$d\lambda(n)=\kappa(n-1)$,
\[
p_s-1\mid \kappa(n-1).
\]
Writing $n=P_{s-1}p_s$ and reducing modulo $p_s-1$ gives
\[
p_s-1\mid \kappa(P_{s-1}-1),
\]
and hence $p_s\le\kappa(P_{s-1}-1)+1$.
\end{proof}


\subsection{A uniform spectral size obstruction}
\label{sec:uniform-size}

Let
\[
L=\frac{n-1}{d},\qquad
\kappa=\frac{\lambda(n)}{L},\qquad
H=\frac{\phi(n)}{\lambda(n)},\qquad
j=\kappa H.
\]
For an odd composite $n$ with $s=\omega(n)$ distinct prime factors,
\[
2^{s-1}\mid H,\qquad j\le d-1,
\]
and hence
\[
2^{s-1}\kappa\le d-1.
\tag{U1}\label{eq:u1}
\]

For $d\ge2$, define
\[
S=S(d):=1+\lfloor\log_2(d-1)\rfloor,
\qquad
M=M(d):=2^{S-1}.
\]
Thus $M$ is the largest power of $2$ not exceeding $d-1$.

\begin{theorem}[Uniform spectral size obstruction]
\label{thm:uniform-U}
Suppose that $n>1$ is odd and
\[
L_b(n)=\frac{n-1}{d}
\]
for some $b\ge2$ coprime to $n$.  If $n$ is composite, then
\[
\boxed{
 n< U_0(d):=\bigl(dS(d)\bigr)^{\,2M(d)} .
}
\tag{U2}\label{eq:U0}
\]
Consequently
\[
n\ge U_0(d)\quad\Longrightarrow\quad n\text{ is prime}.
\]
Here $d\ge2$; for $d=1$ the conclusion is Lucas's criterion.
\end{theorem}

\begin{proof}
The case $d=1$ is immediate from $L_b(n)=n-1$, so assume $d\ge2$.
Write
\[
n=K\prod_{i=1}^{s}p_i,
\qquad
K=\prod_{i=1}^{s}p_i^{a_i-1},
\qquad
3\le p_1<\cdots<p_s.
\]
The repeated-prime restriction gives $K\mid\kappa$.  In particular,
\[
n\le \kappa R_s,
\qquad
R_k:=\prod_{i=1}^{k}p_i,
\tag{U3}\label{eq:radical}
\]
and $s\le S$ by \eqref{eq:u1}.

We first dispose of the prime-power case $s=1$, which is not covered by the
terminal-prime estimate below.  Then $n=p^a$ with $a\ge2$ and
$K=p^{a-1}\mid\kappa\le d-1$.  Since $p\le p^{a-1}=K$,
\[
n=pK\le K^2\le(d-1)^2<d^2\le U_0(d).
\]
Henceforth assume $s\ge2$.

For $2\le k\le s-1$, put
\[
A_{k-1}:=\prod_{i<k}(p_i-1),
\qquad
K_{k-1}:=\prod_{i<k}p_i^{a_i-1}.
\]
Then
$P_{k-1}=K_{k-1}\prod_{i<k}p_i$ has $k-1$ distinct odd prime factors and
$K_{k-1}\mid K\mid\kappa$.  Therefore
\[
\lambda(P_{k-1})
 \le \frac{\phi(P_{k-1})}{2^{k-2}}
 = \frac{K_{k-1}A_{k-1}}{2^{k-2}}
 \le \frac{\kappa A_{k-1}}{2^{k-2}}.
\]
The squarefree-suffix estimate yields
\[
p_k<
\frac{d\kappa(s-k+1)}{2^{k-2}}A_{k-1}
\qquad\text{whenever }p_k>\kappa.
\tag{U4}\label{eq:pk-sharp}
\]
If $p_k\le\kappa$, the weaker inequality
\[
p_k<d\kappa(s-k+1)A_{k-1}
\]
is automatic.  Thus, for all $2\le k\le s-1$,
\[
A_k=A_{k-1}(p_k-1)
<d\kappa S\,A_{k-1}^{\,2}.
\tag{U5}\label{eq:A-rec}
\]

We distinguish two cases.

\smallskip
\noindent\emph{Case 1: $s=S$.}
Since
\[
M=2^{S-1}\le d-1<2M,
\]
\eqref{eq:u1} forces $\kappa=1$.  Moreover $H$ is a positive multiple of
$M$ and $j=H\le d-1<2M$, hence
\[
H=j=M.
\]
The smallest-prime bound gives
\[
p_1<\frac{d+M(S-1)}{d-M}<dS,
\]
so $A_1<dS$.  Iterating \eqref{eq:A-rec} gives
\[
A_{S-1}<(dS)^{M-1}.
\tag{U6}\label{eq:A-top}
\]
By Lemma~\ref{lem:euler-ratio},
\[
\frac{R_{S-1}}{A_{S-1}}
=\prod_{i=1}^{S-1}\frac{p_i}{p_i-1}
\le\frac{S+1}{2}<S,
\]
and therefore
\[
R_{S-1}<S(dS)^{M-1}.
\]
Here $\kappa=K=1$, so $n$ is squarefree.  The terminal-prime divisibility is
\[
p_S-1\mid R_{S-1}-1.
\]
Because $S\ge2$, one has $R_{S-1}-1>0$, hence
$p_S\le R_{S-1}$.  Consequently
\[
n=R_{S-1}p_S
\le R_{S-1}^{\,2}
<S^2(dS)^{2M-2}
<(dS)^{2M}.
\]

\smallskip
\noindent\emph{Case 2: $2\le s<S$.}
Put $M_s=2^{s-1}$, so $M_s\le M/2$.  Using only
$\kappa<d$, $s\le S$, and $p_1<dS$, equation \eqref{eq:A-rec} gives
\[
A_{s-1}<(d^2S)^{M_s-1}.
\]
Lemma~\ref{lem:euler-ratio} gives
\[
R_{s-1}<S(d^2S)^{M_s-1}.
\tag{U7}\label{eq:R-low}
\]
The terminal-prime bound and $P_{s-1}\le\kappa R_{s-1}$ imply
\[
p_s\le\kappa(P_{s-1}-1)+1\le\kappa P_{s-1}
\le\kappa^2R_{s-1}.
\]
Using \eqref{eq:radical}, $\kappa<d$, \eqref{eq:R-low}, and $2M_s\le M$,
\[
n<d^3R_{s-1}^{\,2}
<d^3S^2(d^2S)^{2M_s-2}
\le d^{2M-1}S^M
<(dS)^{2M}.
\]
This completes the proof.
\end{proof}

\begin{corollary}[Asymptotic primality criterion]
\label{cor:asymptotic-U}
For composite spectral points,
\[
\log n
\le
2d\bigl(\log d+\log\log d+O(1)\bigr).
\]
Equivalently,
\[
d\ge
\left(\frac12+o(1)\right)
\frac{\log n}{\log\log n}.
\]
Thus for every fixed $\varepsilon>0$, all sufficiently large $n$ satisfying
\[
d\le
\left(\frac12-\varepsilon\right)
\frac{\log n}{\log\log n}
\]
are prime.
\end{corollary}

\begin{proof}
Since $M(d)\le d-1$ and $S(d)=O(\log d)$, Theorem~\ref{thm:uniform-U}
gives
\[
\log U_0(d)
\le
2d\bigl(\log d+\log\log d+O(1)\bigr).
\]
Fix $\varepsilon>0$ and suppose that
\[
d\le\left(\frac12-\varepsilon\right)
\frac{\log n}{\log\log n}.
\]
Only upper bounds on $d$ are needed.  Uniformly in this range,
\[
\log d\le \log\log n+O(\log\log\log n),
\qquad
\log\log d=O(\log\log\log n)
\]
whenever $d\to\infty$; bounded $d$ are easier and satisfy the same final
estimate.  Therefore
\begin{align*}
\log U_0(d)
&\le 2d\bigl(\log d+\log\log d+O(1)\bigr)\\
&\le (1-2\varepsilon)
   \frac{\log n}{\log\log n}
   \bigl(\log\log n+O(\log\log\log n)\bigr)\\
&=(1-2\varepsilon+o(1))\log n<\log n
\end{align*}
for all sufficiently large $n$.  Hence $n>U_0(d)$, so
Theorem~\ref{thm:uniform-U} forces $n$ to be prime.  This proves the stated
criterion without assuming a lower bound on $d$.
\end{proof}

\begin{remark}
The coefficient $1/2$ is the worst leading coefficient produced by the
uniform bound, rather than a claimed sharp asymptotic constant for actual
composite spectral points.  It is governed by the dyadic extremal values
\[
M(d)=d-1
\quad\Longleftrightarrow\quad
d=2^r+1.
\]
At $d=2^r$ one has $M(d)=d/2$, and the corresponding local leading coefficient
in this bound is $1$ rather than $1/2$.  Retaining the exact branch parameters
$(s,\kappa,H,j)$ gives much smaller finite bounds away from the dyadic spikes.
\end{remark}

\begin{remark}[Comparison with Somer's effective finiteness]
Somer's proof of fixed-$d$ finiteness is effective in the sense that, for
each prescribed $d$, it eventually produces a computable constant beyond
which no Fermat $d$-pseudoprime can occur.  He explicitly remarks, however,
that his argument does not provide a uniform procedure for obtaining this
constant as a function of arbitrary $d$ \cite{Somer1989}.  The point of
$U_0(d)$ is therefore not finiteness itself, which is Somer's theorem, but
uniformity: it is one closed expression valid for every $d$.  The asymptotic
consequence
\[
d\ge\left(\frac12+o(1)\right)\frac{\log n}{\log\log n}
\]
for composite spectral points is a quantitative form of the statement that
small indices are rare.
\end{remark}


\section{Somer's small-index theorem and a unified three-page proof}
\label{sec:small-index}

Somer's Theorem~3 \cite{Somer1989} gives a complete table of his Fermat
$d$-pseudoprimes for $1\le d\le8$.  With his convention excluding bases
$a\equiv\pm1\pmod N$, the table is
\[
\begin{array}{c|c}
 d&\text{composite Fermat $d$-pseudoprimes}\\ \hline
1,2,3,4&\varnothing\\
5&6601\\
6&25\\
7&15,561\\
8&49\\
\end{array}
\]
The present spectral convention allows every integer base $b\ge2$ coprime
to $n$.  It therefore adds the order-two occurrence
$\ord_9(8)=2=(9-1)/4$ and the order-one occurrence
$\ord_9(10)=1=(9-1)/8$.

Somer's published paper gives a detailed proof only for $d=1$ and $d=5$,
describing the latter as illustrative; it refers the six remaining cases to
his dissertation.  We rely here on Somer's published statement of that
reference rather than on an independent inspection of the dissertation pages.  Since a complete proof of the small
range is useful to readers and to later computational checks, we give here
a single argument handling all $1\le d\le8$ at once.  It is independent of
Somer's case-by-case calculation and uses the structural lemmas of the
preceding section.

\begin{theorem}[Small-index classification]
\label{thm:small-index}
Let $n>1$ be odd, let $b\ge2$ with $\gcd(b,n)=1$, and suppose
\[
L_b(n)=\frac{n-1}{d},
\qquad 1\le d\le8.
\]
Then $n$ is prime except for
\[
\begin{array}{c|c}
 d&\text{composite exceptions}\\ \hline
4&9\\
5&6601\\
6&25\\
7&15,561\\
8&9,49.
\end{array}
\]
Every listed exception actually occurs; witnesses are
\[
\ord_9(8)=2,\quad
\ord_{6601}(11)=1320,\quad
\ord_{25}(7)=4,\quad
\ord_{15}(4)=2,
\]
\[
\ord_{561}(5)=80,\quad
\ord_9(10)=1,\quad
\ord_{49}(19)=6.
\]
\end{theorem}

\begin{proof}
Put
\[
\ell:=L_b(n)=\frac{n-1}{d},
\qquad
\phi(n)=j\ell,
\qquad
\lambda(n)=\kappa\ell.
\]
Because $\lambda(n)\mid\phi(n)$,
\[
\kappa\mid j.
\]
For composite $n$, $\phi(n)<n-1=d\ell$, so
\begin{equation}
1\le \kappa\le j\le d-1\le7.
\tag{S1}\label{eq:small-j}
\end{equation}
Write
\[
n=\prod_{i=1}^{s}p_i^{a_i},
\qquad 3\le p_1<\cdots<p_s,
\qquad
K=\prod_i p_i^{a_i-1}.
\]
By Lemma~\ref{lem:Kdivd}, $K\mid \kappa$.  By
Lemma~\ref{lem:H2adic},
\[
2^{s-1}\kappa\mid j.
\]
Together with \eqref{eq:small-j}, this gives $s\le3$.  We treat the three
possibilities separately.

\medskip
\noindent\textit{Case 1: one distinct prime factor.}
Let $n=p^a$, $a\ge2$.  Then $\lambda(n)=\phi(n)$, hence $\kappa=j$, and
$d\phi(n)=j(n-1)$ becomes
\begin{equation}
p^{a-1}\bigl((d-j)p-d\bigr)=-j.
\tag{S2}\label{eq:small-oneprime}
\end{equation}
Thus $p^{a-1}\mid j\le7$.  Hence $a=2$ and
$p\in\{3,5,7\}$.  Substitution in \eqref{eq:small-oneprime} gives only
\[
(d,j,n)=(4,3,9),\quad(6,5,25),\quad(8,6,9),\quad(8,7,49).
\]
These are exactly the prime-power exceptions displayed in the theorem.

\medskip
\noindent\textit{Case 2: two distinct prime factors.}
Write
\[
n=p^a q^c,
\qquad p<q,
\qquad
P=p^{a-1}q^{c-1}=K.
\]
The equation $d\phi(n)=j(n-1)$ is
\begin{equation}
P\bigl((d-j)pq-d(p+q)+d\bigr)=-j.
\tag{S3}\label{eq:small-two-master}
\end{equation}
Here $2\kappa\mid j$ and $P\mid \kappa$.  Since $j\le7$ and $P$ is odd, a
nonsquarefree solution would force $P=3$, hence $n=9q$ with $q>3$ prime.
Here $P=3$ forces $\kappa=3$, and $2\kappa\mid j\le7$ therefore gives
$j=6$.  Equation \eqref{eq:small-two-master} reduces to
\[
6d(q-1)=6(9q-1),
\qquad
q=\frac{d-1}{d-9}<0
\]
for $d\le8$.  Thus no nonsquarefree two-prime solution occurs.

Putting $P=1$ in \eqref{eq:small-two-master} and completing the product
gives the useful factorization
\begin{equation}
\bigl((d-j)p-d\bigr)
\bigl((d-j)q-d\bigr)=j^2.
\tag{S4}\label{eq:small-two-factor}
\end{equation}
For a valid solution the two factors in \eqref{eq:small-two-factor} are positive: solving the master equation for the larger prime gives a positive denominator.  The two-adic overlap budget should also be used before any further enumeration: since $2\kappa\mid j$, the integer $j$ is even.  This parity condition is what removes the formal odd-$j$ master-equation solutions (such as $21,33,65$ on other spectral lines).  Together with
$j\le7$, this leaves only $j\in\{2,4,6\}$.  Factoring
\eqref{eq:small-two-factor} over these three small squares gives exactly
\[
\begin{array}{c|c|c|c|c|c}
 d&j&(p,q)&n&\lambda(n)&\ell=(n-1)/d\\ \hline
5&4&(7,13)&91&12&18\\
7&4&(3,5)&15&4&2\\
8&6&(5,13)&65&12&8.
\end{array}
\]
The defining order can exist only if $\ell\mid\lambda(n)$.  This removes
the first and third rows, leaving only $n=15$ at $d=7$, and indeed
$\ord_{15}(4)=2$.

\medskip
\noindent\textit{Case 3: three distinct prime factors.}
Now
\[
4\kappa\mid j.
\]
Since $j\le7$, necessarily
\begin{equation}
\kappa=1,\qquad j=4.
\tag{S5}\label{eq:small-three-kappa}
\end{equation}
Because $K\mid \kappa$, we immediately get $K=1$: the integer is squarefree,
\[
n=pqr,
\qquad p<q<r.
\]
The Euler equation is
\begin{equation}
d(p-1)(q-1)(r-1)=j(pqr-1).
\tag{S6}\label{eq:small-three-master}
\end{equation}
For fixed $p$, put
\[
A=(d-j)p-d,
\qquad
B=d(p-1).
\]
Equation \eqref{eq:small-three-master} becomes
\[
Aqr-B(q+r)+B+j=0,
\]
and hence
\begin{equation}
(Aq-B)(Ar-B)
=j\bigl(d(p-1)^2+jp\bigr).
\tag{S7}\label{eq:small-three-factor}
\end{equation}
Furthermore, \eqref{eq:small-three-master} implies
\[
(d-j)pqr<d(pq+pr+qr)<3dqr,
\]
and a valid solution also requires $A=(d-j)p-d>0$, hence $p>d/(d-j)$.  Thus
\begin{equation}
\frac{d}{d-j}<p<\frac{3d}{d-j}.
\tag{S8}\label{eq:small-pbound}
\end{equation}
Thus only a handful of primes $p$ can occur, and for each one
\eqref{eq:small-three-factor} reduces the remaining search to divisor pairs
of one explicit integer.  Concretely, the bound gives $p=7,11,13$ for $d=5$,
$p=5,7$ for $d=6$, $p=3,5$ for $d=7$, and $p=3,5$ for $d=8$; checking the
positive divisor pairs in \eqref{eq:small-three-factor} leaves precisely the
three prime triples displayed below (all other pairs contain a composite
entry or violate the required ordering).  Since $j=4$ already, the complete
candidate list has only three rows:
\[
\begin{array}{c|c|c|c|c|c}
 d&j&(p,q,r)&n&\lambda(n)&\ell=(n-1)/d\\ \hline
5&4&(7,19,67)&8911&198&1782\\
5&4&(7,23,41)&6601&1320&1320\\
7&4&(3,11,17)&561&80&80.
\end{array}
\]
By \eqref{eq:small-three-kappa} one must have
$\lambda(n)=\ell$.  Exactly two rows survive:
\[
6601\quad(d=5),
\qquad
561\quad(d=7).
\]

Combining the three cases gives
\[
\begin{array}{c|c}
 d&\mathcal E(d)\\ \hline
1,2,3&\varnothing\\
4&\{9\}\\
5&\{6601\}\\
6&\{25\}\\
7&\{15,561\}\\
8&\{9,49\}.
\end{array}
\]
The displayed witness bases verify every surviving entry.
\end{proof}

\subsection{Comparison with Somer's proof}

The two proofs emphasize different structures.  Somer introduces
$\lambda'(N)=\operatorname{lcm}_{p\mid N}(p-1)$ and an auxiliary product
$\lambda''(N)$, proves the factor-count bound
$t\le\lfloor\log_2 d+1\rfloor$, bounds repeated prime powers, and then
handles each small value of $d$ separately.  In the 1989 paper only the
$d=5$ case is written out in detail: after several reductions it isolates
$7\cdot23\cdot41=6601$.

Our proof packages the same small-index scarcity into the two integers
\[
\kappa=\lambda(n)/L,
\qquad
H=\phi(n)/\lambda(n),
\]
with $j=\kappa H$.  The exact divisibilities $K\mid \kappa$ and
$2^{s-1}\kappa\mid j$ reduce the whole range at once to $s\le3$.  The cases
$s=2$ and $s=3$ are then controlled by the two explicit factorizations
\eqref{eq:small-two-factor} and \eqref{eq:small-three-factor}.  The resulting
proof is short enough to reproduce in full, while retaining a transparent
candidate list that can be checked by hand.

The classification itself should therefore be credited to Somer for his
nontrivial-base convention.  What is new here in the small range is the
unified proof, the inclusion of the two trivial-base spectral occurrences,
and the integration of the result into the effective spectral framework used
for larger $d$.

The later Carlip--Jacobson--Somer treatment \cite{CarlipJacobsonSomer1998}
also recovers the general Fermat factor-count bound $t<\log_2d+1$ and controls
prime powers.  The present argument should therefore not be read as claiming
novelty for those qualitative restrictions.  Its sharper bookkeeping retains
$\kappa$ and $H$ separately: $K\mid\kappa$ records the complete repeated-prime
part, while $2^{s-1}\kappa\mid j\le d-1$ couples repetition and overlap.  This
is what makes the uniform branch bounds and the later counting argument
possible.

\section{Certified computer-assisted classification through $d=64$}
\label{sec:certified-search}
\label{sec:computer-assisted}

For larger $d$ the number of possible distinct prime factors grows:
\[
s\le4\quad(9\le d\le16),
\qquad
s\le5\quad(17\le d\le32).
\]
A hand proof by separate Diophantine branches becomes increasingly
unattractive.  We therefore use a certified exact enumeration of the
finite factorization space.

Write
\[
n=P\prod_{i=1}^s p_i,
\qquad
P=\prod_{i=1}^s p_i^{a_i-1}.
\]
Since
\[
d\phi(n)=j(n-1),
\qquad
P\mid j,
\]
write
\[
j=cP.
\]
Then every candidate satisfies the master equation
\begin{equation}
d\prod_{i=1}^s(p_i-1)
=
c\left(P\prod_{i=1}^s p_i-1\right).
\label{eq:master-search}
\end{equation}
The ranges
\[
s\le1+\lfloor\log_2(d-1)\rfloor,
\qquad
1\le P\le d-1,
\qquad
1\le c\le\left\lfloor\frac{d-1}{P}\right\rfloor
\]
are finite.

Suppose the first $s-1$ primes have been chosen and put
\[
A=\prod_{i<s}(p_i-1),
\qquad
B=\prod_{i<s}p_i.
\]
Equation~\eqref{eq:master-search} is linear in the final prime:
\begin{equation}
p_s=
\frac{dA-c}{dA-cPB}.
\label{eq:lastprime}
\end{equation}
Thus the last prime is solved algebraically rather than scanned.

For counting purposes it is useful to stop one step earlier and solve the
last \emph{two} primes simultaneously.  This gives an exact divisor problem
rather than a second prime scan.

\begin{lemma}[Divisor-pair completion of the last two primes]
\label{lem:two-prime-completion}
Fix a master-equation branch $(d,s,j,P)$, put $c=j/P$, and suppose the first
$r=s-2$ primes have been chosen.  Write
\[
A=\prod_{i\le r}(p_i-1),\qquad
B=\prod_{i\le r}p_i,
\]
with the empty products equal to $1$.  If the two remaining primes are
$x<y$, define
\[
a=dA,\qquad b=cPB=jB,\qquad \Delta=a-b.
\]
Every valid completion has $\Delta>0$ and satisfies
\begin{equation}
\boxed{
(\Delta x-a)(\Delta y-a)=T,
\qquad
T:=a(b-c)+bc.
}
\label{eq:two-prime-divisor-pair}
\end{equation}
Consequently, for a fixed branch and prefix, there are at most
$\lfloor\tau(T)/2\rfloor$ ordered prime completions.
\end{lemma}

\begin{proof}
The master equation becomes
\[
a(x-1)(y-1)=bxy-c,
\]
or
\[
\Delta xy-a(x+y)+(a+c)=0.
\]
Solving for $y$ gives
\[
y=\frac{a(x-1)-c}{\Delta x-a}.
\]
The numerator is positive: $a=dA\ge d$, $x\ge3$, and $c\le j\le d-1$.
Hence the denominator is positive for every valid completion, and therefore
$\Delta>0$.  Multiplying the preceding bilinear equation by $\Delta$ and
completing the product gives
\[
(\Delta x-a)(\Delta y-a)
=a^2-\Delta(a+c)
=a(b-c)+bc.
\]
The two positive factors on the left form a divisor pair of $T$; because
$x<y$, at most half of the positive divisors can occur as the first factor.
Integrality, primality, the support condition for $P$, and the final
$\lambda$-divisibility checks can only reduce this number.
\end{proof}

The principal computational application of this lemma is the difficult
$(d,s,\kappa,H,P)=(49,5,3,16,1)$ branch in
Subsection~\ref{subsec:d64}: after the first three primes are exhausted,
exactly $49\,496$ prefixes remain and the last two primes are then completed
entirely by \eqref{eq:two-prime-divisor-pair}.

At an intermediate prefix, put
\[
A_r=\prod_{i\le r}(p_i-1),
\qquad
B_r=\prod_{i\le r}p_i,
\qquad
R=\frac{dA_r}{cPB_r}.
\]
If $R\le1$, the branch is impossible.  Suppose that $R>1$, that $t$
primes remain, and that the next prime is $x$.  Dividing
\eqref{eq:master-search} by $cPB_r$ gives the necessary inequality
\[
R\prod_{\text{remaining }q}\left(1-\frac1q\right)<1.
\]
Every remaining prime satisfies $q\ge x$, so
\[
\prod_{\text{remaining }q}\left(1-\frac1q\right)
\ge
\left(\frac{x-1}{x}\right)^t.
\]
Hence every surviving branch must satisfy
\begin{equation}
\boxed{
x^t\,cPB_r
>
(x-1)^t\,dA_r.
}
\label{eq:exactbound}
\end{equation}
The predicate is monotone in $x$, so the largest admissible $x$ is obtained
by binary search using exact integer comparisons.  No floating-point
approximation is allowed to decide whether a branch is discarded.

Every completed candidate is checked for
\[
d\mid n-1,
\qquad
\frac{n-1}{d}\mid\lambda(n).
\]
The first of these checks is essential: an earlier exploratory script
used integer division without the divisibility test and thereby created
several false positives in preliminary tables.

The overlap pruning used below can be updated recursively.  Put
\[
c_i=p_i^{a_i-1}(p_i-1),
\qquad
\Lambda_k=\operatorname{lcm}(c_1,\ldots,c_k),
\qquad
H_k=\frac{\prod_{i=1}^k c_i}{\Lambda_k}.
\]
Then
\begin{equation}
H_k=H_{k-1}\gcd(\Lambda_{k-1},c_k).
\label{eq:overlap}
\end{equation}
Thus each new prime consumes part of the finite overlap budget
$H_s=\phi(n)/\lambda(n)$.  In the minimal-overlap case $H_s=2^{s-1}$ this
forces $\gcd(c_i,c_j)=2$ for $i\ne j$; for squarefree $n$,
$\gcd(p_i-1,p_j-1)=2$.

\begin{proposition}[Finite and complete search specification]
\label{prop:search-spec}
For fixed spectral index $d$, every composite solution occurs in the following
finite search.
\begin{enumerate}[label=(\arabic*)]
\item Handle $s=1$ separately: enumerate odd prime powers $p^a$, $a\ge2$,
with $p^{a-1}\le d-1$, and retain exactly those satisfying
$d\mid p^a-1$ and $(p^a-1)/d\mid\lambda(p^a)$.
\item For $s\ge2$, enumerate
\[
2\le s\le S(d),\quad
1\le\kappa\le\frac{d-1}{2^{s-1}},\quad
H\in\{2^{s-1},2\cdot2^{s-1},\ldots\},\quad
\kappa H\le d-1.
\]
For each branch set $j=\kappa H$, enumerate $P\mid\kappa$, and put $c=j/P$.
\item Factor $P$.  Every prime dividing $P=n/\rad(n)$ is a required support
prime and $v_q(P)=a_q-1$.  Recursion uses strictly increasing distinct primes.
\item Apply only necessary prunes: required-prime ordering and remaining-slot
checks; $H_k\mid H$ and the remaining overlap budget from
\eqref{eq:overlap}; $R>1$; and the exact ratio bound
\eqref{eq:exactbound}.
\item Solve the final prime by \eqref{eq:lastprime}.  Reject it if it is not
an integer prime strictly larger than the preceding prime, if required support
is missing, or if the exact terminal overlap is not $H$.  Finally recompute
$n,\phi(n),\lambda(n),L$ and all defining divisibilities.
\end{enumerate}
For $s\ge2$ the number of parameter branches before prefix recursion is
\[
B(d)=\sum_{s=2}^{S(d)}
\sum_{\kappa\le(d-1)/2^{s-1}}
\left\lfloor\frac{d-1}{2^{s-1}\kappa}\right\rfloor\tau(\kappa).
\]
\end{proposition}

\begin{proof}
The prime-power restriction in step (1) follows from
$p^{a-1}\mid\kappa\le d-1$.  For $s\ge2$, Corollary~\ref{cor:sbound},
Lemma~\ref{lem:H2adic}, and Lemma~\ref{lem:Kdivd} place every genuine
solution in one enumerated $(s,\kappa,H,P)$ branch.  The identity
$P=\prod p_i^{a_i-1}$ proves the support rule.

No recursive prune can remove a prefix of a genuine solution.  If $R\le1$,
appending another prime multiplies $R$ by $(q-1)/q<1$, whereas positivity of
the denominator in \eqref{eq:lastprime} requires $R>1$ immediately before
the last prime.  The ratio bound is the necessary inequality derived above.
For the overlap tests, $H_k$ must divide the target $H$, and after the first
local factor every future multiplier
$\gcd(\Lambda_{i-1},c_i)$ is at least $2$ because both arguments are even;
this proves the remaining-budget test.  The final prime in
\eqref{eq:lastprime} is the unique possible algebraic completion.  Requiring
it to exceed the prefix loses nothing because every factorization is enumerated
in increasing order.  Finally $H_s=H$ is an identity for the fixed branch, so
the exact terminal check is necessary.  The final arithmetic tests are exactly
the conditions of Proposition~\ref{prop:composite-spectrum}.
\end{proof}

The implementation follows the proposition literally.  In pseudocode:
\begin{verbatim}
prime_power_solutions(d)
for s, kappa, H, P in all admissible branches:
    required := prime_support(P)
    recurse(increasing_prefix):
        reject if required-prime order/slots are impossible
        reject unless H_prefix divides H and overlap budget can finish
        reject if R <= 1
        bound next prime by the exact ratio inequality
        recurse over every prime within that exact bound
        at length s-1 solve the last prime algebraically
        reject unless last prime > prefix and required support is complete
        reject unless H_full == H
        recompute phi, lambda, L and all defining divisibilities
\end{verbatim}
Thus the written specification includes every mathematical prune used by the
reference search.

\begin{theorem}[Computer-assisted classification for $9\le d\le32$]
\label{thm:d32}
For every $9\le d\le32$, $\Eset(d)$ is exactly the corresponding row of
Table~\ref{tab:exceptions}.
\end{theorem}

\begin{proof}
Proposition~\ref{prop:search-spec} proves completeness of the finite search.
The reference implementation exhausts every branch with exact integer
arithmetic and an arithmetic verifier independently recomputes every reported
solution.  No integer-size cutoff is used in the Diophantine search.
\end{proof}

\subsection{The dyadic transition at $d=33$}
\label{subsec:d33}

At $d=33$, six distinct prime factors become structurally possible.  The two
largest $\kappa=1$ branches can be removed analytically.

\begin{lemma}[Odd-index $2$-adic parity]
\label{lem:odd-index-parity}
Let $n=\prod_{i=1}^s p_i$ be squarefree and
$\lambda(n)=(n-1)/d$ with $d$ odd.  Put
\[
e_i=v_2(p_i-1),\qquad E=\max_i e_i,
\qquad a=\#\{i:e_i=1\}.
\]
Then $v_2(n-1)=E$.  If $E=1$, then $a=s$ is odd; if $E\ge2$, then $a$ is
even.
\end{lemma}

\begin{proof}
Because $d$ is odd,
$v_2(n-1)=v_2(\lambda(n))=E$.  Also $e_i=1$ exactly when
$p_i\equiv3\pmod4$, so $n\equiv(-1)^a\pmod4$.  The two cases follow.
\end{proof}

\begin{theorem}[Exact spectral line at $d=33$]
\label{thm:d33}
\[
\mathcal E(33)=\{1047619,2210671,3494701,15197491,25158871,93755971\}.
\]
In particular $N_{\max}(33)=93\,755\,971$.
\end{theorem}

\begin{proof}
The branch formula gives $B(33)=166=31+135$, separating $31$ branches with
$\kappa=1$ from $135$ with $\kappa>1$.  In the $s=6$, $\kappa=1$ branch,
$H=32$ and $\sum e_i-E=5$, so the five nonmaximal exponents are all $1$;
Lemma~\ref{lem:odd-index-parity} gives a contradiction.  In the $s=5$,
$\kappa=1$, $H=32$ branch, four positive nonmaximal exponents sum to $5$,
so their multiset is $\{2,1,1,1\}$; here $E\ge2$, leaving exactly three
primes congruent to $3\pmod4$, again contradicting the lemma.  The program
\texttt{certify\_d33\_reduced.jl} exhausts the remaining $29$ $\kappa=1$
branches and finds none.  Independently,
\texttt{certify\_d33\_noncar.jl} exhausts all $135$ branches with
$\kappa>1$ and returns exactly the six displayed integers.  Step (1) of
Proposition~\ref{prop:search-spec} gives no additional prime-power solution
at $d=33$.  This exhausts all strata.
\end{proof}

The six surviving points are all squarefree and have exactly three prime
factors.  Their branch invariants are:
\begin{center}
\small
\begin{tabular}{r l r r r r r}
\toprule
$n$ & factorization & $\kappa$ & $H$ & $j$ & $\lambda(n)$ & $L=(n-1)/33$\\
\midrule
$1\,047\,619$  & $79\cdot89\cdot149$    & $4$ & $8$ & $32$ & $126\,984$    & $31\,746$\\
$2\,210\,671$  & $59\cdot89\cdot421$    & $4$ & $8$ & $32$ & $267\,960$    & $66\,990$\\
$3\,494\,701$  & $7\cdot101\cdot4943$   & $7$ & $4$ & $28$ & $741\,300$    & $105\,900$\\
$15\,197\,491$ & $37\cdot431\cdot953$  & $4$ & $8$ & $32$ & $1\,842\,120$ & $460\,530$\\
$25\,158\,871$ & $41\cdot173\cdot3547$ & $4$ & $8$ & $32$ & $3\,049\,560$ & $762\,390$\\
$93\,755\,971$ & $41\cdot167\cdot13693$& $4$ & $8$ & $32$ & $11\,364\,360$& $2\,841\,090$\\
\bottomrule
\end{tabular}
\end{center}
Thus this dyadic transition permits six prime factors structurally, but
none of the actual exceptions uses more than three; in particular the entire
$\kappa=1$ (Carmichael-index) slice is empty at $d=33$.

The branch $\kappa=1$ is exactly the Carmichael-index problem
$n-1=d\lambda(n)$ studied in the small-index Carmichael literature; compare
Pinch~\cite{PinchSmallIndex}.  For external comparison, Pinch's exhaustive
Carmichael-number tables were extended from the earlier $10^{15}$ computation
to a complete enumeration through $10^{21}$ \cite{Pinch2007}.  The present exclusion of the
$\kappa=1$ line at $d=33$ does not rely on absence from Pinch's table:
the two largest branches are eliminated analytically by
Lemma~\ref{lem:odd-index-parity}, and the remaining branches are exhausted by
the certified exact search below.  Pinch's table is therefore an independent
large-range check rather than an input to the proof.


\subsection{Extension from $d=34$ through $d=64$}
\label{subsec:d64}

The interval $34\le d\le64$ has a useful structural split.
If $\kappa>1$, then
\[
2^{s-1}\kappa\le d-1\le63
\]
forces $s\le5$.  If $\kappa=1$, then $K=1$ and
\[
n-1=d\lambda(n),
\]
so the number is precisely a Carmichael number of Carmichael index $d$.
We independently implement the finite prescribed-index recursion for this
$\kappa=1$ stratum.  At each prefix the exact ratio bound and the necessary
overlap condition restrict the next prime; with two primes left we use
Lemma~\ref{lem:two-prime-completion}, factoring the resulting integer $T$
and reconstructing every admissible divisor pair.  The two largest empty
cases, $d=34$ and $d=35$, are recorded as disjoint exact prefix partitions in
the supplement; the union of each partition is the full branch, not a size
cutoff.  All factor outputs are reconstructed and their prime factors are
independently certified.  The completed computation agrees exactly with
Pinch's published prescribed-index list through $64$ \cite{PinchSmallIndex},
so Pinch is an external cross-check rather than a proof dependency.  The
nonempty new $\kappa=1$ rows occur at
\[
37,39,43,44,45,47,48,49,50,52,53,54,55,60,61.
\]

For $\kappa>1$ we used a second exact implementation of
Proposition~\ref{prop:search-spec}, written in C++ with arbitrary-precision
integers for the branch algebra.  Every exact recursive prime bound encountered
in the completed $34\le d\le64$ runs lies below $2^{64}$; the program aborts
rather than truncating a branch if this condition fails.  Primality is tested
with the first twelve prime Miller--Rabin bases $2,3,\ldots,37$.
Sorenson--Webster proved that the smallest composite passing all twelve is
$318665857834031151167461>2^{64}$ \cite{SorensonWebster2017}; hence every
C++ primality decision in the recorded 64-bit domain is deterministic.  The
implementation exhausts all such branches.  The only noticeably expensive
branch is
\[
(d,s,\kappa,H,P)=(49,5,3,16,1).
\]
For it we stop after three primes and apply Lemma~\ref{lem:two-prime-completion}.
There are exactly $49\,496$ surviving three-prime prefixes; exact divisor-pair
completion yields no solution.  The source and captured log are supplied.
Prime powers are enumerated separately by step (1) of
Proposition~\ref{prop:search-spec}.

\begin{theorem}[Certified composite spectrum through $d=64$]
\label{thm:d64}
For every $1\le d\le64$, the inclusive composite spectral line $\Eset(d)$ is
exactly the corresponding row of Table~\ref{tab:exceptions}.  Across the
first $64$ lines there are $351$ incidences representing $234$ distinct odd
composites.  The largest row size is
\[
\max_{d\le64} e(d)=13,
\]
attained at $d=53$ and $d=62$.
\end{theorem}

\begin{proof}
The range $d\le8$ is Theorem~\ref{thm:small-index}, the range $9\le d\le32$
is Theorem~\ref{thm:d32}, and $d=33$ is Theorem~\ref{thm:d33}.  For
$34\le d\le64$, the prime-power stratum is exhausted by step (1) of
Proposition~\ref{prop:search-spec}; the $\kappa=1$ stratum is exhausted by
the independent prescribed-index certificate just described (with Pinch used
only as a cross-check); and every $\kappa>1$ branch is exhausted by
the exact search just described, with the single $d=49$ branch completed by
the divisor-pair identity.  The accompanying verifier recomputes
$\lambda(n)$ and both defining divisibilities for all $351$ incidences and
returns $234$ distinct integers.  Hence no stratum is omitted.
\end{proof}


\subsection{Reproducibility of the computer-assisted part}
\label{subsec:reproducibility}

The computational claims are supplied with machine-readable supplementary
material rather than being treated as an opaque numerical experiment.  The
reference implementation uses exact integer arithmetic for every pruning
decision and performs the final checks
\[
d\mid n-1,\qquad (n-1)/d\mid\lambda(n)
\]
for every reported candidate.  Failed tasks are required to abort a run.

The original frozen $d\le33$ subarchive has SHA--256
\[
\text{\scriptsize\ttfamily bc7b993c4a58678de2aa2303b4e011c8d5a319c24ad805ff00515452f47881b2}.
\]
It is retained unchanged inside the extended supplement.  The present arXiv
supplement additionally contains the exact $34\le d\le64$ C++ search,
the specialized $d=49$ divisor-pair certificate, the independent $\kappa=1$
prescribed-index certificate and its exact split logs, the Pinch cross-check,
the combined machine-readable $1\le d\le64$ table, its arithmetic verifier,
the exact prime-mass calculation through $64$, and the arXiv-only low-$\omega$
and OEIS data.  A machine-readable SHA manifest identifies every file.  The
legacy subarchive contains the following proof-critical files:
\begin{center}
\scriptsize
\begin{tabularx}{\textwidth}{lX}
\toprule
file & SHA--256\\
\midrule
\texttt{overlap\_U\_search.jl} &
\texttt{df97f12dbf322229\allowbreak 91131f36482fece2\allowbreak ab752643638691c3\allowbreak 6a1fbf37c45ae426}\\
\texttt{verify\_overlap\_U\_output.jl} &
\texttt{670adbd590865eca\allowbreak 6eb4b9f1154ccdee\allowbreak 87e9518a8c2180ec\allowbreak 051cfe9ee113bf17}\\
\texttt{certify\_d33\_reduced.jl} &
\texttt{b023717245e965b1\allowbreak 6559edc4a0e6d211\allowbreak c34bd82274d3409f\allowbreak a510f1dc93741da7}\\
\texttt{certify\_d33\_noncar.jl} &
\texttt{57a26bb49947a99a\allowbreak 6a712ac2e0c25bf6\allowbreak 5b033bff3bb0ccb5\allowbreak 8167647df03d0f08}\\
\texttt{prime\_mass\_d33\_check.jl} &
\texttt{dd3fd20d28a51606\allowbreak 15da798fa40d0a7d\allowbreak 642ac6ae9191b5ec\allowbreak 41ad0b94d8eda212}\\
\texttt{branch\_parameter\_counts\_d9\_d33.tsv} &
\texttt{60ead43be1cea3c1\allowbreak af6ab85d8590dd3c\allowbreak f5330864f5b74a8c\allowbreak 763585ee8440508d}\\
\bottomrule
\end{tabularx}
\end{center}

The frozen reference computation was executed under Julia~1.12.6 on
August~23, 2026.  The $d\le8$ regression self-test passed.  The general exact
search then exhausted every parameter branch for $9\le d\le32$ and reproduced
the corresponding rows of Table~\ref{tab:exceptions}; the separate arithmetic verifier
checked all $96$ resulting incidences.  It uses no external packages,
deterministically certifies every reported prime factor below $2^{64}$, and
recomputes $n$, $\phi(n)$, $\lambda(n)$, $L$, $\kappa$, $H$, the repeated-prime
factor $P$, and all defining divisibilities from the reported factorizations.

The dyadic transition $d=33$ is certified by the split proof described above,
rather than by an unspecialized rerun of the general search.  After the two
$\kappa=1$ branches eliminated analytically by Lemma~\ref{lem:odd-index-parity}
are removed, \texttt{certify\_d33\_reduced.jl} exhausts all $29$ remaining
$\kappa=1$ parameter branches and finds none.  Independently,
\texttt{certify\_d33\_noncar.jl} exhausts all $135$ branches with $\kappa>1$
and returns exactly the six integers in Theorem~\ref{thm:d33}; it then checks
the factorizations and the $(\kappa,H,j,\lambda,L)$ data displayed above.  The
reference runs used the command-line option \texttt{--resume=false} for a clean non-resuming run.  The supplementary
branch-count table records the exact parameter-branch counts for every
$9\le d\le33$.

For $34\le d\le64$, the supplement contains the exact C++ branch search,
the independent $\kappa=1$ prescribed-index certificate, its exact split logs
for the two dominant empty indices $34$ and $35$, the Pinch cross-check, the
specialized $d=49$ last-two-prime certificate, and an independent Python
arithmetic verifier.  The verifier recomputes $\lambda(n)$ and the defining
divisibilities for all $351$ incidences and returns $234$ distinct integers.
Every recursive prime bound encountered in the completed $\kappa>1$ runs lies
below $2^{64}$; this is not an upper bound on $n$, and the search aborts rather
than truncating a branch if a required prime bound exceeds the deterministic
primality range.  The recorded runs also print the maximum exact recursive
prime bound and the maximum primality-test input, together with an explicit
\texttt{overflow=false} status.

The density calculation used in Theorem~\ref{thm:prime-mass-64} is certified
separately by \path{prime_mass_d64_check_exact.py}.  It checks all
$2^{18}-1=262143$ squarefree kernels supported on primes through $61$ using
exact integer/rational arithmetic, proves the $q\ge67$ tail exclusion by
$G_{64}-\mathfrak M_2(64)>G_{64}/4421$, and reproduces both quoted masses and
the unique minimizing squarefree kernel $2$.  Captured output is included in
the supplement.

The legacy Julia source uses only standard libraries.  The new C++ search
tests primality with the first twelve prime bases $2,3,\ldots,37$ and aborts
instead of silently changing to a probable-prime decision outside its 64-bit
search domain.  Sorenson--Webster's value
$\psi_{12}=318665857834031151167461$ proves every such C++ decision
\cite{SorensonWebster2017}.  The specialized $d=49$ divisor-pair certificate
checks complete factorizations of the integers $T$ factor by factor with the
same certificate.  The independent $\kappa=1$ program uses $\psi_{12}$ and,
for a few larger auxiliary factors, Sorenson--Webster's
$\psi_{13}=3317044064679887385961981$; above that it falls back to a recursive
Pocklington certificate after an exact factorization of $n-1$.  Thus no
probable-prime assumption occurs in either certificate.  The separate verifier
is an arithmetic verifier of reported outputs; it does not replace the
completeness proof in Proposition~\ref{prop:search-spec}.  This explicit separation
between theorem, search engine, specialized $d=33$ certification, and
independent arithmetic verification is especially important because an earlier
exploratory program used integer division before checking $d\mid n-1$ and
consequently produced several false positives.  Those exploratory outputs are
not used in any theorem of the present paper.

The exact supplement accompanying this revision is frozen and identified by
its SHA--256 manifest.  The unchanged archive is intended for deposition in
Zenodo or an institutional repository; the resulting persistent identifier
will replace the placeholder \texttt{[SUPPLEMENT-DOI]} in the archival version.
The deposited archive itself will not subsequently be modified.

\begin{table}[htbp]
\centering
\small
\caption{Certificate summary for the new $34\le d\le64$ computation.
``Max bound'' is the largest exact recursive prime bound recorded by the
component; every recorded overflow/abort flag is false.}
\label{tab:certificate-summary}
\scriptsize
\resizebox{\textwidth}{!}{%
\begin{tabular}{llllrl}
\toprule
component & program / theorem & range & work & max bound & status\\
\midrule
$\kappa=1$ prescribed-index & ratio recursion + Lemma~\ref{lem:two-prime-completion} & $d=34$ & $15$ branches; $344284$ prefixes & $2140591$ & pass\\
$\kappa=1$ prescribed-index & same & $d=35$ & $15$ branches; $381524$ prefixes & $7296343$ & pass\\
$\kappa=1$ prescribed-index & same & $36\le d\le48$ & $227$ branches; $202725$ prefixes & $1113271$ & Pinch match\\
$\kappa=1$ prescribed-index & same & $49\le d\le64$ & $384$ branches; $194688$ prefixes & $305758$ & Pinch match\\
$\kappa>1$ C++ & Proposition~\ref{prop:search-spec}; 12-base MR & $34\le d\le48$ & $2645$ branches & $4864229$ & overflow=false\\
$\kappa>1$ C++ & same & $50\le d\le54$ & $1360$ branches & $22551385$ & overflow=false\\
$\kappa>1$ C++ & same & $55\le d\le59$ & $1523$ branches & $1368683$ & overflow=false\\
$\kappa>1$ C++ & same & $60\le d\le64$ & $1725$ branches & $31170$ & overflow=false\\
$d=49$ special branch & Lemma~\ref{lem:two-prime-completion} & $(49,5,3,16,1)$ & $49496$ three-prime prefixes & --- & pass\\
\bottomrule
\end{tabular}%
}
\end{table}

\begin{longtable}{r r r r >{\raggedright\arraybackslash}p{0.58\textwidth}}
\caption{Certified inclusive composite spectrum through $d=64$.  Here
$e(d)=|\Eset(d)|$, $g(d)$ counts first (genuine) occurrences, $C(d)$ is the
cumulative number of distinct composites, and $\ast$ marks inheritance from a
smaller spectral index.}\label{tab:exceptions}\\
\toprule
$d$ & $e(d)$ & $g(d)$ & $C(d)$ & $\Eset(d)$\\
\midrule
\endfirsthead
\toprule
$d$ & $e(d)$ & $g(d)$ & $C(d)$ & $\Eset(d)$\\
\midrule
\endhead
1&0&0&0&---\\
2&0&0&0&---\\
3&0&0&0&---\\
4&1&1&1&$9$\\
5&1&1&2&$6601$\\
6&1&1&3&$25$\\
7&2&2&5&$15,\allowbreak\ 561$\\
8&2&1&6&$9^{\ast},\allowbreak\ 49$\\
9&3&3&9&$30889,\allowbreak\ 88561,\allowbreak\ 534061$\\
10&2&1&10&$21,\allowbreak\ 6601^{\ast}$\\
11&1&1&11&$45$\\
12&2&1&12&$25^{\ast},\allowbreak\ 121$\\
13&3&3&15&$27,\allowbreak\ 126673,\allowbreak\ 1333333$\\
14&3&1&16&$15^{\ast},\allowbreak\ 169,\allowbreak\ 561^{\ast}$\\
15&2&1&17&$91,\allowbreak\ 6601^{\ast}$\\
16&3&2&19&$33,\allowbreak\ 49^{\ast},\allowbreak\ 65$\\
17&2&2&21&$35,\allowbreak\ 139231$\\
18&7&4&25&$289,\allowbreak\ 30889^{\ast},\allowbreak\ 88561^{\ast},\allowbreak\ 534061^{\ast},\allowbreak\ 55462177,\allowbreak\ 8885251441,\allowbreak\ 42018333841$\\
19&4&4&29&$39,\allowbreak\ 153,\allowbreak\ 4371,\allowbreak\ 7107$\\
20&4&2&31&$21^{\ast},\allowbreak\ 361,\allowbreak\ 3201,\allowbreak\ 6601^{\ast}$\\
21&7&7&38&$85,\allowbreak\ 8695,\allowbreak\ 10585,\allowbreak\ 26335,\allowbreak\ 9717247,\allowbreak\ 16666651,\allowbreak\ 37769887$\\
22&3&2&40&$45^{\ast},\allowbreak\ 133,\allowbreak\ 2465$\\
23&5&5&45&$231,\allowbreak\ 645,\allowbreak\ 1105,\allowbreak\ 24013,\allowbreak\ 34133$\\
24&4&1&46&$25^{\ast},\allowbreak\ 49^{\ast},\allowbreak\ 121^{\ast},\allowbreak\ 529$\\
25&4&3&49&$51,\allowbreak\ 6601^{\ast},\allowbreak\ 15051,\allowbreak\ 11921001$\\
26&4&1&50&$27^{\ast},\allowbreak\ 105,\allowbreak\ 126673^{\ast},\allowbreak\ 1333333^{\ast}$\\
27&6&3&53&$55,\allowbreak\ 325,\allowbreak\ 3565,\allowbreak\ 30889^{\ast},\allowbreak\ 88561^{\ast},\allowbreak\ 534061^{\ast}$\\
28&3&1&54&$57,\allowbreak\ 169^{\ast},\allowbreak\ 561^{\ast}$\\
29&6&6&60&$117,\allowbreak\ 175,\allowbreak\ 1045,\allowbreak\ 1447681,\allowbreak\ 2121061,\allowbreak\ 164424781$\\
30&4&2&62&$91^{\ast},\allowbreak\ 841,\allowbreak\ 6601^{\ast},\allowbreak\ 158401$\\
31&8&8&70&$63,\allowbreak\ 125,\allowbreak\ 435,\allowbreak\ 62745,\allowbreak\ 1154131,\allowbreak\ 2609581,\allowbreak\ 3936691,\allowbreak\ 5624641$\\
32&6&4&74&$33^{\ast},\allowbreak\ 65^{\ast},\allowbreak\ 385,\allowbreak\ 961,\allowbreak\ 12673,\allowbreak\ 12801$\\
33&6&6&80&$1047619,\allowbreak\ 2210671,\allowbreak\ 3494701,\allowbreak\ 15197491,\allowbreak\ 25158871,\allowbreak\ 93755971$\\
34&4&2&82&$35^{\ast},\allowbreak\ 69,\allowbreak\ 341,\allowbreak\ 139231^{\ast}$\\
35&4&3&85&$561^{\ast},\allowbreak\ 91001,\allowbreak\ 610051,\allowbreak\ 683761$\\
36&9&2&87&$145,\allowbreak\ 217,\allowbreak\ 289^{\ast},\allowbreak\ 30889^{\ast},\allowbreak\ 88561^{\ast},\allowbreak\ 534061^{\ast},\allowbreak\ 55462177^{\ast},\allowbreak\ 8885251441^{\ast},\allowbreak\ 42018333841^{\ast}$\\
37&11&11&98&$75,\allowbreak\ 17169,\allowbreak\ 33153,\allowbreak\ 45141,\allowbreak\ 62605,\allowbreak\ 2415953,\allowbreak\ 2929291,\allowbreak\ 11972017,\allowbreak\ 16894571,\allowbreak\ 67902031,\allowbreak\ 191834567$\\
38&9&5&103&$39^{\ast},\allowbreak\ 77,\allowbreak\ 153^{\ast},\allowbreak\ 1065,\allowbreak\ 1369,\allowbreak\ 4371^{\ast},\allowbreak\ 7107^{\ast},\allowbreak\ 13833,\allowbreak\ 1170401$\\
39&5&3&106&$703,\allowbreak\ 15211,\allowbreak\ 126673^{\ast},\allowbreak\ 334153,\allowbreak\ 1333333^{\ast}$\\
40&7&4&110&$81,\allowbreak\ 361^{\ast},\allowbreak\ 481,\allowbreak\ 3201^{\ast},\allowbreak\ 6601^{\ast},\allowbreak\ 460801,\allowbreak\ 1372801$\\
41&8&8&118&$165,\allowbreak\ 247,\allowbreak\ 2871,\allowbreak\ 8365,\allowbreak\ 73555,\allowbreak\ 184255,\allowbreak\ 9345541,\allowbreak\ 128950249$\\
42&9&1&119&$85^{\ast},\allowbreak\ 169^{\ast},\allowbreak\ 1681,\allowbreak\ 8695^{\ast},\allowbreak\ 10585^{\ast},\allowbreak\ 26335^{\ast},\allowbreak\ 9717247^{\ast},\allowbreak\ 16666651^{\ast},\allowbreak\ 37769887^{\ast}$\\
43&7&7&126&$87,\allowbreak\ 259,\allowbreak\ 861,\allowbreak\ 22791,\allowbreak\ 52633,\allowbreak\ 69231,\allowbreak\ 82733$\\
44&5&2&128&$45^{\ast},\allowbreak\ 133^{\ast},\allowbreak\ 1849,\allowbreak\ 2465^{\ast},\allowbreak\ 15841$\\
45&7&4&132&$91^{\ast},\allowbreak\ 451,\allowbreak\ 8911,\allowbreak\ 88561^{\ast},\allowbreak\ 534061^{\ast},\allowbreak\ 3763801,\allowbreak\ 14245741$\\
46&8&3&135&$93,\allowbreak\ 185,\allowbreak\ 231^{\ast},\allowbreak\ 369,\allowbreak\ 645^{\ast},\allowbreak\ 1105^{\ast},\allowbreak\ 24013^{\ast},\allowbreak\ 34133^{\ast}$\\
47&9&9&144&$95,\allowbreak\ 1035,\allowbreak\ 2821,\allowbreak\ 4795,\allowbreak\ 204263,\allowbreak\ 223721,\allowbreak\ 308039,\allowbreak\ 559489,\allowbreak\ 38649041$\\
48&5&3&147&$49^{\ast},\allowbreak\ 529^{\ast},\allowbreak\ 1729,\allowbreak\ 2209,\allowbreak\ 52417$\\
49&12&12&159&$99,\allowbreak\ 87676681,\allowbreak\ 133861729,\allowbreak\ 161005573,\allowbreak\ 258844951,\allowbreak\ 1411040359,\allowbreak\ 3462426241,\allowbreak\ 7950140448251,\allowbreak\ 131097700645913,\allowbreak\ 798178197617711,\allowbreak\ 989473937059427,\allowbreak\ 1208361237478669$\\
50&9&5&164&$51^{\ast},\allowbreak\ 301,\allowbreak\ 1001,\allowbreak\ 6601^{\ast},\allowbreak\ 15051^{\ast},\allowbreak\ 3679201,\allowbreak\ 11921001^{\ast},\allowbreak\ 4199932801,\allowbreak\ 1353385637801$\\
51&3&2&166&$205,\allowbreak\ 32131,\allowbreak\ 139231^{\ast}$\\
52&8&5&171&$105^{\ast},\allowbreak\ 833,\allowbreak\ 9361,\allowbreak\ 126673^{\ast},\allowbreak\ 741937,\allowbreak\ 1333333^{\ast},\allowbreak\ 206955841,\allowbreak\ 560630848961$\\
53&13&13&184&$425,\allowbreak\ 637,\allowbreak\ 3605,\allowbreak\ 8481,\allowbreak\ 55651,\allowbreak\ 3547291,\allowbreak\ 3984541,\allowbreak\ 43523707,\allowbreak\ 52519291,\allowbreak\ 260410837,\allowbreak\ 1271325841,\allowbreak\ 152674725071,\allowbreak\ 3916268927681$\\
54&8&2&186&$55^{\ast},\allowbreak\ 325^{\ast},\allowbreak\ 2809,\allowbreak\ 3565^{\ast},\allowbreak\ 30889^{\ast},\allowbreak\ 88561^{\ast},\allowbreak\ 534061^{\ast},\allowbreak\ 4169867689$\\
55&10&9&195&$111,\allowbreak\ 221,\allowbreak\ 1431,\allowbreak\ 6601^{\ast},\allowbreak\ 15181,\allowbreak\ 100101,\allowbreak\ 689701,\allowbreak\ 1269621,\allowbreak\ 271794601,\allowbreak\ 568898551$\\
56&6&3&198&$57^{\ast},\allowbreak\ 169^{\ast},\allowbreak\ 225,\allowbreak\ 561^{\ast},\allowbreak\ 25761,\allowbreak\ 376993$\\
57&6&6&204&$115,\allowbreak\ 343,\allowbreak\ 2737,\allowbreak\ 95761,\allowbreak\ 37883341,\allowbreak\ 713798461$\\
58&9&3&207&$117^{\ast},\allowbreak\ 175^{\ast},\allowbreak\ 1045^{\ast},\allowbreak\ 2553,\allowbreak\ 80041,\allowbreak\ 1447681^{\ast},\allowbreak\ 2121061^{\ast},\allowbreak\ 8372881,\allowbreak\ 164424781^{\ast}$\\
59&4&4&211&$119,\allowbreak\ 1653,\allowbreak\ 1771,\allowbreak\ 42353741$\\
60&8&3&214&$121^{\ast},\allowbreak\ 361^{\ast},\allowbreak\ 841^{\ast},\allowbreak\ 3481,\allowbreak\ 6601^{\ast},\allowbreak\ 25201,\allowbreak\ 158401^{\ast},\allowbreak\ 6840001$\\
61&12&12&226&$123,\allowbreak\ 245,\allowbreak\ 1221,\allowbreak\ 44409,\allowbreak\ 50997,\allowbreak\ 74665,\allowbreak\ 631351,\allowbreak\ 964411,\allowbreak\ 1102759,\allowbreak\ 940724555,\allowbreak\ 1962804565,\allowbreak\ 2417562005$\\
62&13&5&231&$63^{\ast},\allowbreak\ 125^{\ast},\allowbreak\ 435^{\ast},\allowbreak\ 3721,\allowbreak\ 19345,\allowbreak\ 62745^{\ast},\allowbreak\ 264121,\allowbreak\ 828073,\allowbreak\ 1154131^{\ast},\allowbreak\ 2609581^{\ast},\allowbreak\ 3936691^{\ast},\allowbreak\ 5624641^{\ast},\allowbreak\ 7697921$\\
63&7&1&232&$1891,\allowbreak\ 8695^{\ast},\allowbreak\ 10585^{\ast},\allowbreak\ 26335^{\ast},\allowbreak\ 9717247^{\ast},\allowbreak\ 16666651^{\ast},\allowbreak\ 37769887^{\ast}$\\
64&7&2&234&$65^{\ast},\allowbreak\ 129,\allowbreak\ 385^{\ast},\allowbreak\ 961^{\ast},\allowbreak\ 1281,\allowbreak\ 12673^{\ast},\allowbreak\ 12801^{\ast}$\\
\bottomrule\end{longtable}

\begin{theorem}[Exact primality criterion through spectral index $64$]
\label{thm:primality64}
Let $n>1$ be odd, let $\gcd(b,n)=1$, and suppose the exact order $L=\ord_n(b)$
satisfies
\[
L=\frac{n-1}{d},\qquad 1\le d\le64.
\]
Then either $n$ is prime or $n$ belongs to the row $\Eset(d)$ of
Table~\ref{tab:exceptions}.  In particular, if the row information is ignored,
there are $234$ distinct composite possibilities in the whole certified range.
\end{theorem}

\begin{proof}
For $d\le8$ this is the small-index classification of
Section~\ref{sec:small-index}; for $9\le d\le32$ it is
Theorem~\ref{thm:d32}; $d=33$ is Theorem~\ref{thm:d33}; and $34\le d\le64$ is covered by Theorem~\ref{thm:d64}.
\end{proof}

\begin{remark}[How to use the criterion]
This statement is stronger than a probabilistic test.  Once $L$ is known
exactly and $d\le64$, primality is decided by a finite table lookup: if a
composite candidate is not in the relevant row, it is prime.  The theorem does
not claim that computing $L$ from scratch is cheap.
\end{remark}

\section{Quantitative rarity and prime--pseudoprime spectral separation}
\label{sec:rarity}

Fixed-index finiteness is only the starting point.  To measure how sparse the
low-index composite spectrum actually is, define
\[
C(D)=\left|\bigcup_{1\le d\le D}\Eset(d)\right|,
\qquad
C^\ast(D)=\left|\bigcup_{1\le d\le D}\Eset^\ast(d)\right|.
\]
The first quantity uses the inclusive spectral convention; the second uses
Somer's nontrivial-base convention.

\begin{theorem}[Quantitative rarity through $d=64$]
\label{thm:rarity64}
The certified table gives
\[
C(64)=234,\qquad C^\ast(64)=231.
\]
Counting inherited reappearances separately, the first $64$ inclusive lines
contain $351$ spectral incidences, while Somer's nontrivial-base convention
contains $329$.
\end{theorem}

The difference between the inclusive and nontrivial-base counts can again be
audited directly.  The $22$ excluded incidences are
\[
\begin{array}{c|l}
d&\text{excluded occurrence}\\ \hline
4&9\ (L=2)\\
8&9\ (L=1)\\
12&25\ (L=2)\\
13&27\ (L=2)\\
14&15\ (L=1)\\
20&21\ (L=1)\\
24&25\ (L=1),\ 49\ (L=2)\\
26&27\ (L=1)\\
32&33\ (L=1)\\
34&35\ (L=1)\\
38&39\ (L=1)\\
40&81\ (L=2)\\
44&45\ (L=1)\\
48&49\ (L=1)\\
50&51\ (L=1)\\
54&55\ (L=1)\\
56&57\ (L=1)\\
60&121\ (L=2)\\
62&63\ (L=1),\ 125\ (L=2)\\
64&65\ (L=1).
\end{array}
\]
Order $1$ is represented only by $1\bmod n$, and for an odd prime power the
unique involution is $-1\bmod n$.  Thus these are precisely the incidences
removed by Somer's convention.  Only $9,27,81$ disappear completely from the
union through $64$; all other affected integers have another nontrivial
spectral occurrence by that cutoff.  Hence $351-22=329$ incidences and
$234-3=231$ distinct composites.

The distinction between $234$ and $351$ is important: a composite can occupy
several spectral lines by the inheritance theorem.  The title of this paper refers to the certified scarcity of the first quantity across a low-index range; the general upper bounds proved later are much weaker than the observed values.  ``Rare'' is therefore not being used as a
synonym for Somer's theorem that each single line is finite; it refers to the
observed and certified scarcity of distinct composites across an entire
low-index interval.

Figure~\ref{fig:composite-cumulative} shows the cumulative composite count.
The growth is irregular, with visible changes near the dyadic thresholds at
which one more distinct prime factor becomes structurally possible.

\begin{figure}[htbp]
\centering
\includegraphics[width=.92\textwidth]{composite_rarity_cumulative_d64.png}
\caption{Cumulative number of distinct odd composite exceptions through
spectral index $D\le64$.  The upper curve uses the inclusive group-theoretic
spectrum $\Eset(d)$; the lower curve uses Somer's nontrivial-base convention
$\Eset^\ast(d)$.}
\label{fig:composite-cumulative}
\end{figure}

On the prime side the relevant comparison is not the number of primes on one
line, since every existential prime line is infinite, but the \emph{fixed-base
prime mass} carried by the same low-index interval.  This is quantified in
Theorem~\ref{thm:prime-mass-64} below.  The resulting juxtaposition is a structural spectral comparison: the existential-over-bases composite union remains thin in the certified range, while each fixed non-perfect-power base places most of its prime residual-index mass in the same range.

The two sides of this display are deliberately not probabilities on the same
sample space.  The composite set $\Eset(d)$ is enlarged by allowing the base to
vary with $n$, whereas $\delta(b,d)$ fixes one base.  The comparison is therefore
structural rather than a direct density identity: even the union over all
possible bases contains very few low-index composite exceptions, while each
individual positive non-perfect-power base places most of its prime
residual-index mass in the same spectral range.  The line-sampling probability
introduced in Section~\ref{sec:counting-probability} later places primes and
composites in one common sample space.

The computations also suggest a dyadic stratification.  Composite spectral
lines become combinatorially richer when $d$ passes
\[
d=2^r+1,
\]
because another distinct prime factor becomes compatible with
\[
2^{s-1}\le d-1.
\]

\section{Counting composites on a spectral line and a primality probability}
\label{sec:counting-probability}

The exact low-index table suggests a natural quantitative question stronger
than fixed-$d$ finiteness.  Put
\[
e(d):=|\Eset(d)|,
\qquad
g(d):=|\mathcal G(d)|,
\]
and, for the stratification by the number of distinct prime factors,
\[
e_s(d):=
\bigl|\{n\in\Eset(d):\omega(n)=s\}\bigr|.
\]
No convincing asymptotic formula for $e(d)$ or $g(d)$ is presently known.
The data through $64$ are much too short and irregular to justify one.
Nevertheless the structural theory gives considerably more than the crude
consequence of fixed-index finiteness.

\subsection{The crude bound inherited from the size theorem}

For $d\ge2$, every composite in $\Eset(d)$ satisfies
\[
n<U_0(d):=(dS(d))^{2M(d)}
\]
and $n\equiv1\pmod d$.  Hence
\[
e(d)\ll\frac{U_0(d)}d,
\qquad
\log e(d)\le
2d\bigl(\log d+\log\log d+O(1)\bigr).
\]
This is uniform but enormously larger than the observed counts.  We now
obtain a sharper bound by counting factorization patterns directly rather
than counting all integers below the size obstruction.

\subsection{Overlap-resolved branches versus master-equation branches}

The exact search is organized by $(s,\kappa,H,P)$.  Its number of
\emph{overlap-resolved} parameter branches is
\[
B(d)=
\sum_{s=2}^{S(d)}
\sum_{\kappa\le(d-1)/2^{s-1}}
\left\lfloor\frac{d-1}{2^{s-1}\kappa}\right\rfloor\tau(\kappa).
\]
This decomposition is computationally valuable because fixing $\kappa$ and
$H$ permits exact $\lambda$-overlap pruning.

\begin{proposition}[Asymptotic overlap-resolved branch count]
\label{prop:branch-asymptotic}
As $d\to\infty$,
\[
\boxed{
B(d)=\frac12d(\log d)^2+O(d\log d).
}
\]
\end{proposition}

\begin{proof}
Put $D=d-1$ and $r=s-1$.  For
$X_r=\lfloor D/2^r\rfloor$, the inner sum is
\[
\sum_{\kappa\le X_r}\tau(\kappa)
\left\lfloor\frac{X_r}{\kappa}\right\rfloor
=
\sum_{m\le X_r}d_3(m),
\]
where $d_3=1*1*1$.  The Piltz divisor estimate
\[
\sum_{m\le x}d_3(m)
=\frac12x(\log x)^2+O(x\log x)
\]
is more than sufficient here; see, for example,
\cite[Chapter~1]{IwaniecKowalski2004}.  Summing over $r\ge1$ and using
$\sum2^{-r}=1$ together with convergence of
$\sum r2^{-r}$ and $\sum r^22^{-r}$ gives the result.
\end{proof}

For the Diophantine equation itself, however, $\kappa$ and $H$ occur only
through
\[
j=\kappa H,
\qquad
c=\frac jP.
\]
This collapses one logarithmic layer of the parameter space.

\begin{proposition}[Relaxed collapsed master-equation count]
\label{prop:master-branch-asymptotic}
Fix $s\ge2$ and put $D=d-1$.  Every actual master-equation parameter pair
$(j,P)$ satisfies
\[
1\le j\le d-1,
\qquad
2^{s-1}P\mid j.
\]
If we deliberately drop the additional arithmetic restrictions on $P$---in
particular that $P=n/\rad(n)$ is odd and is supported on the primes dividing
$n$---the number of resulting \emph{relaxed} candidate pairs is
\[
B^{+}_{\mathrm{mast},s}(d)=
\sum_{m\le D/2^{s-1}}\tau(m).
\]
Consequently the number of actual $(j,P)$ pairs is at most
$B^{+}_{\mathrm{mast},s}(d)$, and the total relaxed count is
\[
B^{+}_{\mathrm{mast}}(d)
=\sum_{s=2}^{S(d)}B^{+}_{\mathrm{mast},s}(d)
=\sum_{r\ge1}\sum_{m\le D/2^r}\tau(m).
\]
Moreover
\begin{equation}
\boxed{
B^{+}_{\mathrm{mast}}(d)
=D\log D+(2\gamma-1-2\log2)D+O(\sqrt D)
=d\log d+O(d),
}
\label{eq:master-branch-asymptotic}
\end{equation}
where $\gamma$ is Euler's constant.
\end{proposition}

\begin{proof}
Every actual branch has $P\mid\kappa$ and $2^{s-1}\mid H$, hence
$2^{s-1}P\mid j=\kappa H$.  Writing $j=2^{s-1}m$, every actual $P$ is
therefore a divisor of $m$.  If all divisors of $m$ are retained, without
checking oddness, prime support, or whether a compatible $\lambda$-overlap
really exists, there are exactly $\tau(m)$ relaxed choices.  This proves the
upper count.

The classical divisor summatory formula
\[
\sum_{m\le x}\tau(m)
=x\log x+(2\gamma-1)x+O(\sqrt x)
\]
now gives, after summing over $r=s-1\ge1$,
\[
D\log D\sum_{r\ge1}2^{-r}
-D\log2\sum_{r\ge1}r2^{-r}
+(2\gamma-1)D\sum_{r\ge1}2^{-r}
+O\!\left(\sqrt D\sum_{r\ge1}2^{-r/2}\right).
\]
Since the first two geometric sums are $1$ and $2$, respectively,
\eqref{eq:master-branch-asymptotic} follows.
\end{proof}

Thus the extra factor of order $\log d$ in the overlap-resolved search count
$B(d)$ disappears already after forgetting the particular decomposition
$j=\kappa H$.  The relaxed quantity $B^{+}_{\mathrm{mast}}(d)$ is not claimed
to count realizable arithmetic branches exactly; it is a convenient and safe
upper layer for estimating how many Diophantine candidates a spectral line can
support.

\subsection{Exact low-$\omega$ strata}

The first three strata admit particularly transparent estimates.

\begin{proposition}[Prime-power stratum]
\label{prop:e1-bound}
For an odd prime power $n=p^a$, $a\ge2$, put
\[
R_a(p)=1+p+\cdots+p^{a-1}=\frac{p^a-1}{p-1}.
\]
Then
\[
d_\lambda(p^a)=R_a(p),
\qquad
\Spec_L(p^a)=\{R_a(p)r:r\mid p-1\}.
\]
Consequently
\[
e_1(d)=
\#\{(p,a):p\text{ odd prime},\ a\ge2,
\ R_a(p)\mid d,\ d/R_a(p)\mid p-1\},
\]
and
\[
\boxed{e_1(d)=O(\tau(d)\log d)=d^{o(1)}.}
\]
\end{proposition}

\begin{proof}
Since
\[
\lambda(p^a)=p^{a-1}(p-1),
\qquad
p^a-1=(p-1)R_a(p),
\]
and $R_a(p)\equiv1\pmod p$, one has
$\gcd(p^a-1,\lambda(p^a))=p-1$.  The formula for the genuine index follows,
and spectral inheritance gives the displayed spectrum.  Finally
$R_a(p)\ge(3^a-1)/2$, so only $O(\log d)$ exponents can occur; for each
exponent the strictly increasing function $R_a(p)$ can take at most one
prime value for each divisor of $d$.
\end{proof}

For two distinct prime factors, Lemma~\ref{lem:two-prime-completion} has no
prefix: $A=B=1$, $a=d$, $b=j$, and $\Delta=d-j$.  It therefore becomes
\begin{equation}
\boxed{
((d-j)p-d)((d-j)q-d)
=c\,[d(P-1)+j].
}
\label{eq:two-prime-special}
\end{equation}
Indeed, since $cP=j\le d-1$ and $c,j\le d-1$,
\[
T=c[d(P-1)+j]=dj-cd+cj=O(d^2).
\]

\begin{theorem}[One-, two-, and three-prime counting bounds]
\label{thm:low-omega-counts}
As $d\to\infty$,
\[
\boxed{
e_1(d)\le d^{o(1)},
\qquad
e_2(d)\le d^{1+o(1)},
\qquad
e_3(d)\le d^{1+o(1)}.
}
\]
\end{theorem}

\begin{proof}
The first assertion is Proposition~\ref{prop:e1-bound}.  For $s=2$,
Proposition~\ref{prop:master-branch-asymptotic} gives $O(d\log d)$ relaxed master
parameter pairs, and hence at most that many actual pairs.  Equation~\eqref{eq:two-prime-special} has $T=O(d^2)$ as noted above, so it
supplies at most $\tau(T)/2=d^{o(1)}$ divisor-pair completions on each branch,
by the standard maximal-order estimate
\[
\tau(m)=\exp\!\left(O\!\left(\frac{\log m}{\log\log m}\right)\right).
\]
This proves $e_2(d)\le d^{1+o(1)}$.

For $s=3$, $4\mid H$ and $P\mid\kappa$, hence
\[
P\mid\kappa\mid\frac{\kappa H}{4}=\frac j4.
\]
Write $j=4m$; then $P\mid m$.  If $p=p_1$, the smallest-prime bound gives
\[
p<\frac{d+2j}{d-j},
\qquad\text{hence}\qquad
j>\frac{d(p-1)}{p+2}.
\]
Also $p<3d$.  For fixed $p$, the number of multiples of $4$ in the allowed
$j$-interval is $O(d/p+1)$, and for each $j$ there are at most
$\tau(j/4)=d^{o(1)}$ choices for $P$.  Mertens' estimate for reciprocal
primes therefore gives at most
\[
d^{o(1)}\sum_{p<3d}\left(\frac d p+1\right)=d^{1+o(1)}
\]
possible triples $(j,P,p_1)$.  For each such prefix $p_1<3d$, so
$A=p_1-1=O(d)$ and $B=p_1=O(d)$.  Hence in
Lemma~\ref{lem:two-prime-completion}
\[
a=O(d^2),\qquad b=O(d^2),\qquad c=O(d),\qquad T=O(d^4),
\]
and there are only $d^{o(1)}$ divisor-pair completions.  This proves the third
assertion.
\end{proof}

These estimates are intentionally one-sided.  They do not suggest that
$e_2(d)$ or $e_3(d)$ actually has polynomial order; the certified small-index
data are far smaller.  Their point is that the first nontrivial strata already
admit nearly linear rigorous upper bounds.

\subsection{An exploratory exact count of the low-$\omega$ strata}
\label{subsec:low-omega-experiment}

To see the scale hidden by the upper bounds, we implemented the preceding
prime-power formula and divisor-pair identities independently in Python/SymPy.
Every generated candidate is rechecked from $n$, $d$, and $\lambda(n)$ rather
than accepted from the parametrization alone.  As a regression test, the
program reproduces every $\omega(n)\le3$ incidence in the certified spectrum through $d=64$.  It was then run for every $2\le d\le500$.

\begin{center}
\begin{tabular}{c|rrrr}
\toprule
range of $d$ & mean $e_1(d)$ & mean $e_2(d)$ & mean $e_3(d)$
& mean $\sum_{s\le3}e_s(d)$\\
\midrule
$2$--$100$   & $0.556$ & $1.424$ & $4.475$ & $6.455$\\
$101$--$200$ & $0.510$ & $2.020$ & $9.710$ & $12.240$\\
$201$--$500$ & $0.420$ & $2.350$ & $13.817$ & $16.587$\\
\bottomrule
\end{tabular}
\end{center}

These are \emph{not} values of the full function $e(d)$ beyond the certified
range: strata with four or more distinct prime divisors are omitted.  Through $d=64$, however, the low-$\omega$ calculation accounts for $320$ of the $351$ certified spectral incidences: $40$ have one distinct prime factor, $71$ have two, and $209$ have three.  The remaining $31$ incidences have $\omega(n)=4$; none in the certified range has five or more distinct prime factors.  The experiment therefore
supports the qualitative impression that the actual low-order strata grow far
more slowly than the worst-case bounds.  We deliberately do not fit an
asymptotic law to the range $d\le500$: it is too short, and the dyadic changes
in the allowed value of $\omega(n)$ make extrapolation particularly unsafe.
The script and captured output are supplied separately as exploratory material;
they are not used in any proof.

\subsection{A uniform upper bound for the full composite line}

The same two-prime completion identity improves the global count even when
$s$ is allowed to grow with $d$.

\begin{theorem}[Uniform counting bound for a spectral line]
\label{thm:uniform-e-count}
As $d\to\infty$,
\begin{equation}
\boxed{
\log e(d)
\le
\frac d2\bigl(\log d+\log\log d+O(1)\bigr).
}
\label{eq:uniform-e-count}
\end{equation}
Equivalently,
\[
e(d)\le
\exp\!\left\{
\frac d2\bigl(\log d+\log\log d+O(1)\bigr)
\right\}.
\]
\end{theorem}

\begin{proof}
The strata $s=1,2$ are already smaller than the claimed bound, so suppose
$s\ge3$ and fix one overlap-resolved branch $(s,\kappa,H,P)$.  Let
\[
Q_r:=\prod_{i=1}^r p_i
\]
be the radical of a prime prefix.  The smallest-prime estimate gives
$Q_1=p_1<sd$.  For $k\ge2$, if $p_k>\kappa$, the squarefree-suffix bound applies.  Since
$K_{k-1}\mid\kappa$,
\[
\lambda(P_{k-1})\le\phi(P_{k-1})
=K_{k-1}\prod_{i<k}(p_i-1)
< K_{k-1}Q_{k-1}\le\kappa Q_{k-1}.
\]
Therefore
\[
p_k<d\lambda(P_{k-1})(s-k+1)
\le d\kappa s\,Q_{k-1}.
\]
If $p_k\le\kappa$, the same weaker inequality is automatic.  Hence
\[
Q_k<d\kappa s\,Q_{k-1}^2.
\]
Induction yields
\begin{equation}
Q_r<(sd)^{2^r-1}\kappa^{2^{r-1}-1}.
\label{eq:prefix-radical-bound}
\end{equation}

Stop with $r=s-2$ primes selected.  Since
$2^{s-1}\kappa\le d-1$, one has
\[
2^r=2^{s-2}\le\frac{d-1}{2\kappa}<\frac d{2\kappa},
\qquad
2^{r-1}=2^{s-3}<\frac d{4\kappa}.
\]
Consequently the two leading logarithmic contributions in
\eqref{eq:prefix-radical-bound} satisfy
\[
(2^r-1)\log(sd)\le
\frac d{2\kappa}\bigl(\log d+\log s\bigr),
\qquad
(2^{r-1}-1)\log\kappa\le
\frac d4\frac{\log\kappa}{\kappa}.
\]
Now $s=O(\log d)$ and $(\log\kappa)/\kappa$ is bounded for
$\kappa\ge1$; the first term is largest in leading order at $\kappa=1$.
Thus \eqref{eq:prefix-radical-bound} gives
\begin{equation}
\log Q_r
\le
\frac d2\bigl(\log d+\log\log d+O(1)\bigr).
\label{eq:prefix-count-log}
\end{equation}
On the fixed branch, the numerical value of $P$ is already fixed.  Hence an
increasing prime prefix is determined by the squarefree integer $Q_r$; for any
prefix prime dividing $P$, its exponent is fixed by $v_p(P)+1$.  The number of
possible prefixes on this branch is therefore at most the right-hand size bound
for $Q_r$ itself.  Required prime factors of $P$ that are not yet in the prefix
must occur among the final two primes and are handled by the support check in
the completion step.

For a fixed prefix, Lemma~\ref{lem:two-prime-completion} gives at most
$\tau(T)/2$ completions.  Here
\[
a=dA\le dQ_r,
\qquad
b=jB\le dQ_r,
\qquad
c\le d,
\]
so
\[
T=a(b-c)+bc\le2d^2Q_r^2.
\]
Thus $\log T=O(d\log d)$ and the maximal-order bound for the divisor
function gives $\log\tau(T)=O(d)$, lower order relative to
\eqref{eq:prefix-count-log}.  Finally the total number of overlap-resolved
branches is only $B(d)=O(d(\log d)^2)$, whose logarithm is also lower order.
Summing over all branches proves \eqref{eq:uniform-e-count}.
\end{proof}

Compared with the crude estimate inherited from the size theorem, the leading
coefficient in the exponent has fallen from $2$ to $1/2$.  This is not merely
a cosmetic improvement: after inversion it enlarges the range in which a
spectral line is overwhelmingly prime under the sampling model below.

\subsection{A natural conditional probability of primality}

There is no intrinsic probability that an integer is prime until a sampling
model is specified.  We choose uniformly from the integers on a fixed
existential spectral line $d$ up to a cutoff $X$:
\[
\mathcal S_d(X)=
\{n\le X:\exists b,\ \gcd(b,n)=1,\ \ord_n(b)=(n-1)/d\}.
\]
Its prime part is $\{p\le X:p\equiv1\pmod d\}$, while its composite part is
$\Eset(d)\cap[1,X]$.  Put
\[
\pi(X;d,1)=\#\{p\le X:p\equiv1\pmod d\}.
\]
Then
\begin{equation}
\Pr_X(n\text{ prime}\mid d)
=
\frac{\pi(X;d,1)}
{\pi(X;d,1)+|\Eset(d)\cap[1,X]|}.
\label{eq:conditional-prime-prob}
\end{equation}

\begin{theorem}[Primality probability on a fixed existential line]
\label{thm:conditional-prime-prob}
For every fixed $d\ge1$,
\[
\Pr_X(n\text{ prime}\mid d)
=
1-e(d)\phi(d)\frac{\log X}{X}
+o\!\left(\frac{\log X}{X}\right).
\]
In particular
\[
\boxed{\Pr_X(n\text{ prime}\mid d)\longrightarrow1.}
\]
\end{theorem}

\begin{proof}
For fixed $d$ and all sufficiently large $X$,
$|\Eset(d)\cap[1,X]|=e(d)$.  The prime number theorem in arithmetic
progressions gives
\[
\pi(X;d,1)\sim\frac{\operatorname{Li}(X)}{\phi(d)}
\sim\frac{X}{\phi(d)\log X}.
\]
Substitution in \eqref{eq:conditional-prime-prob} and expansion of
$1/(1+y)$ prove the formula.
\end{proof}

For fixed $d$, the convergence to $1$ is essentially the prime-number-theorem
consequence of Somer's finiteness theorem; the first-order coefficient records
how the exact finite composite count enters.  The stronger analytic use of the
new counting theorem is the growing-$d$ result below.

For $d\le64$ the leading coefficient
\[
\rho(d):=e(d)\phi(d)
\]
is known exactly.  Its maximum is
\[
\max_{1\le d\le64}\rho(d)=720,
\qquad\text{at }d=61,
\]
since $e(61)=12$ and $\phi(61)=60$.

\begin{corollary}[Uniform line-sampling estimate through $64$]
\label{cor:probability64}
For the finite family $1\le d\le64$,
\[
\min_{1\le d\le64}\Pr_X(n\text{ prime}\mid d)
=
1-720\frac{\log X}{X}
+o\!\left(\frac{\log X}{X}\right).
\]
The lines $d=1,2,3$ have $e(d)=0$ and therefore contain no odd composites.
\end{corollary}

The new counting theorem also permits $d$ to grow with the sampling cutoff.

\begin{theorem}[A uniform high-probability spectral range]
\label{thm:uniform-probability-growing-d}
For every fixed $0<\varepsilon<2$,
\begin{equation}
\boxed{
\inf_{1\le d\le(2-\varepsilon)\log X/\log\log X}
\Pr_X(n\text{ prime}\mid d)
\longrightarrow1
\qquad(X\to\infty).
}
\label{eq:uniform-probability-growing-d}
\end{equation}
\end{theorem}

\begin{proof}
Put
\[
D_X=(2-\varepsilon)\frac{\log X}{\log\log X}.
\]
Since $D_X\ll\log X$, Siegel--Walfisz gives, uniformly for $d\le D_X$,
\[
\pi(X;d,1)=\frac{X}{\phi(d)\log X}(1+o(1));
\]
see, for example, \cite[Chapter~5]{IwaniecKowalski2004}.  Choose a constant
$C$ and $d_0$ so that Theorem~\ref{thm:uniform-e-count} gives
\[
\log e(d)\le\frac d2\bigl(\log d+\log\log d+C\bigr)
\qquad(d\ge d_0).
\]
The right-hand side is increasing for sufficiently large $d$.  Hence, uniformly
for $d_0\le d\le D_X$,
\begin{align*}
\log e(d)
&\le \frac{D_X}{2}
   \bigl(\log D_X+\log\log D_X+C\bigr)\\
&=\left(1-\frac\varepsilon2+o(1)\right)\log X.
\end{align*}
The finitely many $d<d_0$ contribute only bounded constants $e(d)$ and satisfy
the same weaker estimate for large $X$.  Therefore
\[
e(d)\le X^{1-\varepsilon/2+o(1)}
\]
uniformly throughout the displayed range.  Since $\phi(d)\le d$,
\[
\frac{|\Eset(d)\cap[1,X]|}{\pi(X;d,1)}
\le X^{-\varepsilon/2+o(1)}d\log X\longrightarrow0
\]
uniformly.  Equation~\eqref{eq:conditional-prime-prob} proves the claim.
\end{proof}

\begin{remark}[Certainty versus high probability]
The constants $1/2$ and $2$ describe different statements.  The asymptotic
size theorem gives the deterministic implication
\[
d\le\left(\frac12-\varepsilon\right)
\frac{\log n}{\log\log n}
\quad\Longrightarrow\quad n\text{ is prime}
\]
for every sufficiently large candidate.  Theorem~\ref{thm:uniform-probability-growing-d}
extends to the four-times-wider leading range
\[
d\le(2-\varepsilon)\frac{\log X}{\log\log X},
\]
but only in the explicitly specified uniform line-sampling sense: composites
are asymptotically negligible rather than impossible.  Neither constant is
claimed sharp.
\end{remark}

For $d\le64$, Theorem~\ref{thm:primality64} is stronger still: once $n$ and
its exact period are known, the explicit exceptional table decides the
question rather than assigning a probability.

For a fixed base $b$, put $e_b(d)=|\Eset_b(d)|$.  Since
$\Eset_b(d)\subseteq\Eset(d)$, this number is finite.  Under the GRH
hypotheses of Theorem~\ref{thm:moree-density}, if $\delta(b,d)>0$, the same
fixed-$d$ argument gives
\[
\Pr_X(n\text{ prime}\mid b,d)
=
1-\frac{e_b(d)}{\delta(b,d)}\frac{\log X}{X}
+o\!\left(\frac{\log X}{X}\right).
\]
This formula makes explicit how the composite count and the fixed-base prime
density combine in a spectral estimate of primality.

\section{Fixed-base prime spectra and near-primitive roots}
\label{sec:fixed-base}

The existential prime spectrum of Remark~\ref{thm:prime-spectrum} allows the
base to vary with the prime.  A finer spectral problem fixes an integer
$b\notin\{-1,0,1\}$ and considers
\[
\Pset_b(d)
=
\left\{
p\nmid b:
\ord_p(b)=\frac{p-1}{d}
\right\}.
\]
In the terminology of the Artin literature, $b$ is then an
$d$-\emph{near-primitive root} modulo $p$, and $d$ is the residual index of
$b\bmod p$.

The relevant density theory is classical.  Following Hooley's method for
Artin's conjecture, Wagstaff proved the near-primitive-root density statement
under GRH; Moree subsequently derived explicit Euler-product and entanglement
correction formulas \cite{Wagstaff1982,Moree2000,Moree2013}.

\begin{theorem}[Wagstaff density theorem, with Moree correction formulas, under GRH]
\label{thm:moree-density}
Let $b\in\mathbb Q\setminus\{-1,0,1\}$ and $d\ge1$.  Assume the Riemann Hypothesis for the Dedekind zeta functions of every Kummer field
$\mathbb Q(\zeta_{kd},b^{1/(kd)})$, $k\ge1$.  Then the set
$\Pset_b(d)$ has a natural density
\[
\delta(b,d)
=
\sum_{k=1}^{\infty}
\frac{\mu(k)}
{[\mathbb Q(\zeta_{kd},b^{1/(kd)}):\mathbb Q]}.
\tag{A1}\label{eq:moree-series}
\]
This density is a rational multiple of Artin's constant.  Its vanishing is
governed by explicit entanglement conditions.
\end{theorem}

For a fixed base $b$ and a cutoff $D$, define the cumulative prime spectral
mass
\[
\mathfrak M_b(D)=\sum_{1\le d\le D}\delta(b,d).
\]
Since each prime $p\nmid b$ has one residual index, $\mathfrak M_b(D)$ is the
natural density, relative to all primes, of those primes for which
\[
\frac{p-1}{\ord_p(b)}\le D.
\]

For the spectral comparison it is useful that a large class of bases has an
especially simple odd-index spectrum.  Write
\[
b=b_0^h,
\]
where $b_0>0$ is not a perfect power in $\mathbb Q$ and $h\ge1$.  For a
positive integer $b$ that is not a perfect power, define its squarefree part by
\[
\operatorname{sf}(b):=\prod_{v_p(b)\ {\rm odd}}p.
\]
This is the squarefree representative of $b$ modulo rational squares and
should not be confused with the radical $\rad(b)$.

\begin{definition}[Odd-index generic base]
A positive integer $b>1$ is called \emph{odd-index generic} if it is not a
perfect power and
\[
\operatorname{sf}(b)\not\equiv1\pmod4.
\]
Equivalently, the quadratic discriminant of $\mathbb Q(\sqrt b)$ is even.
\end{definition}

\begin{lemma}[Explicit positive non-perfect-power correction factor, under GRH]
\label{lem:moree-h1-table}
Let $b>1$ be a positive integer that is not a perfect power, put
$g=\operatorname{sf}(b)$, and write
\[
d_2=2^{v_2(d)},
\qquad
B(g,d)=\prod_{\substack{q\mid g\\q\nmid d}}
\frac{-1}{q^2-q-1}.
\]
Then Moree's $h=1$ formula depends on $b$ only through the quadratic field
$\mathbb Q(\sqrt g)$, and the ratio $\delta(b,d)/E(d)$ is given by
\[
\begin{array}{c|c|c}
 g\pmod4 & d_2 & \delta(b,d)/E(d)\\ \hline
 1 & 1 & 1-B(g,d)\\
 1 & d_2\ge2 & 1+B(g,d)\\
 2 & d_2<4 & 1\\
 2 & d_2=4 & 1-B(g,d)/3\\
 2 & d_2>4 & 1+B(g,d)\\
 3 & d_2=1 & 1\\
 3 & d_2=2 & 1-B(g,d)/3\\
 3 & d_2\ge4 & 1+B(g,d).
\end{array}
\]
This is the positive $h=1$ specialization of Moree's explicit
near-primitive-root density formula; equivalently it is Table~1 of
\cite{Moree2013} applied to the squarefree representative of $b$ modulo
rational squares.
\end{lemma}

\begin{proof}[Derivation of Lemma~\ref{lem:moree-h1-table}]
For $b$ not a perfect power, Moree's degree calculation gives
\[
[\mathbb Q(\zeta_m,b^{1/m}):\mathbb Q]
=\frac{m\phi(m)}{\epsilon(m)},
\]
where $\epsilon(m)=2$ precisely when $2\mid m$ and the quadratic discriminant
$D=\operatorname{disc}\mathbb Q(\sqrt{\operatorname{sf}(b)})$ divides $m$;
otherwise $\epsilon(m)=1$.  Insert this dichotomy into
\eqref{eq:moree-series}.  Relative to the generic Euler product $E(d)$, the
terms for which the discriminant condition is active force a finite set $S$
of primes into the squarefree Möbius variable.  The resulting Euler-product
ratio is
\[
\prod_{p\in S}\frac{-1}{p c_p-1},
\qquad
c_p=\begin{cases}
p,&p\mid d,\\ p-1,&p\nmid d.\end{cases}
\]
The $2$-adic condition determines the eight cases in the table.  In particular
when $p=2\mid d$ the factor is $-1/(2^2-1)=-1/3$.  This gives exactly the
listed correction factors.
\end{proof}

\begin{corollary}[Universal odd-index prime spectrum, under GRH]
\label{cor:odd-universal}
Let $b$ be odd-index generic and let $d$ be odd.  Then
\[
\delta(b,d)=E(d)
=
\frac{A}{d^2}
\prod_{q\mid d}
\frac{q^2-1}{q^2-q-1}.
\tag{A2}\label{eq:Em}
\]
where $A$ is Artin's constant.  In particular, for odd spectral indices the
prime density is independent of the choice of odd-index generic base.
\end{corollary}

\begin{proof}
If $d$ is odd, then $d_2=1$.  For an odd-index generic base one has
$g=\operatorname{sf}(b)\equiv2$ or $3\pmod4$.  The corresponding $d_2=1$
rows of Lemma~\ref{lem:moree-h1-table} both give
$\delta(b,d)/E(d)=1$.
\end{proof}

\begin{theorem}[Prime-mass minimum and all-base infimum through $d=64$, under GRH]
\label{thm:prime-mass-64}
Let
\[
G_{64}:=\sum_{1\le d\le64}E(d),
\qquad
\mathfrak M_b(64):=\sum_{1\le d\le64}\delta(b,d).
\]
Then:
\begin{enumerate}[label=(\roman*)]
\item The correction-free Wagstaff mass is
\[
G_{64}=0.984529145953884713\ldots .
\]

\item Among all odd-index generic bases,
\[
\min_b\mathfrak M_b(64)
=0.982709160681865379\ldots .
\]
The minimum is attained exactly by those bases in this class whose squarefree
part is $2$.

\item More strongly, the same sharp minimum holds over all positive integer
bases $b>1$ that are not perfect powers:
\[
\mathfrak M_b(64)\ge\mathfrak M_2(64)
=0.982709160681865379\ldots .
\]
Equality holds precisely when the squarefree part of $b$ is $2$.

\item If arbitrary perfect-power bases are admitted, no positive minimum
exists.  Instead
\[
\inf_{b\ge2}\mathfrak M_b(64)=0.
\]
A sequence of powers of $2$ approaches this infimum.
\end{enumerate}
\end{theorem}

\begin{proof}
Part (i) is the direct sum of Moree's generic factors
\[
E(d)=\frac{A}{d^2}
\prod_{q\mid d}\frac{q^2-1}{q^2-q-1}.
\]

For (iii), let $g=\operatorname{sf}(b)$ for a positive non-perfect-power base.
Lemma~\ref{lem:moree-h1-table} reduces the problem to squarefree kernels
$g>1$.  Every correction has the form
\[
\delta(b,d)=E(d)\{1+\alpha(g,d)B(g,d)\},
\qquad \alpha(g,d)\in\{0,\pm\tfrac13,\pm1\}.
\]
If $g$ contains a prime $q\ge67$, then $q\nmid d$ for all $d\le64$ and
\[
|B(g,d)|\le\frac1{67^2-67-1}=\frac1{4421}.
\]
Hence
\[
|\mathfrak M_b(64)-G_{64}|\le\frac{G_{64}}{4421}.
\]
Exact rational arithmetic gives
\[
G_{64}-\mathfrak M_2(64)>\frac{G_{64}}{4421},
\]
so no minimizing kernel contains a prime at least $67$.

It remains to inspect the $2^{18}-1=262143$ nonempty squarefree products of
\[
2,3,5,7,11,13,17,19,23,29,31,37,41,43,47,53,59,61.
\]
The exact computation evaluates Lemma~\ref{lem:moree-h1-table} for all of
these kernels and finds the unique minimizer $g=2$.  The calculation is
implemented without floating-point comparisons in
\path{prime_mass_d64_check_exact.py}; its captured output is included
in the supplement.  This proves (iii), and (ii) follows because the odd-index
generic class contains the kernel $2$.

For (iv), put
\[
H_y=\prod_{\substack{67\le q\le y\\q\ {\rm prime}}}q,
\qquad b_y=2^{H_y}.
\]
For every $d\le64$ and every prime $q\mid H_y$ one has $q\nmid2d$.
In Moree's Euler-product formula this is exactly the perfect-power factor
$q(q-2)/(q^2-q-1)$; see \cite[Eq.~(8)]{Moree2013}, specifically the factor
attached to $q\mid h$ and $q\nmid mt$.  Consequently
\[
\mathfrak M_{b_y}(64)=\mathfrak M_2(64)
\prod_{\substack{67\le q\le y\\q\ {\rm prime}}}
\frac{q(q-2)}{q^2-q-1}.
\]
The factors are $1-1/q+O(q^{-2})$, so the product tends to zero.  This proves
(iv).
\end{proof}

\begin{remark}[A fixed-base prime-side comparison through $64$]
\label{cor:low-index-separation}
Under GRH, the same interval $1\le d\le64$ which contains only
\[
C^\ast(64)=231
\]
distinct odd composite exceptions in Somer's nontrivial-base convention
contains at least
\[
98.2709160681865\ldots\%
\]
of the fixed-base prime mass for every positive non-perfect-power integer
base.  For correction-free Wagstaff behavior the corresponding prime mass is
\[
98.4529145953885\ldots\%.
\]
These percentages are fixed-base prime densities; $C^\ast(64)$ is an existential-over-bases composite cardinality.  The display is therefore a side-by-side structural comparison, not a probability theorem.
\end{remark}

\begin{remark}
The prime-mass theorem clarifies two different uses of the word ``generic''.  For
correction-free Wagstaff behavior the mass through $64$ is
$98.452915\ldots\%$.  Over the larger odd-index generic class the worst
case is the squarefree part $2$, giving the sharp
$98.270916\ldots\%$.  Over all bases, including arbitrarily high perfect
powers, the infimum is $0$.  Thus the correction-free value $98.452915\ldots\%$ is not a universal lower bound over all non-perfect-power bases, while the sharp universal value $98.270916\ldots\%$ is.
\end{remark}

\begin{figure}[htbp]
\centering
\includegraphics[width=.92\textwidth]{prime_mass_cumulative_d64.png}
\caption{Cumulative fixed-base prime mass through spectral index $D$.
The upper curve is the correction-free Wagstaff mass
$\sum_{d\le D}E(d)$.  The lower curve is $\mathfrak M_2(D)$, the cumulative
mass for squarefree kernel $2$; its endpoint at $D=64$ is the proven global
minimum among positive non-perfect-power bases.  No claim is made here that
kernel $2$ minimizes every intermediate cutoff $D$.  The dashed horizontal
line, labeled $0.98$, marks $98\%$.}
\label{fig:prime-mass-cumulative}
\end{figure}

For example,
\[
E(5)=
\frac{A}{25}\frac{24}{19}
\approx0.0189,
\]
so approximately $1.89\%$ of all primes lie on the fixed-base spectral line
$d=5$ for every odd-index generic base.

The distinction between the two prime spectra should be kept explicit:
\[
\underbrace{\frac1{\phi(d)}}_{\text{existential over bases}}
\qquad\text{versus}\qquad
\underbrace{\delta(b,d)}_{\text{one fixed base}}.
\]
The first quantity is unconditional and is the correct comparison with the
union-over-bases composite spectrum $\Eset(d)$.  The second is a refined,
fixed-base Artin spectrum and is conditional on GRH in the present state of
the theory.

\section{The range beyond $d=64$}

The certified theorem in this paper now stops at $d=64$.  No exploratory table for
larger $d$ is used as evidence or as an input to any proof.  Extending the exact
branch search to $34\le d\le128$ is a natural next computation, with the next
dyadic transition occurring at $d=65$.

\section{Further structural refinements}

The overlap recurrence \eqref{eq:overlap}, introduced above as a certified
search invariant, also has a useful structural interpretation.  With
\[
c_i=p_i^{a_i-1}(p_i-1),
\qquad
\Lambda_k=\operatorname{lcm}(c_1,\ldots,c_k),
\qquad
H_k=\frac{\prod_{i=1}^k c_i}{\Lambda_k},
\]
one has $H_s=\phi(n)/\lambda(n)$, so every prefix consumes part of a finite
``overlap budget''.  In the minimal-overlap case
\[
H_s=2^{s-1},
\]
one obtains
\[
\gcd(c_i,c_j)=2
\qquad(i\ne j).
\]
For squarefree $n$ this reduces to
\[
\gcd(p_i-1,p_j-1)=2.
\]

The terminal prime is also strongly constrained.  If $p_s>\kappa$, then
\[
p_s-1\mid \kappa(P_{s-1}-1).
\]
Together with \eqref{eq:overlap}, this turns final-prime completion into
a finite divisor problem.  These refinements are central to the ongoing
extension beyond $d=64$.

\section{Discussion and open problems}

The spectral formulation suggests several problems.

\begin{enumerate}
\item
Determine the exact growth of
\[
N_{\max}(d)
=
\max\Eset(d),
\]
with the convention $N_{\max}(d)=0$ when $\Eset(d)=\varnothing$.
The explicit function $U_0(d)$ of
Theorem~\ref{thm:uniform-U} proves finiteness but is far
from optimal.

\item
Determine the asymptotic behavior of
\[
e(d)=|\Eset(d)|.
\]
The data for $d\le64$ show strong fluctuations and repeated inherited
spectral lines.

\item
Determine the asymptotic behavior of the genuine count
\[
g(d)=|\mathcal G(d)|
\]
and compare it with the inclusive incidence count $e(d)$.  By
Corollary~\ref{cor:genuine-partition}, $g(d)$ is the distribution of the
canonical genuine index across all odd composites.

\item
Extend the certified classification beyond $d=64$.  The next dyadic threshold is $65$, followed by $129,\ldots$; these are natural complexity transitions because another distinct prime factor becomes possible.

\item
Develop the $\phi$-spectrum.  Its finiteness is substantially more
delicate and contains Lehmer-type problems.

\item
For fixed base $b$, determine and compare the prime density
$\delta_b(d)$ and the fixed-base pseudoprime count
\[
|\Eset_b(d)|.
\]
This is the most direct fixed-base version of the spectral-separation
program.
\end{enumerate}


\subsection{Four natural OEIS sequences}
\label{subsec:oeis}

The certified rows now define substantially longer prefixes of four natural
integer sequences.  A search of the current OEIS finds many fixed-base
pseudoprime, Carmichael and fixed-order sequences, but we have not located a
direct entry for Somer's variable-base Fermat $d$-pseudoprime spectrum in the
inclusive form used here.  For example, OEIS A086249 counts \emph{base-$2$}
Fermat pseudoprimes by the value of $\ord_n(2)$; that fixed-base, fixed-order
statistic differs from the variable-base reciprocal index studied here.  The
following prefixes are now certified through $d=64$ with no numerical cutoff
on $n$.

Let
\[
\mathcal G(d)=\{n:d_\lambda(n)=d\}.
\]
The four sequences are as follows.

\begin{enumerate}
\item \textbf{Inclusive spectral list, retaining inherited repetitions.}
Read the rows $\Eset(1),\Eset(2),\ldots$ in increasing $d$, sorting each row
increasingly.  Through $d=64$ the complete certified prefix is
\begin{quote}\small\sloppy
9,\allowbreak\ 6601,\allowbreak\ 25,\allowbreak\ 15,\allowbreak\ 561,\allowbreak\ 9,\allowbreak\
49,\allowbreak\ 30889,\allowbreak\ 88561,\allowbreak\ 534061,\allowbreak\ 21,\allowbreak\ 6601,\allowbreak\
45,\allowbreak\ 25,\allowbreak\ 121,\allowbreak\ 27,\allowbreak\ 126673,\allowbreak\ 1333333,\allowbreak\
15,\allowbreak\ 169,\allowbreak\ 561,\allowbreak\ 91,\allowbreak\ 6601,\allowbreak\ 33,\allowbreak\
49,\allowbreak\ 65,\allowbreak\ 35,\allowbreak\ 139231,\allowbreak\ 289,\allowbreak\ 30889,\allowbreak\
88561,\allowbreak\ 534061,\allowbreak\ 55462177,\allowbreak\ 8885251441,\allowbreak\ 42018333841,\allowbreak\
39,\allowbreak\ 153,\allowbreak\ 4371,\allowbreak\ 7107,\allowbreak\ 21,\allowbreak\ 361,\allowbreak\
3201,\allowbreak\ 6601,\allowbreak\ 85,\allowbreak\ 8695,\allowbreak\ 10585,\allowbreak\ 26335,\allowbreak\
9717247,\allowbreak\ 16666651,\allowbreak\ 37769887,\allowbreak\ 45,\allowbreak\ 133,\allowbreak\
2465,\allowbreak\ 231,\allowbreak\ 645,\allowbreak\ 1105,\allowbreak\ 24013,\allowbreak\ 34133,\allowbreak\
25,\allowbreak\ 49,\allowbreak\ 121,\allowbreak\ 529,\allowbreak\ 51,\allowbreak\ 6601,\allowbreak\
15051,\allowbreak\ 11921001,\allowbreak\ 27,\allowbreak\ 105,\allowbreak\ 126673,\allowbreak\
1333333,\allowbreak\ 55,\allowbreak\ 325,\allowbreak\ 3565,\allowbreak\ 30889,\allowbreak\ 88561,\allowbreak\
534061,\allowbreak\ 57,\allowbreak\ 169,\allowbreak\ 561,\allowbreak\ 117,\allowbreak\ 175,\allowbreak\
1045,\allowbreak\ 1447681,\allowbreak\ 2121061,\allowbreak\ 164424781,\allowbreak\ 91,\allowbreak\
841,\allowbreak\ 6601,\allowbreak\ 158401,\allowbreak\ 63,\allowbreak\ 125,\allowbreak\ 435,\allowbreak\
62745,\allowbreak\ 1154131,\allowbreak\ 2609581,\allowbreak\ 3936691,\allowbreak\ 5624641,\allowbreak\
33,\allowbreak\ 65,\allowbreak\ 385,\allowbreak\ 961,\allowbreak\ 12673,\allowbreak\ 12801,\allowbreak\
1047619,\allowbreak\ 2210671,\allowbreak\ 3494701,\allowbreak\ 15197491,\allowbreak\ 25158871,\allowbreak\
93755971,\allowbreak\ 35,\allowbreak\ 69,\allowbreak\ 341,\allowbreak\ 139231,\allowbreak\ 561,\allowbreak\
91001,\allowbreak\ 610051,\allowbreak\ 683761,\allowbreak\ 145,\allowbreak\ 217,\allowbreak\ 289,\allowbreak\
30889,\allowbreak\ 88561,\allowbreak\ 534061,\allowbreak\ 55462177,\allowbreak\ 8885251441,\allowbreak\
42018333841,\allowbreak\ 75,\allowbreak\ 17169,\allowbreak\ 33153,\allowbreak\ 45141,\allowbreak\
62605,\allowbreak\ 2415953,\allowbreak\ 2929291,\allowbreak\ 11972017,\allowbreak\ 16894571,\allowbreak\
67902031,\allowbreak\ 191834567,\allowbreak\ 39,\allowbreak\ 77,\allowbreak\ 153,\allowbreak\
1065,\allowbreak\ 1369,\allowbreak\ 4371,\allowbreak\ 7107,\allowbreak\ 13833,\allowbreak\
1170401,\allowbreak\ 703,\allowbreak\ 15211,\allowbreak\ 126673,\allowbreak\ 334153,\allowbreak\
1333333,\allowbreak\ 81,\allowbreak\ 361,\allowbreak\ 481,\allowbreak\ 3201,\allowbreak\ 6601,\allowbreak\
460801,\allowbreak\ 1372801,\allowbreak\ 165,\allowbreak\ 247,\allowbreak\ 2871,\allowbreak\
8365,\allowbreak\ 73555,\allowbreak\ 184255,\allowbreak\ 9345541,\allowbreak\ 128950249,\allowbreak\
85,\allowbreak\ 169,\allowbreak\ 1681,\allowbreak\ 8695,\allowbreak\ 10585,\allowbreak\ 26335,\allowbreak\
9717247,\allowbreak\ 16666651,\allowbreak\ 37769887,\allowbreak\ 87,\allowbreak\ 259,\allowbreak\
861,\allowbreak\ 22791,\allowbreak\ 52633,\allowbreak\ 69231,\allowbreak\ 82733,\allowbreak\ 45,\allowbreak\
133,\allowbreak\ 1849,\allowbreak\ 2465,\allowbreak\ 15841,\allowbreak\ 91,\allowbreak\ 451,\allowbreak\
8911,\allowbreak\ 88561,\allowbreak\ 534061,\allowbreak\ 3763801,\allowbreak\ 14245741,\allowbreak\
93,\allowbreak\ 185,\allowbreak\ 231,\allowbreak\ 369,\allowbreak\ 645,\allowbreak\ 1105,\allowbreak\
24013,\allowbreak\ 34133,\allowbreak\ 95,\allowbreak\ 1035,\allowbreak\ 2821,\allowbreak\ 4795,\allowbreak\
204263,\allowbreak\ 223721,\allowbreak\ 308039,\allowbreak\ 559489,\allowbreak\ 38649041,\allowbreak\
49,\allowbreak\ 529,\allowbreak\ 1729,\allowbreak\ 2209,\allowbreak\ 52417,\allowbreak\ 99,\allowbreak\
87676681,\allowbreak\ 133861729,\allowbreak\ 161005573,\allowbreak\ 258844951,\allowbreak\
1411040359,\allowbreak\ 3462426241,\allowbreak\ 7950140448251,\allowbreak\ 131097700645913,\allowbreak\
798178197617711,\allowbreak\ 989473937059427,\allowbreak\ 1208361237478669,\allowbreak\ 51,\allowbreak\
301,\allowbreak\ 1001,\allowbreak\ 6601,\allowbreak\ 15051,\allowbreak\ 3679201,\allowbreak\
11921001,\allowbreak\ 4199932801,\allowbreak\ 1353385637801,\allowbreak\ 205,\allowbreak\ 32131,\allowbreak\
139231,\allowbreak\ 105,\allowbreak\ 833,\allowbreak\ 9361,\allowbreak\ 126673,\allowbreak\
741937,\allowbreak\ 1333333,\allowbreak\ 206955841,\allowbreak\ 560630848961,\allowbreak\ 425,\allowbreak\
637,\allowbreak\ 3605,\allowbreak\ 8481,\allowbreak\ 55651,\allowbreak\ 3547291,\allowbreak\
3984541,\allowbreak\ 43523707,\allowbreak\ 52519291,\allowbreak\ 260410837,\allowbreak\
1271325841,\allowbreak\ 152674725071,\allowbreak\ 3916268927681,\allowbreak\ 55,\allowbreak\ 325,\allowbreak\
2809,\allowbreak\ 3565,\allowbreak\ 30889,\allowbreak\ 88561,\allowbreak\ 534061,\allowbreak\
4169867689,\allowbreak\ 111,\allowbreak\ 221,\allowbreak\ 1431,\allowbreak\ 6601,\allowbreak\
15181,\allowbreak\ 100101,\allowbreak\ 689701,\allowbreak\ 1269621,\allowbreak\ 271794601,\allowbreak\
568898551,\allowbreak\ 57,\allowbreak\ 169,\allowbreak\ 225,\allowbreak\ 561,\allowbreak\ 25761,\allowbreak\
376993,\allowbreak\ 115,\allowbreak\ 343,\allowbreak\ 2737,\allowbreak\ 95761,\allowbreak\
37883341,\allowbreak\ 713798461,\allowbreak\ 117,\allowbreak\ 175,\allowbreak\ 1045,\allowbreak\
2553,\allowbreak\ 80041,\allowbreak\ 1447681,\allowbreak\ 2121061,\allowbreak\ 8372881,\allowbreak\
164424781,\allowbreak\ 119,\allowbreak\ 1653,\allowbreak\ 1771,\allowbreak\ 42353741,\allowbreak\
121,\allowbreak\ 361,\allowbreak\ 841,\allowbreak\ 3481,\allowbreak\ 6601,\allowbreak\ 25201,\allowbreak\
158401,\allowbreak\ 6840001,\allowbreak\ 123,\allowbreak\ 245,\allowbreak\ 1221,\allowbreak\
44409,\allowbreak\ 50997,\allowbreak\ 74665,\allowbreak\ 631351,\allowbreak\ 964411,\allowbreak\
1102759,\allowbreak\ 940724555,\allowbreak\ 1962804565,\allowbreak\ 2417562005,\allowbreak\ 63,\allowbreak\
125,\allowbreak\ 435,\allowbreak\ 3721,\allowbreak\ 19345,\allowbreak\ 62745,\allowbreak\ 264121,\allowbreak\
828073,\allowbreak\ 1154131,\allowbreak\ 2609581,\allowbreak\ 3936691,\allowbreak\ 5624641,\allowbreak\
7697921,\allowbreak\ 1891,\allowbreak\ 8695,\allowbreak\ 10585,\allowbreak\ 26335,\allowbreak\
9717247,\allowbreak\ 16666651,\allowbreak\ 37769887,\allowbreak\ 65,\allowbreak\ 129,\allowbreak\
385,\allowbreak\ 961,\allowbreak\ 1281,\allowbreak\ 12673,\allowbreak\ 12801.
\end{quote}
There are $351$ displayed terms because inherited occurrences are retained.

\item \textbf{Genuine spectral list.}
Read the rows $\mathcal G(1),\mathcal G(2),\ldots$ in increasing $d$.  By
Corollary~\ref{cor:genuine-partition}, no odd composite occurs twice in this
sequence.  Through $d=64$ the complete certified prefix is
\begin{quote}\small\sloppy
9,\allowbreak\ 6601,\allowbreak\ 25,\allowbreak\ 15,\allowbreak\ 561,\allowbreak\ 49,\allowbreak\
30889,\allowbreak\ 88561,\allowbreak\ 534061,\allowbreak\ 21,\allowbreak\ 45,\allowbreak\ 121,\allowbreak\
27,\allowbreak\ 126673,\allowbreak\ 1333333,\allowbreak\ 169,\allowbreak\ 91,\allowbreak\ 33,\allowbreak\
65,\allowbreak\ 35,\allowbreak\ 139231,\allowbreak\ 289,\allowbreak\ 55462177,\allowbreak\
8885251441,\allowbreak\ 42018333841,\allowbreak\ 39,\allowbreak\ 153,\allowbreak\ 4371,\allowbreak\
7107,\allowbreak\ 361,\allowbreak\ 3201,\allowbreak\ 85,\allowbreak\ 8695,\allowbreak\ 10585,\allowbreak\
26335,\allowbreak\ 9717247,\allowbreak\ 16666651,\allowbreak\ 37769887,\allowbreak\ 133,\allowbreak\
2465,\allowbreak\ 231,\allowbreak\ 645,\allowbreak\ 1105,\allowbreak\ 24013,\allowbreak\ 34133,\allowbreak\
529,\allowbreak\ 51,\allowbreak\ 15051,\allowbreak\ 11921001,\allowbreak\ 105,\allowbreak\ 55,\allowbreak\
325,\allowbreak\ 3565,\allowbreak\ 57,\allowbreak\ 117,\allowbreak\ 175,\allowbreak\ 1045,\allowbreak\
1447681,\allowbreak\ 2121061,\allowbreak\ 164424781,\allowbreak\ 841,\allowbreak\ 158401,\allowbreak\
63,\allowbreak\ 125,\allowbreak\ 435,\allowbreak\ 62745,\allowbreak\ 1154131,\allowbreak\
2609581,\allowbreak\ 3936691,\allowbreak\ 5624641,\allowbreak\ 385,\allowbreak\ 961,\allowbreak\
12673,\allowbreak\ 12801,\allowbreak\ 1047619,\allowbreak\ 2210671,\allowbreak\ 3494701,\allowbreak\
15197491,\allowbreak\ 25158871,\allowbreak\ 93755971,\allowbreak\ 69,\allowbreak\ 341,\allowbreak\
91001,\allowbreak\ 610051,\allowbreak\ 683761,\allowbreak\ 145,\allowbreak\ 217,\allowbreak\ 75,\allowbreak\
17169,\allowbreak\ 33153,\allowbreak\ 45141,\allowbreak\ 62605,\allowbreak\ 2415953,\allowbreak\
2929291,\allowbreak\ 11972017,\allowbreak\ 16894571,\allowbreak\ 67902031,\allowbreak\ 191834567,\allowbreak\
77,\allowbreak\ 1065,\allowbreak\ 1369,\allowbreak\ 13833,\allowbreak\ 1170401,\allowbreak\ 703,\allowbreak\
15211,\allowbreak\ 334153,\allowbreak\ 81,\allowbreak\ 481,\allowbreak\ 460801,\allowbreak\
1372801,\allowbreak\ 165,\allowbreak\ 247,\allowbreak\ 2871,\allowbreak\ 8365,\allowbreak\ 73555,\allowbreak\
184255,\allowbreak\ 9345541,\allowbreak\ 128950249,\allowbreak\ 1681,\allowbreak\ 87,\allowbreak\
259,\allowbreak\ 861,\allowbreak\ 22791,\allowbreak\ 52633,\allowbreak\ 69231,\allowbreak\ 82733,\allowbreak\
1849,\allowbreak\ 15841,\allowbreak\ 451,\allowbreak\ 8911,\allowbreak\ 3763801,\allowbreak\
14245741,\allowbreak\ 93,\allowbreak\ 185,\allowbreak\ 369,\allowbreak\ 95,\allowbreak\ 1035,\allowbreak\
2821,\allowbreak\ 4795,\allowbreak\ 204263,\allowbreak\ 223721,\allowbreak\ 308039,\allowbreak\
559489,\allowbreak\ 38649041,\allowbreak\ 1729,\allowbreak\ 2209,\allowbreak\ 52417,\allowbreak\
99,\allowbreak\ 87676681,\allowbreak\ 133861729,\allowbreak\ 161005573,\allowbreak\ 258844951,\allowbreak\
1411040359,\allowbreak\ 3462426241,\allowbreak\ 7950140448251,\allowbreak\ 131097700645913,\allowbreak\
798178197617711,\allowbreak\ 989473937059427,\allowbreak\ 1208361237478669,\allowbreak\ 301,\allowbreak\
1001,\allowbreak\ 3679201,\allowbreak\ 4199932801,\allowbreak\ 1353385637801,\allowbreak\ 205,\allowbreak\
32131,\allowbreak\ 833,\allowbreak\ 9361,\allowbreak\ 741937,\allowbreak\ 206955841,\allowbreak\
560630848961,\allowbreak\ 425,\allowbreak\ 637,\allowbreak\ 3605,\allowbreak\ 8481,\allowbreak\
55651,\allowbreak\ 3547291,\allowbreak\ 3984541,\allowbreak\ 43523707,\allowbreak\ 52519291,\allowbreak\
260410837,\allowbreak\ 1271325841,\allowbreak\ 152674725071,\allowbreak\ 3916268927681,\allowbreak\
2809,\allowbreak\ 4169867689,\allowbreak\ 111,\allowbreak\ 221,\allowbreak\ 1431,\allowbreak\
15181,\allowbreak\ 100101,\allowbreak\ 689701,\allowbreak\ 1269621,\allowbreak\ 271794601,\allowbreak\
568898551,\allowbreak\ 225,\allowbreak\ 25761,\allowbreak\ 376993,\allowbreak\ 115,\allowbreak\
343,\allowbreak\ 2737,\allowbreak\ 95761,\allowbreak\ 37883341,\allowbreak\ 713798461,\allowbreak\
2553,\allowbreak\ 80041,\allowbreak\ 8372881,\allowbreak\ 119,\allowbreak\ 1653,\allowbreak\
1771,\allowbreak\ 42353741,\allowbreak\ 3481,\allowbreak\ 25201,\allowbreak\ 6840001,\allowbreak\
123,\allowbreak\ 245,\allowbreak\ 1221,\allowbreak\ 44409,\allowbreak\ 50997,\allowbreak\ 74665,\allowbreak\
631351,\allowbreak\ 964411,\allowbreak\ 1102759,\allowbreak\ 940724555,\allowbreak\ 1962804565,\allowbreak\
2417562005,\allowbreak\ 3721,\allowbreak\ 19345,\allowbreak\ 264121,\allowbreak\ 828073,\allowbreak\
7697921,\allowbreak\ 1891,\allowbreak\ 129,\allowbreak\ 1281.
\end{quote}
It contains the $234$ distinct composites that occur in the inclusive spectrum
through $d=64$.

\item \textbf{Inclusive spectral counts.}
\[
e(d)=|\Eset(d)|.
\]
For $1\le d\le64$:
\[
\begin{split}
0,0,0,1,1,1,2,2,3,2,1,2,3,3,2,3,\\
2,7,4,4,7,3,5,4,4,4,6,3,6,4,8,6,\\
6,4,4,9,11,9,5,7,8,9,7,5,7,8,9,5,\\
12,9,3,8,13,8,10,6,6,9,4,8,12,13,7,7
\end{split}
\]

\item \textbf{Genuine spectral counts.}
\[
g(d)=|\mathcal G(d)|.
\]
For $1\le d\le64$:
\[
\begin{split}
0,0,0,1,1,1,2,1,3,1,1,1,3,1,1,2,\\
2,4,4,2,7,2,5,1,3,1,3,1,6,2,8,4,\\
6,2,3,2,11,5,3,4,8,1,7,2,4,3,9,3,\\
12,5,2,5,13,2,9,3,6,3,4,3,12,5,1,2
\end{split}
\]
\end{enumerate}

The two list sequences should be cross-referenced to the two count sequences,
since the counts recover the row boundaries, including the empty rows
$d=1,2,3$.  All four use the \emph{inclusive} spectrum.  An OEIS submission
should state explicitly that the witnessing base may vary with $n$, that
inherited and genuine appearances are different notions, and that the
inclusive convention allows the residue classes $\pm1$, unlike Somer's
nontrivial-base convention.  The supplementary machine-readable
\texttt{combined\_d1\_64.json} is the preferred source for preparing OEIS
b-files, because it preserves the row boundaries and exact integer values.

\section{Conclusion}

The starting question of this extended version is not initially spectral:
if an exact reciprocal period
\[
L=\ord_n(b)
\]
is known, what can it tell us about primality?  The answer is a collection of
classical and newer tools.  At one extreme, $L=n-1$ gives Lucas' immediate
primality criterion.  More generally, the exact order supplies Fermat and
strong-pseudoprime information, Midy decompositions, cyclotomic structure,
and sometimes enough of $n-1$ to complete a Pocklington or
Brillhart--Lehmer--Selfridge certificate.  The toolbox also makes clear its
own limitation: obtaining an exact order can be substantial computational
input.

The spectral investigation begins when $L\mid n-1$ and
\[
d=\frac{n-1}{L}.
\]
It then gives two direct extensions of the Lucas idea.  The exact theorem through $d=64$ says that an odd $n$ with $L=(n-1)/d$ is prime unless it belongs to one explicitly certified row of Table~\ref{tab:exceptions}; the $64$ rows contain $234$ distinct composite integers in $351$ spectral incidences.  The uniform theorem says
that every composite spectral point satisfies
\[
n<(dS(d))^{2M(d)},
\]
and therefore sufficiently small $d$ relative to $\log n/\log\log n$ forces
primality without any finite table.

The genuine spectrum adds a complementary base-independent structure.  Every
odd composite has exactly one genuine index
\[
d_\lambda(n)=\frac{n-1}{\gcd(n-1,\lambda(n))},
\]
so the sets $\mathcal G(d)$ partition all odd composites.  The counting
problem is now more quantitative than in the preceding drafts.  The
$\lambda$-overlap search has
$B(d)=\frac12d(\log d)^2+O(d\log d)$ parameter branches, while a relaxed collapsed master-equation count has only
$B_{\mathrm{mast}}^{+}(d)=d\log d+O(d)$ candidate pairs.  Solving
the last two primes by a divisor-pair identity gives
\[
\log e(d)\le
\frac d2(\log d+\log\log d+O(1)).
\]
Under the natural line-sampling model this implies not only that the
conditional primality probability tends to $1$ for every fixed $d$, but
uniformly, for each fixed $0<\varepsilon<2$, throughout
\[
d\le(2-\varepsilon)\frac{\log X}{\log\log X}.
\]
Thus the deterministic certainty threshold with leading constant $1/2$ and
the high-probability threshold with leading constant $2$ give complementary
ways in which a long reciprocal period signals primality.

The prime side exhibits a different kind of concentration.  Under GRH, the
fixed-base Wagstaff--Moree theory places at least
$98.2709160681865\ldots\%$ of the prime mass in $d\le64$ for every positive
integer base that is not a perfect power; the correction-free mass is
$98.4529145953885\ldots\%$.  Thus the same low-index interval contains most of
the prime mass but only a very small certified composite population.

Historically, the reciprocal index was developed here from the period--primality
question before Somer's 1985 dissertation and 1989 paper were located through
an AI-assisted literature search.  Somer's earlier work establishes the
correct priority for fixed-$d$ finiteness and the $d\le8$ classification.
A targeted citation search shows real but relatively narrow independent uptake,
most visibly in Ribenboim's dedicated ``Somer Pseudoprimes'' subsection and
Pinch's bounded-index Carmichael work.  We therefore regard the Fermat
$d$-pseudoprime theory as under-recognized rather than forgotten.  The
independent rediscovery from a different motivation nevertheless suggests that
$(n-1)/L$ is a natural arithmetic invariant.  The main questions now are
quantitative: determine the true growth of $e(d)$ and $g(d)$, sharpen the
uniform primality boundary, extend the exact spectrum beyond $33$, and
understand how the $L$-, $\lambda$-, $\phi$-, and fixed-base Artin spectra fit
together.

\section*{Declaration on computational and AI assistance}

During the development of this work, generative-AI tools, principally OpenAI
ChatGPT, were used as interactive aids for exploratory algebra, code drafting,
literature discovery, and manuscript editing.  In particular, an AI-assisted
literature search first brought Somer's 1985 dissertation and 1989 Fermat
$d$-pseudoprime paper to the authors' attention after the reciprocal-index
viewpoint and initial small-index arguments had already been developed
independently.  This historical sequence is stated to explain the development
of the project, not to claim priority over Somer's results.

All literature claims used mathematically are tied to human-authored sources.
All computer-assisted classification claims are tied to the frozen Julia
source, captured execution logs, and separate arithmetic verifier described in
Subsection~\ref{subsec:reproducibility}.  The human authors retain full
responsibility for the mathematical statements, citations, and the
verification of every theorem and computational claim in this manuscript.

\begin{thebibliography}{99}

\bibitem{APR1983}
L.~M. Adleman, C.~Pomerance, and R.~S. Rumely,
\emph{On distinguishing prime numbers from composite numbers},
Ann. of Math. (2) \textbf{117} (1983), 173--206.

\bibitem{AKS2004}
M.~Agrawal, N.~Kayal, and N.~Saxena,
\emph{PRIMES is in P},
Ann. of Math. \textbf{160} (2004), 781--793.

\bibitem{AGP1994}
W.~R. Alford, A.~Granville, and C.~Pomerance,
\emph{There are infinitely many Carmichael numbers},
Ann. of Math. 140 (1994), 703--722.

\bibitem{AtkinMorain1993}
A.~O.~L. Atkin and F.~Morain,
\emph{Elliptic curves and primality proving},
Math. Comp. \textbf{61} (1993), 29--68,
doi:10.1090/S0025-5718-1993-1199989-X.

\bibitem{BLS1975}
J.~Brillhart, D.~H. Lehmer, and J.~L. Selfridge,
\emph{New primality criteria and factorizations of $2^m\pm1$},
Math. Comp. 29 (1975), 620--647.

\bibitem{CarlipJacobsonSomer1998}
W.~Carlip, E.~Jacobson, and L.~Somer,
\emph{Pseudoprimes, perfect numbers, and a problem of Lehmer},
Fibonacci Quart. \textbf{36} (1998), 361--371.

\bibitem{CarlipSomer2007}
W.~Carlip and L.~Somer,
\emph{Square-free Lucas $d$-pseudoprimes and Carmichael--Lucas numbers},
Czechoslovak Math. J. \textbf{57} (2007), 447--463.

\bibitem{Castillo2012}
J.~H. Castillo, G.~Garc\'ia-Pulgar\'in, and J.~M. Vel\'asquez-Soto,
\emph{Pseudoprimes stronger than strong pseudoprimes},
arXiv:1202.3428, 2012.

\bibitem{Castillo2015}
J.~H. Castillo, G.~Garc\'ia-Pulgar\'in, and J.~M. Vel\'asquez-Soto,
\emph{De los n\'umeros de Midy a la primalidad},
Rev. Integr. Temas Mat. \textbf{33} (2015), 1--10,
doi:10.18273/revint.v33n1-2015001.

\bibitem{CohenLenstra1984}
H.~Cohen and H.~W. Lenstra, Jr.,
\emph{Primality testing and Jacobi sums},
Math. Comp. \textbf{42} (1984), 297--330.

\bibitem{CrandallPomerance}
R.~Crandall and C.~Pomerance,
\emph{Prime Numbers: A Computational Perspective},
2nd ed., Springer, 2005.

\bibitem{GarciaGiraldo2009}
G.~Garc\'ia-Pulgar\'in and H.~Giraldo,
\emph{Characterizations of Midy's property},
Integers 9 (2009), 191--197.

\bibitem{GuptaSury2005}
A.~Gupta and B.~Sury,
\emph{Decimal expansion of $1/p$ and subgroup sums},
Integers 5 (2005), A19.

\bibitem{Hooley1967}
C.~Hooley,
\emph{On Artin's conjecture},
J. Reine Angew. Math. \textbf{225} (1967), 209--220.

\bibitem{IwaniecKowalski2004}
H.~Iwaniec and E.~Kowalski,
\emph{Analytic Number Theory},
American Mathematical Society Colloquium Publications, Vol.~53,
AMS, Providence, RI, 2004.

\bibitem{KamadaPhi10}
M.~Kamada,
\emph{Factorizations of $\Phi_n(10)$},
STUDIO KAMADA, continuously updated online factor tables,\newline\url{https://stdkmd.net/nrr/repunit/phin10.htm}.

\bibitem{KamadaPRP}
M.~Kamada,
\emph{Probable prime factors of repunit},
STUDIO KAMADA, continuously updated online appendix,\newline\url{https://stdkmd.net/nrr/repunit/prpfactors.htm}.

\bibitem{Lehmer1932}
D.~H. Lehmer,
\emph{On Euler's totient function},
Bull. Amer. Math. Soc. 38 (1932), 745--751.

\bibitem{Lewittes2007}
J.~Lewittes,
\emph{Midy's theorem for periodic decimals},
Integers 7 (2007), A02.

\bibitem{Martin2007}
H.~W. Martin,
\emph{Generalizations of Midy's theorem on repeating decimals},
Integers 7 (2007), A03.

\bibitem{Moree2000}
P.~Moree,
\emph{Asymptotically exact heuristics for (near) primitive roots},
J. Number Theory 83 (2000), 155--181.

\bibitem{Moree2013}
P.~Moree,
\emph{Near-primitive roots},
Funct. Approx. Comment. Math. 48 (2013), 133--145,
doi:10.7169/facm/2013.48.1.11.

\bibitem{Pinch2006e18}
R.~G.~E. Pinch,
\emph{The Carmichael numbers up to $10^{18}$},
arXiv:math/0604376, 2006.

\bibitem{Pinch2007}
R.~G.~E. Pinch,
\emph{The Carmichael numbers up to $10^{21}$},
in A.-M. Ernvall-Hyt\"onen, M. Jutila, J. Karhum\"aki, and A. Lepist\"o (eds.),
\emph{Proceedings of the Conference on Algorithmic Number Theory},
TUCS General Publication 46, Turku Centre for Computer Science, 2007,
pp.~129--131.

\bibitem{PinchSmallIndex}
R.~G.~E. Pinch,
\emph{Carmichael numbers with small index},
ANTS VII poster, 2006,
\url{https://antsmath.org/ANTSVII/Poster/Pinch_Poster2.pdf}.

\bibitem{Pomerance1977}
C.~Pomerance,
\emph{On composite $n$ for which $\phi(n)\mid n-1$},
Acta Arith. \textbf{28} (1976), 387--389,
doi:10.4064/aa-28-4-387-389.

\bibitem{Pomerance1996}
C.~Pomerance,
\emph{A tale of two sieves},
Notices Amer. Math. Soc. \textbf{43} (1996), 1473--1485.

\bibitem{Pratt1975}
V.~R. Pratt,
\emph{Every prime has a succinct certificate},
SIAM J. Comput. 4 (1975), 214--220.

\bibitem{Ribenboim1996}
P.~Ribenboim,
\emph{The New Book of Prime Number Records},
Springer, New York, 1996.

\bibitem{Shevelev2012}
V.~Shevelev, G.~Garc\'ia-Pulgar\'in, J.~M. Vel\'asquez-Soto, and
J.~H. Castillo,
\emph{Overpseudoprimes, and Mersenne and Fermat numbers as primover numbers},
J. Integer Seq. 15 (2012), Article 12.7.7.

\bibitem{Shor1997}
P.~W. Shor,
\emph{Polynomial-time algorithms for prime factorization and discrete
logarithms on a quantum computer},
SIAM J. Comput. \textbf{26} (1997), 1484--1509.

\bibitem{Somer1985}
L.~E. Somer,
\emph{The Divisibility and Modular Properties of $K$th-Order Linear
Recurrences over the Ring of Integers of an Algebraic Number Field with
Respect to Prime Ideals},
Ph.D. dissertation, University of Illinois at Urbana--Champaign, 1985.

\bibitem{Somer1989}
L.~Somer,
\emph{On Fermat $d$-pseudoprimes},
in J.-M. De Koninck and C. Levesque (eds.),
\emph{Number Theory: Proceedings of the International Number Theory
Conference, Universit\'e Laval, 1987},
de Gruyter, 1989, pp.~841--860,
doi:10.1515/9783110852790.841.

\bibitem{SomerLucas1998}
L.~Somer,
\emph{On Lucas $d$-pseudoprimes},
in G.~E. Bergum, A.~N. Philippou, and A.~F. Horadam (eds.),
\emph{Applications of Fibonacci Numbers}, Vol.~7,
Kluwer Academic Publishers, 1998, pp.~369--375.

\bibitem{SorensonWebster2017}
J.~P.~Sorenson and J.~Webster,
\emph{Strong pseudoprimes to twelve prime bases},
Math. Comp. \textbf{86} (2017), 985--1003,
doi:10.1090/mcom/3134.

\bibitem{Sutherland2007}
A.~V. Sutherland,
\emph{Order Computations in Generic Groups},
Ph.D. thesis, Massachusetts Institute of Technology, 2007.

\bibitem{Wagstaff1982}
S.~S. Wagstaff, Jr.,
\emph{Pseudoprimes and a generalization of Artin's conjecture},
Acta Arith. 41 (1982), 141--150,
doi:10.4064/aa-41-2-141-150.

\end{thebibliography}

\end{document}
